\documentclass[11pt,a4paper]{report}

% =========================
% Packages (minimum required)
% =========================
\usepackage[a4paper,margin=2.5cm]{geometry}
\usepackage{microtype}
\usepackage{setspace}
\usepackage[T1]{fontenc}
\usepackage[utf8]{inputenc}
\usepackage{lmodern}
\usepackage{cmap}
\usepackage[french]{babel}
\usepackage{newunicodechar}
\newunicodechar{«}{``}
\newunicodechar{»}{''}
\renewcommand{\og}{``}
\renewcommand{\fg}{''}

\usepackage{amsmath,amssymb,amsthm,mathtools}
\usepackage{bm}
\usepackage{graphicx}
\usepackage{booktabs}
\usepackage{longtable}
\usepackage{enumitem}
\usepackage{listings}
\usepackage{comment}
\usepackage{xcolor}
\usepackage{tikz}
\usetikzlibrary{arrows.meta,positioning,calc}
\usepackage{pgfplots}
\pgfplotsset{compat=1.18}

\usepackage[hidelinks]{hyperref}
\usepackage[nameinlink,capitalise,noabbrev]{cleveref}

\setstretch{1.16}

% =========================
% Colors / Listings
% =========================
\definecolor{codebg}{RGB}{248,248,248}
\definecolor{codeframe}{RGB}{210,210,210}
\definecolor{codetxt}{RGB}{20,20,20}
\definecolor{kw}{RGB}{0,92,170}
\definecolor{cm}{RGB}{90,110,90}
\definecolor{st}{RGB}{140,60,60}

\lstdefinestyle{py}{
language=Python,
basicstyle=\ttfamily\small\color{codetxt},
backgroundcolor=\color{codebg},
frame=single,
rulecolor=\color{codeframe},
numbers=left,
numberstyle=\ttfamily\tiny\color{black!60},
xleftmargin=2.2em,
framexleftmargin=1.8em,
showstringspaces=false,
tabsize=4,
breaklines=true,
keywordstyle=\color{kw}\bfseries,
commentstyle=\color{cm}\itshape,
stringstyle=\color{st},
upquote=true
}
\lstset{style=py}

% =========================
% Theorem-like environments
% =========================
\theoremstyle{definition}
\newtheorem{definition}{Définition}[chapter]
\newtheorem{example}{Exemple}[chapter]
\theoremstyle{plain}
\newtheorem{theorem}{Théorème}[chapter]
\newtheorem{proposition}{Proposition}[chapter]
\newtheorem{lemma}{Lemme}[chapter]
\theoremstyle{remark}
\newtheorem{remark}{Remarque}[chapter]

% =========================
% Macros (notations)
% =========================
\newcommand{\dd}{\mathrm{d}}
\newcommand{\E}{\mathbb{E}}
\newcommand{\Var}{\mathrm{Var}}
\newcommand{\Cov}{\mathrm{Cov}}
\newcommand{\Corr}{\mathrm{Corr}}
\newcommand{\Prob}{\mathbb{P}}
\newcommand{\1}{\mathbf{1}}
\newcommand{\Normal}{\mathcal{N}}
\newcommand{\Pmeas}{\mathbb{P}}
\newcommand{\Qmeas}{\mathbb{Q}}
\newcommand{\R}{\mathbb{R}}
\newcommand{\F}{\mathcal{F}}
\newcommand{\G}{\mathcal{G}}
\newcommand{\argmin}{\operatorname*{arg\,min}}
\newcommand{\argmax}{\operatorname*{arg\,max}}

\newcommand{\filepath}[1]{\texttt{\detokenize{#1}}}
\newcommand{\codeid}[1]{\texttt{\detokenize{#1}}}
\newcommand{\NonImplem}{\textbf{non implémenté dans le code actuel}}

% =========================
% Boxes
% =========================
\newsavebox{\juryboxbox}
\newenvironment{jurybox}{
\par\medskip
\begin{lrbox}{\juryboxbox}
\begin{minipage}{0.97\textwidth}
\textbf{Ce que le jury doit retenir}\par\smallskip
\ignorespaces
}{
\unskip\end{minipage}
\end{lrbox}
\noindent\fcolorbox{black!30}{black!3}{\usebox{\juryboxbox}}
\par\medskip
}

\newcommand{\qa}[4]{
\begin{itemize}[leftmargin=*]
\item \textbf{À quoi ça sert ?} #1
\item \textbf{Hypothèses.} #2
\item \textbf{Limites.} #3
\item \textbf{Applications.} #4
\end{itemize}
}

% =========================
% Document
% =========================
\begin{document}

% -------------------------
% Title page
% -------------------------
\begin{titlepage}
\centering
\vspace*{2cm}
{\Large\bfseries PaperTradingApp\par}
\vspace{0.4cm}
{\large Dossier théorique \& dérivations pédagogiques (from ground up)\par}
\vspace{1.2cm}
{\large\itshape Public : jury financier (Master)\par}
\vspace{1.0cm}

\noindent\rule{0.9\textwidth}{0.4pt}\par
\vspace{0.6cm}

\begin{minipage}{0.9\textwidth}
\small
\textbf{Positionnement de ce dossier (rappel)}\par
\begin{itemize}[leftmargin=*]
\item Ce document est un \textbf{dossier pédagogique} centré sur les \textbf{hypothèses}, la \textbf{théorie} et les \textbf{dérivations}.
\item Les sections \textbf{B. Théorie} et \textbf{C. Dérivations} sont développées en détail (pas à pas, avec intuition et vérifications).
\item Le \textbf{lien avec le projet} est regroupé \textbf{en annexe} pour ne pas interrompre le fil théorique.
\end{itemize}
\end{minipage}

\vfill
{\small Généré le \today\par}
\end{titlepage}

% -------------------------
% Abstract / toc / figures
% -------------------------
\chapter*{Résumé}
Ce dossier regroupe les fondations du \textbf{calcul stochastique} et du \textbf{pricing d'options} utilisées en finance quantitative, puis détaille les dérivations standards (au niveau Master) qui relient \emph{hypothèses} $\to$ \emph{dynamiques} $\to$ \emph{prix} $\to$ \emph{paramètres}.

Le fil conducteur est volontairement \textbf{pédagogique} : chaque brique est reconstruite \textbf{à partir des hypothèses} (probabilités \& absence d'arbitrage) et \textbf{dérivée} étape par étape, avec intuition et contrôles de cohérence. En particulier :
\begin{itemize}[leftmargin=*]
\item Pricing d'options : Black--Scholes (PDE, formule fermée, grecs, volatilité implicite), arbres, Monte Carlo et principe de Longstaff--Schwartz.
\item Volatilité : surface d'IV, interpolation et arbitrage statique, calibration (problème inverse) et modèles paramétriques (SABR, Heston, sauts, rough/Volterra).
\item Couverture : principe de réplication, limites en marché incomplet, et formulation séquentielle (MDP) d'un hedge discret.
\item Taux : courbe de taux, discount factors, taux zéros, forwards, paramétrisation Nelson--Siegel.
\item Allocation \& risque : Markowitz, risk parity, eigen-portfolios, et métriques de risque (exposition, PnL, quantiles).
\end{itemize}

\tableofcontents
\listoffigures

% =========================================================
% CHAPITRE 0 — Architecture \& Notations
% =========================================================
\setcounter{chapter}{-1}
\chapter{Architecture \& Notations}

\section{A. Problème financier (objectif, inputs, outputs)}
L'objectif global est un pipeline \emph{data} $\to$ \emph{modèles} $\to$ \emph{décision} $\to$ \emph{exécution} :
\begin{itemize}[leftmargin=*]
\item \textbf{Inputs} : prix spot, historiques (OHLC), chaîne d'options (prix/IV par strike et maturité), structure de taux (zéros/DF), positions et contraintes (cash, limites de risque, coûts).
\item \textbf{Outputs} : prix d'options (et payoffs), grecs, surfaces de volatilité (marché vs modèle), paramètres calibrés, scénarios de couverture, plans de rebalancement et indicateurs de risque.
\end{itemize}

\section{B. Théorie (définitions, hypothèses)}
\subsection*{Séparation des responsabilités (idée générale)}
Pour garder le raisonnement financier lisible et \emph{audit-able}, on sépare conceptuellement :
\begin{itemize}[leftmargin=*]
\item \textbf{Données} : observation et nettoyage (spots, taux, options).
\item \textbf{Noyau quantitatif} : modèles et valorisation (mesure risque-neutre, pricers, grecs).
\item \textbf{Inverse problem} : calibration (choix de métrique, contraintes, régularisation).
\item \textbf{Décision} : couverture/allocation sous contraintes (discrétisation, coûts, limites).
\item \textbf{Exécution} : traduction d'une décision en trades (avec contrôle de risque opérationnel).
\end{itemize}

\subsection*{Intuition (pourquoi c'est crucial en finance)}
Une chaîne de calcul quantitative doit pouvoir :
\begin{itemize}[leftmargin=*]
\item \textbf{reproduire} un résultat (mêmes hypothèses $\Rightarrow$ même prix) ;
\item \textbf{auditer} le raisonnement (quelles données ? quel modèle ? quels paramètres ?) ;
\item \textbf{remplacer} une source de données ou une méthode numérique sans «casser» la théorie ;
\item \textbf{tester} séparément les hypothèses (absence d'arbitrage, bornes, monotonicité, diagnostics).
\end{itemize}

\qa
{Garantir qu'un résultat numérique (prix, calibration) est traçable, testable, et reproductible.}
{Séparation stricte des responsabilités : la couche de calcul ne dépend pas de la présentation, et la décision n'altère pas les données observées.}
{MVC ne garantit pas la qualité des modèles, seulement l'industrialisation et l'auditabilité.}
{Industrialisation : validation, reproductibilité, et articulation vers une boucle décision/exécution (définie conceptuellement).}

\subsection*{Notations transversales (utilisées dans tout le dossier)}
On fixe un horizon temporel $t\in[0,T]$ et un espace probabilisé filtré $(\Omega,\F,(\F_t)_{t\ge 0},\Pmeas)$.

\paragraph{Sous-jacent, numéraire, taux.}
\begin{itemize}[leftmargin=*]
\item $S_t$ : prix du sous-jacent (action/ETF) à l'instant $t$.
\item $r$ : taux sans risque (continu). Quand une courbe est disponible, on écrit plutôt $r(t)$ et on actualise via une intégrale.
\item $q$ : \emph{dividend yield} continu (ou \emph{repo rate} dans certains marchés) : il représente un «rendement» payé par l'actif, ce qui diminue la croissance du prix \emph{ex-dividend}.
\item $B_t$ : numéraire monétaire (compte rémunéré) ; en taux constant, $B_t=e^{rt}$.
\item $D(0,t)$ : facteur d'actualisation, $D(0,t)=e^{-rt}$ (ou $D(0,t)=\exp(-\int_0^t r(u)\,\dd u)$ si $r$ dépend du temps).
\end{itemize}

\paragraph{Mesures $\Pmeas$ et $\Qmeas$.}
\begin{itemize}[leftmargin=*]
\item $\Pmeas$ : mesure «réelle» (statistique) qui gouverne les probabilités physiques (historiques).
\item $\Qmeas$ : mesure risque-neutre associée au numéraire $B$ : sous $\Qmeas$, les prix \emph{actualisés} des actifs (corrigés de $q$) sont des martingales (absence d'arbitrage).
\item La différence $\Pmeas$ vs $\Qmeas$ est une différence de \emph{drift} (prime de risque) ; la diffusion (volatilité instantanée) est souvent la même dans les modèles de base, mais pas toujours en modèles plus riches.
\end{itemize}

\paragraph{Contrats d'options et fonctions de payoff.}
Un payoff européen s'écrit comme une fonction mesurable $\Phi(S_T)$ et son prix à $t$ est noté $V(t,S_t)$.
Exemples :
\[
\Phi_{\text{call}}(S_T)=(S_T-K)^+,\qquad \Phi_{\text{put}}(S_T)=(K-S_T)^+.
\]
On utilise systématiquement la convention $x^+=\max(x,0)$.

\paragraph{Volatilité et variables de moneyness.}
\begin{itemize}[leftmargin=*]
\item $\sigma$ : volatilité (constante en Black--Scholes, stochastique en Heston, «rough» en rBergomi, etc.).
\item $K$ : strike ; $T$ : maturité.
\item On introduit souvent le \emph{log-moneyness} $k=\ln(K/F_{0,T})$ où $F_{0,T}=S_0e^{(r-q)T}$ est le forward (taux constants).
\item Total variance : $w(k,T)=\sigma_{\text{impl}}(k,T)^2\,T$ est un objet plus régulier que $\sigma_{\text{impl}}$ et central pour discuter l'arbitrage statique.
\end{itemize}

\subsection*{Préliminaires probabilistes (rappel L3 $\to$ M1)}
Le calcul stochastique se formalise sur un espace probabilisé filtré $(\Omega,\F,(\F_t)_{t\ge 0},\Pmeas)$. On rappelle ici les points \emph{indispensables} pour comprendre les dérivations (et éviter les confusions de notation).

\begin{definition}[$\sigma$-algèbre, variable aléatoire]
Une $\sigma$-algèbre $\F$ sur $\Omega$ est une famille d'événements stable par complémentaire et réunion dénombrable.
Une variable aléatoire $X$ est une application mesurable $(\Omega,\F)\to(\R,\mathcal{B}(\R))$.
\end{definition}

\begin{definition}[Intégrabilité et espaces $L^p$]
Pour $p\ge 1$, on définit
\[
L^p(\Omega,\F,\Pmeas)=\{X \text{ mesurable }:\ \E[|X|^p]<\infty\}.
\]
On utilisera surtout $L^1$ (pour parler d'espérance) et $L^2$ (pour l'intégrale d'Itô, via des arguments de projection et d'isométrie).
\end{definition}

\subsection*{Espérance conditionnelle : la notion centrale derrière «prix = espérance»}
Intuitivement, $\E[X\mid \G]$ est la \emph{meilleure estimation} de $X$ lorsqu'on ne connaît que l'information $\G\subseteq\F$.

\begin{definition}[Espérance conditionnelle]
Soit $X\in L^1$ et $\G\subseteq\F$ une sous-$\sigma$-algèbre.
Une version de $\E[X\mid \G]$ est une variable aléatoire $\G$-mesurable telle que, pour tout événement $A\in\G$,
\[
\int_A \E[X\mid \G]\,\dd\Pmeas=\int_A X\,\dd\Pmeas.
\]
\end{definition}

\paragraph{Propriétés à connaître (et à utiliser sans hésiter).}
Pour $X,Y\in L^1$ et $a,b\in\R$ :
\begin{itemize}[leftmargin=*]
\item \textbf{Linéarité} : $\E[aX+bY\mid \G]=a\E[X\mid \G]+b\E[Y\mid \G]$ ;
\item \textbf{Positivité} : $X\ge 0\Rightarrow \E[X\mid \G]\ge 0$ ;
\item \textbf{Règle de la tour} : si $\mathcal{H}\subseteq\G$, alors $\E[\E[X\mid \G]\mid \mathcal{H}]=\E[X\mid \mathcal{H}]$ ;
\item \textbf{Extraction} : si $Y$ est $\G$-mesurable et $XY\in L^1$, alors $\E[XY\mid \G]=Y\E[X\mid \G]$.
\end{itemize}

\subsection*{Lecture géométrique en $L^2$ : projection orthogonale}
Lorsque $X\in L^2$, l'espérance conditionnelle est la projection orthogonale sur le sous-espace des variables $\G$-mesurables :

\begin{proposition}[Caractérisation $L^2$]
Si $X\in L^2$, alors $Y=\E[X\mid \G]$ est l'unique variable $\G$-mesurable dans $L^2$ telle que
\[
\E\big[(X-Y)Z\big]=0\qquad \text{pour tout } Z\in L^2 \text{ et } \G\text{-mesurable}.
\]
\end{proposition}

\begin{remark}
Cette propriété est la bonne façon de retenir le rôle de $L^2$ dans l'intégrale d'Itô : on construit d'abord l'intégrale sur des processus simples, puis on \emph{prolonge par densité} en contrôlant la norme $L^2$ grâce à l'isométrie.
\end{remark}

\subsection*{Filtration, adaptativité, prévisibilité : formaliser «l'information disponible»}
La filtration $(\F_t)_{t\ge 0}$ représente l'information disponible au cours du temps.

\begin{definition}[Adapté]
Un processus $(X_t)_{t\ge 0}$ est \emph{adapté} si $X_t$ est $\F_t$-mesurable pour tout $t$.
\end{definition}

\paragraph{Pourquoi ce point est non négociable en finance.}
Une stratégie de couverture $\Delta_t$ doit être déterminée à partir des informations observables à $t$ : elle ne peut pas dépendre de $S_{t+\varepsilon}$.
Cette contrainte est précisément l'adaptativité (et, en continu, une notion légèrement plus fine : la prévisibilité).

\begin{definition}[Prévisible (idée)]
En temps discret, «prévisible» signifie typiquement $\Delta_{t_{k}}$ mesurable par rapport à $\F_{t_{k-1}}$.
En temps continu, la définition passe par la \emph{$\sigma$-algèbre prévisible} et garantit que l'intégrande ne regarde pas «l'instant même» de manière anticipative.
\end{definition}

\subsection*{Martingales : le langage naturel de l'absence d'arbitrage}
\begin{definition}[Martingale]
Un processus $(M_t)_{t\ge 0}$ est une martingale (par rapport à $(\F_t)$ et $\Pmeas$) si :
(i) $M_t$ est adapté, (ii) $\E[|M_t|]<\infty$, et (iii) pour $s\le t$,
\[
\E[M_t\mid \F_s]=M_s.
\]
\end{definition}

\begin{remark}[Intuition]
Une martingale est un «jeu équitable» : la meilleure prévision de demain, sachant aujourd'hui, est la valeur d'aujourd'hui.
En finance, ce n'est pas une hypothèse psychologique : c'est une conséquence de l'absence d'arbitrage \emph{après actualisation et choix d'un numéraire}.
\end{remark}

\paragraph{Sous-martingale, super-martingale.}
On parlera aussi de sous-martingale (resp. super-martingale) lorsque $\E[M_t\mid \F_s]\ge M_s$ (resp. $\le M_s$). Les inégalités de Doob et les arguments de contrôle de maxima se formulent naturellement à ce niveau.

\subsection*{Temps d'arrêt et arrêt optionnel (rappel utile)}
Un \emph{temps d'arrêt} $\tau$ est une variable aléatoire à valeurs dans $[0,\infty]$ telle que $\{\tau\le t\}\in \F_t$ pour tout $t$.
Il formalise des événements «détectables» au moment où ils se produisent (ex. franchissement d'une barrière).

\begin{theorem}[Arrêt optionnel, version simple]
Si $(M_t)$ est une martingale et si $\tau$ est un temps d'arrêt borné, alors $\E[M_\tau]=\E[M_0]$.
\end{theorem}

\begin{remark}
Les versions générales exigent des conditions d'intégrabilité (uniform integrability). On retiendra surtout l'idée : \emph{on ne fabrique pas un profit espéré en arrêtant une martingale «au bon moment»}, ce qui est cohérent avec l'absence d'arbitrage.
\end{remark}

\subsection*{Changement de mesure : le pont entre $\Pmeas$ et $\Qmeas$}
On utilise très souvent deux mesures :
\begin{itemize}[leftmargin=*]
\item $\Pmeas$ : mesure statistique (historiques) ;
\item $\Qmeas$ : mesure de valorisation (risk-neutral), sous laquelle le prix actualisé devient une martingale.
\end{itemize}
Le passage $\Pmeas\to\Qmeas$ s'exprime par une densité de Radon--Nikodym.

\begin{definition}[Densité de Radon--Nikodym]
On dit que $\Qmeas$ est \emph{absolument continue} par rapport à $\Pmeas$ (noté $\Qmeas\ll\Pmeas$) s'il existe une variable $Z\ge 0$ telle que, pour tout $A\in\F$,
\[
\Qmeas(A)=\E^{\Pmeas}[\1_A Z].
\]
On note alors $Z=\frac{\dd\Qmeas}{\dd\Pmeas}$.
\end{definition}

\paragraph{Formule de transfert des espérances.}
Si $Z=\frac{\dd\Qmeas}{\dd\Pmeas}$ et $X\in L^1(\Qmeas)$, alors
\[
\E^{\Qmeas}[X]=\E^{\Pmeas}[ZX].
\]
Et (point crucial en finance dynamique) pour toute $\F_t$,
\[
\E^{\Qmeas}[X\mid \F_t]=\frac{\E^{\Pmeas}[ZX\mid \F_t]}{\E^{\Pmeas}[Z\mid \F_t]}\qquad (\text{Bayes}).
\]
Ces identités seront utilisées implicitement dans Girsanov et dans le changement de numéraire.

\subsection*{Exemple fil rouge : modèle 1 période et probabilité risque-neutre}
Avant les diffusions, il est sain de revoir le cas discret (1 période), car il contient déjà l'essentiel : \emph{prix = espérance sous une mesure ajustée}.

\begin{example}[Binomial 1 période]
On a un actif sans risque $B_0=1$ et $B_1=1+r$, et un actif risqué $S_0$ tel que
\[
S_1=\begin{cases}
S^u & \text{avec probabilité } p,\\
S^d & \text{avec probabilité } 1-p,
\end{cases}
\qquad S^d<S^u.
\]
Une stratégie auto-financée $(\Delta,\beta)$ a une valeur $V_1=\Delta S_1+\beta B_1$.
Pour un payoff $H(S_1)$, l'absence d'arbitrage équivaut (ici) à l'existence d'un $q\in(0,1)$ tel que
\[
S_0=\frac{1}{1+r}\big(qS^u+(1-q)S^d\big).
\]
Alors le \textbf{prix sans arbitrage} d'un payoff réplicable est
\[
V_0=\frac{1}{1+r}\E^{\Qmeas}[H(S_1)]\quad \text{avec}\quad \Qmeas(S_1=S^u)=q.
\]
\end{example}

\begin{remark}
Ce résultat n'est pas «magique» : on peut le dériver (i) par réplication (résoudre deux équations linéaires), ou (ii) par changement de mesure.
Dans les modèles continus, on remplace $q$ par une densité (Girsanov) et on remplace la réplication exacte par une PDE (marché complet) ou par un critère de pricing (marché incomplet).
\end{remark}

\subsection*{Mini-exercices (pour consolider)}
\begin{enumerate}[leftmargin=*]
\item Montrer que si $M$ est une martingale intégrable alors $\E[M_t]=\E[M_0]$.
\item Soit $\G\subseteq\F$ et $X\in L^2$. Montrer que $Y=\E[X\mid \G]$ minimise $\E[(X-Z)^2]$ parmi les variables $\G$-mesurables $Z$.
\item Dans le binomial 1 période, exprimer $(\Delta,\beta)$ qui réplique un call $H(S_1)=(S_1-K)^+$.
\end{enumerate}

\subsection*{Hypothèses transversales : ce qui est supposé (et ce qui ne l'est pas)}
Pour rendre les dérivations lisibles, on commence par des hypothèses «de base» (marché frictionless, trading continu, pas de coûts). Elles sont ensuite relâchées :
\begin{itemize}[leftmargin=*]
\item \textbf{Incomplétude / non-réplication} : en présence de sauts, de volatilité stochastique, de contraintes ou de rebalancement discret, la réplication parfaite échoue ; le pricing devient une question de mesure risque-neutre \emph{non unique} ou de choix de critère.
\item \textbf{Calibrage} : une surface de volatilité implicite est un \emph{résumé} des prix de marché ; calibrer un modèle revient à résoudre un problème inverse potentiellement mal posé.
\item \textbf{Numerics} : une formule fermée est rare ; on alterne entre PDE, Monte Carlo, FFT et approximations, d'où l'importance de la stabilité et des diagnostics.
\end{itemize}

\section{C. Dérivations (ultra détaillées)}
Ce chapitre formalise la chaîne \emph{données $\to$ modèles $\to$ prix $\to$ décisions} à un niveau conceptuel, indépendamment de toute implémentation.

\subsection*{Données $\to$ variables d'état : ce qu'on retient vraiment}
Un point pédagogique crucial est qu'une «donnée de marché» brute n'est pas directement un objet de pricing.
On commence par construire une \emph{variable d'état} $X_t$ (ou un petit nombre de variables) qui résume l'information pertinente à l'instant $t$.
Exemples :
\begin{itemize}[leftmargin=*]
\item option vanille : $X_t=(t,S_t)$ suffit en Black--Scholes ;
\item volatilité stochastique : $X_t=(t,S_t,V_t)$ (ex. Heston) ;
\item portefeuille : $X_t=(t,\text{positions},\text{cash},\text{paramètres de risque},\text{contraintes})$.
\end{itemize}
Mathématiquement, $X_t$ doit être \emph{$\F_t$-mesurable} : l'état à $t$ ne doit dépendre que de l'information disponible à $t$.

\subsection*{Prix : un opérateur conditionnel sous mesure risque-neutre}
Sous absence d'arbitrage (et dans un cadre standard), il existe une mesure risque-neutre $\Qmeas$ telle que, pour un payoff européen $\Phi(S_T)$,
\[
V(t,X_t)=\E^{\Qmeas}\big[D(t,T)\,\Phi(S_T)\mid \F_t\big],
\qquad D(t,T)=\exp\Big(-\int_t^T r(u)\,\dd u\Big).
\]
Ce point n'est pas seulement une «formule» : c'est un \emph{principe d'identification} qui force des propriétés structurelles :
\begin{itemize}[leftmargin=*]
\item \textbf{linéarité} : $V(\alpha\Phi_1+\beta\Phi_2)=\alpha V(\Phi_1)+\beta V(\Phi_2)$ ;
\item \textbf{positivité} : $\Phi\ge 0 \Rightarrow V\ge 0$ ;
\item \textbf{monotonicité} : $\Phi_1\le \Phi_2 \Rightarrow V(\Phi_1)\le V(\Phi_2)$ ;
\item \textbf{bornes d'arbitrage} : par exemple $0\le C\le S_0e^{-qT}$ pour un call européen.
\end{itemize}
Ces propriétés servent de \emph{tests de cohérence} pour les calculs numériques et pour la calibration.

\subsection*{Calibration : un problème inverse (et souvent mal posé)}
Un modèle paramétrique $\mathcal{M}(\theta)$ fournit des prix (ou IV) théoriques $V^{\mathcal{M}}(\theta)$.
Le marché fournit des observations $V^{\text{mkt}}$ sur un ensemble d'options (strikes, maturités).
Calibrer revient à résoudre
\[
\hat\theta=\argmin_\theta\ \mathcal{L}\big(V^{\mathcal{M}}(\theta),V^{\text{mkt}}\big),
\]
où $\mathcal{L}$ encode le choix \emph{métier} : erreur en prix vs en IV, pondérations par Vega, robustesse aux outliers (Huber), et contraintes (ex. Feller en Heston).
Au niveau M1, il faut retenir que :
\begin{itemize}[leftmargin=*]
\item le problème est \textbf{non convexe} et peut avoir plusieurs minima locaux ;
\item il peut être \textbf{mal conditionné} (paramètres peu identifiables sur une fenêtre de marché) ;
\item il nécessite souvent \textbf{régularisation} et \textbf{diagnostics} (stabilité temporelle des paramètres, sensibilité).
\end{itemize}

\subsection*{Décision : du pricing à l'optimisation séquentielle}
Une fois un prix et des sensibilités disponibles, une «décision» (couverture, rebalancement) s'écrit naturellement comme un \emph{problème de contrôle} en temps discret :
choisir une action $a_t$ (quantités tradées) à chaque date sur la base de l'état $X_t$ pour minimiser une fonction de coût :
\[
\min_{\pi}\ \E\Big[\sum_{t\in\{t_0,\ldots,t_{n-1}\}} \text{coût}(X_t,a_t)+\text{risque final}(X_T)\Big],
\]
avec contraintes (budget, limites, coûts de transaction). Cette formulation explique pourquoi un cadre MDP (et, éventuellement, RL) est un \emph{langage naturel} pour décrire un hedge discret.

\subsection*{Dérivation conceptuelle : pourquoi ce graphe est le bon (en finance)}
Le découpage «données $\to$ modèles $\to$ décision $\to$ exécution» n'est pas qu'un choix d'organisation : c'est la transcription d'un raisonnement financier standard.
\begin{enumerate}[leftmargin=*]
\item \textbf{Étape 1 : définir la variable d'état.} Exemple : pour une option vanille, l'état minimal est $(t,S_t)$ ; pour Heston, c'est $(t,S_t,V_t)$ ; pour un portefeuille, c'est (positions, cash, contraintes, paramètres de risque).
\item \textbf{Étape 2 : choisir une hypothèse dynamique.} Une hypothèse de diffusion (GBM) ou de volatilité stochastique (Heston) ne vaut pas «vérité» : elle fixe une \emph{famille de lois} pour $S_T$ et donc une famille de prix pour $\Phi(S_T)$.
\item \textbf{Étape 3 : fixer un principe de valorisation.} En absence d'arbitrage, le prix est une espérance actualisée sous une mesure $\Qmeas$ (ou une solution de PDE équivalente) ; dans un marché incomplet, il faut choisir un critère (min-variance, entropie minimale, etc.).
\item \textbf{Étape 4 : calibrer.} Le marché fournit des prix (ou des IV). Le modèle fournit des prix (ou des IV) comme fonction de paramètres. On résout un problème inverse (souvent non convexe) : \emph{trouver des paramètres qui reproduisent au mieux le marché}.
\item \textbf{Étape 5 : décider et gérer le risque.} Un modèle donne un prix «théorique», des grecs, une distribution de PnL, etc. La décision (hedge/rebalancement) se fait sous contraintes (discrétisation, coûts, limites de risque).
\item \textbf{Étape 6 : exécuter.} Sans exécution (même papier), la boucle de rétroaction est incomplète : on ne peut pas confronter modèle et résultats (PNL, slippage, erreurs de calibration).
\end{enumerate}
Ainsi, le graphe de dépendances n'est pas décoratif : il matérialise le chemin \emph{hypothèses $\to$ prix $\to$ décisions}.

\begin{comment}
\section{D. Numérique / algorithmes (stabilité, complexité, choix)}
Le dépôt mélange plusieurs familles d'algorithmes :
\begin{itemize}[leftmargin=*]
\item \textbf{Formules fermées} : Black--Scholes (ex. \filepath{app/model/options/engines/black_scholes.py}).
\item \textbf{PDE} : Crank--Nicolson pour BS en log-spot (classe \codeid{CrankNicolsonBS} dans \filepath{app/model/options/engines/pricing.py}).
\item \textbf{Arbres} : CRR binomial (ex. \codeid{price_american_crr} dans \filepath{app/model/options/engines/crr.py}).
\item \textbf{Monte Carlo} : simulateurs GBM et exotiques (ex. \codeid{simulate_gbm_paths} dans \filepath{app/model/options/core/trees.py}).
\item \textbf{FFT} : Carr--Madan pour prix de calls (ex. \codeid{carr_madan_fft_call_prices} dans \filepath{app/model/volatility_models/common/fft.py}).
\item \textbf{Optimisation} : moindres carrés multi-start (SciPy) et design LHS (ex. \codeid{latin_hypercube_samples} dans \filepath{app/model/calibration/optimizers.py}).
\item \textbf{RL} : DQN numpy (replay buffer + target network) dans \filepath{app/model/hedger_v2/dqn_hedger.py}.
\end{itemize}
Le choix «plusieurs algorithmes» est cohérent en finance : on cherche un compromis précision/temps et une robustesse (fallbacks).

\end{comment}
\begin{comment}
\section{E. Implémentation dans le repo (fichiers, fonctions, pipeline, paramètres)}
\subsection*{Diagramme de flux (imposé)}
\begin{figure}[h]
\centering
\begin{tikzpicture}[
node distance=1.2cm and 1.2cm,
box/.style={draw, rounded corners, align=center, minimum width=3.6cm, minimum height=0.9cm},
arr/.style={-latex, thick}
]
\node[box] (md) {Market Data\\small (Stooq/Yahoo/Alpaca)};
\node[box, right=of md] (opt) {Options \& Vol\\small (IV, pricing)};
\node[box, right=of opt] (cal) {Calibration\\small (SABR/Heston/...)};
\node[box, right=of cal] (hedge) {Hedging RL\\small (DQN)};
\node[box, below=of cal] (yc) {Yield Curve\\small (DF, forwards, NS)};
\node[box, left=of yc] (pa) {Portfolio Allocation\\small (Markowitz, RP, eigen)};
\node[box, left=of pa] (exec) {Trading/Execution\\small (Alpaca orders)};

\draw[arr] (md) -- (opt);
\draw[arr] (opt) -- (cal);
\draw[arr] (cal) -- (hedge);
\draw[arr] (yc) -- (pa);
\draw[arr] (pa) -- (exec);

\draw[arr, dashed] (yc) -- (opt);
\draw[arr, dashed] (cal) -- (opt);

\end{tikzpicture}
\caption{Flux logique cible : données $\to$ modèles $\to$ décision $\to$ exécution. Les flèches en pointillés correspondent à une intégration observable via \filepath{app/vue/components/options/router.py} (taux depuis Yield Curve) et \filepath{app/vue/tabs/tab_advanced_calibration.py} (envoi d'une surface calibrée vers Options).}
\label{fig:global-flow}
\end{figure}

\subsection*{Tableau de cartographie (module / rôle / fichiers / fonctions clés / I/O)}
\begin{longtable}{@{}p{3.0cm}p{3.6cm}p{4.8cm}p{4.2cm}@{}}
\toprule
\textbf{Module} & \textbf{Rôle} & \textbf{Fichiers (exemples)} & \textbf{Fonctions/classes clés (exemples)}\\
\midrule
\endhead
Vue (UI) & Navigation + rendu Streamlit & \filepath{streamlit_app.py}, \filepath{app/vue/main_app.py}, \filepath{app/vue/tabs/tab_.py} & \codeid{main} (\filepath{app/vue/main_app.py}), \codeid{render_tab} (tabs)\\
Contrôleurs & Glue View $\to$ Model & \filepath{app/controller/options_controller.py}, \filepath{app/controller/calibration_controller.py}, \filepath{app/controller/yieldcurve_controller.py}, \filepath{app/controller/hedger_v2_controller.py}, \filepath{app/controller/portfolio_and_risk_controller.py} & \codeid{run_advanced_surface_calibration}, \codeid{get_dqn_hedge_suggestion}, \codeid{compute_allocation}, \ldots\\
Options (Model) & Pricing/grecs/PDE/MC/exotiques & \filepath{app/model/options/service.py}, \filepath{app/model/options/engines/pricing.py}, \filepath{app/model/options/engines/black_scholes.py}, \filepath{app/model/options/core/trees.py} & \codeid{compute_price}, \codeid{CrankNicolsonBS}, \codeid{black_scholes_price}, \codeid{simulate_gbm_paths}\\
Calibration (Model) & Parsing surface, Heston V1, NN Heston & \filepath{app/model/calibration/market_surface.py}, \filepath{app/model/calibration/heston_calibrator.py}, \filepath{app/model/calibration/heston_nn.py} & \codeid{build_fixed_grid}, \codeid{calibrate_heston_least_squares}, \codeid{HestonSurfaceNet}\\
Vol models (Model) & SABR/Heston/jumps/rough/Volterra & \filepath{app/model/volatility_models/} & \codeid{hagan_black_iv}, \codeid{carr_madan_fft_call_prices}, \codeid{rheston_log_return_cf_markovian}, \ldots\\
Hedger RL (Model) & DQN numpy + envs & \filepath{app/model/hedger_v2/dqn_hedger.py}, \filepath{app/model/hedger_v2/env.py} & \codeid{DQNAgent}, \codeid{ReplayBuffer}, \codeid{SyntheticHedgingEnv}\\
Yield curve (Model) & Courbe, DF, forwards, Nelson--Siegel & \filepath{app/model/yieldcurve/service.py}, \filepath{app/model/yieldcurve/engine.py}, \filepath{app/model/yieldcurve/rates_utils.py} & \codeid{get_active_curve}, \codeid{fit_nelson_siegel_grid_search}, \codeid{get_r}\\
Allocation \& risk (Model) & Markowitz/RP/eigen + VaR-lite & \filepath{app/model/portfolio_allocation/engine.py}, \filepath{app/model/risk_management/engine.py} & \codeid{markowitz_optimize}, \codeid{risk_parity_optimize}, \codeid{compute_var_lite}\\
\bottomrule
\end{longtable}

\section{F. Exemple end-to-end (ce qui se passe quand l'utilisateur lance la fonctionnalité)}
\paragraph{Exemple 1 : calibrer une surface puis pricer dans Options.}
\begin{enumerate}[leftmargin=*]
\item L'utilisateur ouvre l'onglet «Calibration avancée» (\filepath{app/vue/tabs/tab_advanced_calibration.py}, \codeid{render_tab}).
\item Il charge une surface (Yahoo/CSV/Cache/Alpaca) puis lance la calibration via \codeid{CalibrationController.run_advanced_surface_calibration} (\filepath{app/controller/calibration_controller.py}).
\item Le contrôleur instancie un calibrateur (ex. \codeid{HestonFFTCalibrator} dans \filepath{app/model/volatility_models/heston/calibrator_fft.py}) et renvoie \codeid{iv_model} et \codeid{params}.
\item Un clic «Envoyer IV modèle vers Options» écrit \codeid{st.session_state["calib_model_surface_df"]} et force la source «Calibration» côté Options (\filepath{app/vue/tabs/tab_advanced_calibration.py}, bouton à proximité du bloc résultat).
\item L'onglet Options (\filepath{app/vue/components/options/router.py}, \codeid{render_options_router}) lit cette surface et l'affiche/consomme.
\end{enumerate}

\paragraph{Exemple 2 : couverture DQN.}
\begin{enumerate}[leftmargin=*]
\item L'utilisateur ouvre «Hedging Systems» (\filepath{app/vue/tabs/tab_hedging_systems.py}).
\item Il entraîne le modèle (bouton «Train / update DQN») : appel \codeid{ctrl.train_dqn_model} $\to$ \codeid{load_or_train_dqn_model} dans \filepath{app/model/hedger_v2/dqn_hedger.py}.
\item Il demande une suggestion : \codeid{ctrl.get_dqn_hedge_suggestion} $\to$ \codeid{suggest_hedge_action}.
\item Optionnellement, il exécute : \codeid{ctrl.execute_dqn_hedge} $\to$ \codeid{AlpacaHedgerClient.submit_market_order} (\filepath{app/model/hedger_v2/alpaca_client.py}).
\end{enumerate}

\section{G. Limites \& extensions crédibles}
\begin{itemize}[leftmargin=*]
\item La présence de plusieurs moteurs (formule/PDE/tree/MC/FFT) exige une gouvernance : quels moteurs sont «référence» par produit ?
\item Certaines briques existent comme \emph{scaffold} : \codeid{app.model.calibration.logic.run_calibration} renvoie explicitement «Calibration non implémentée» (\filepath{app/model/calibration/logic.py}). L'onglet «Calibration avancée» utilise une voie séparée (contrôleur unifié) qui, elle, est opérationnelle.
\item Dividendes : \codeid{get_q} retourne 0.0 (placeholder) dans \filepath{app/model/yieldcurve/rates_utils.py} ; en production, il faut une source robuste (dividendes attendus/forwards).
\end{itemize}

\end{comment}
\begin{jurybox}
Ce dossier ne cherche pas à empiler des recettes : il formalise une \textbf{chaîne de raisonnement} cohérente, de la donnée brute jusqu'à une décision sous contraintes. Le point clé à retenir est la séparation \emph{hypothèses} (modèle et mesure) / \emph{calcul} (pricers, calibration) / \emph{décision} (couverture, allocation, risque).
\end{jurybox}

% =========================================================
% PARTIE I — Stochastic calculus \& pricing
% =========================================================
\part{Stochastic calculus \& pricing (base indispensable)}

% ---------------------------------------------------------
\chapter{Brownien et intégrale stochastique}

\section{A. Problème financier (objectif, inputs, outputs)}
On veut modéliser l'incertitude du prix d'un sous-jacent (action) à court terme pour :
\begin{itemize}[leftmargin=*]
\item pricer des options (prix = espérance actualisée sous une mesure appropriée) ;
\item simuler des trajectoires (Monte Carlo) ;
\item construire des risques (grecs, PnL, mesures de risque).
\end{itemize}
\textbf{Input} : paramètres $(S_0,r,q,\sigma,T)$ et un générateur de bruit. \textbf{Output} : une trajectoire $S_t$ (ou une loi de $S_T$).

\section{B. Théorie (définitions, hypothèses)}
\subsection*{Définition formelle}
\begin{definition}[Mouvement brownien]
Sur un espace probabilisé filtré $(\Omega,\F,(\F_t){t\ge 0},\Pmeas)$, un processus $(W_t){t\ge 0}$ est un \emph{mouvement brownien standard} si :
\begin{enumerate}[leftmargin=*]
\item $W_0=0$ p.s. ;
\item $W_t-W_s \sim \Normal(0,t-s)$ pour $0\le s<t$ ;
\item les accroissements sont indépendants ;
\item $t\mapsto W_t$ est continu.
\end{enumerate}
\end{definition}

L'intégrale stochastique (Itô) est l'opérateur qui donne sens à
\[
\int_0^T H_t,\dd W_t
\]
pour des processus adaptés $H_t$ (typiquement carrés-intégrables).

\subsection*{Intuition financière}
\begin{itemize}[leftmargin=*]
\item Le brownien représente un «bruit» agrégé (micro-chocs) qui, à l'échelle du modèle, a des incréments gaussiens.
\item L'intégrale stochastique est la façon correcte d'accumuler un risque instantané : si l'exposition $H_t$ varie dans le temps, la volatilité cumulée n'est pas une somme classique mais une intégrale par rapport à $W$.
\end{itemize}

\qa
{C'est le socle de la dynamique GBM (Black--Scholes) et des simulations Monte Carlo utilisées partout (pricing, exotiques, calibration rough/Volterra).}
{Marché frictionless pour les modèles de base ; dynamique continue ; volatilité constante si GBM ; indépendance/gaussianité du bruit.}
{Le brownien est une idéalisation : sauts, volatilité stochastique, microstructure, queues épaisses.}
{Simulation Monte Carlo des diffusions : l'approximation opérationnelle est toujours $\Delta W\approx \sqrt{\Delta t}\,Z$ avec $Z\sim\Normal(0,1)$, puis accumulation des incréments.}

\subsection*{Filtration, adaptativité et «information disponible»}
La filtration $(\F_t)_{t\ge 0}$ formalise l'information observable jusqu'au temps $t$.
Dire que $H_t$ est \emph{adapté} signifie que $H_t$ ne dépend que de l'information passée et présente (pas d'anticipation).
En finance, c'est la traduction mathématique de : «un hedge à $t$ ne peut pas dépendre de $S_{t+\varepsilon}$».

\subsection*{Construction (idée) de l'intégrale d'Itô}
On construit d'abord l'intégrale pour des \emph{processus simples} :
\[
H_t = \sum_{k=0}^{n-1} h_k\,\1_{(t_k,t_{k+1}]}(t),\qquad h_k\in\F_{t_k}.
\]
On définit alors
\[
\int_0^T H_t\,\dd W_t := \sum_{k=0}^{n-1} h_k\,(W_{t_{k+1}}-W_{t_k}).
\]
Puis on étend par densité à des processus $H$ dans $L^2(\Omega\times[0,T])$ (adaptés) via une limite en moyenne quadratique.

\begin{theorem}[Isométrie d'Itô]
Si $H$ est adapté et $\E\big[\int_0^T H_t^2\,\dd t\big]<\infty$, alors
\[
\E\Big[\Big(\int_0^T H_t\,\dd W_t\Big)^2\Big] = \E\Big[\int_0^T H_t^2\,\dd t\Big].
\]
En particulier, $\E[\int_0^T H_t\,\dd W_t]=0$ (l'intégrale d'Itô est une martingale).
\end{theorem}

\begin{remark}
L'isométrie est le pont «théorie $\leftrightarrow$ pratique» : elle dit que la variance d'un risque intégré est l'intégrale des variances instantanées.
Dans les modèles de hedging, cela justifie le rôle central de $\int H_t^2\,\dd t$ comme quantité de risque cumulée.
\end{remark}

\subsection*{Propriétés browniennes utiles pour les dérivés}
Au-delà de la définition, certaines propriétés structurent l'intuition de produits exotiques :
\begin{itemize}[leftmargin=*]
\item \textbf{Scaling} : $(W_{ct})_{t\ge 0}\overset{d}{=}(\sqrt{c}\,W_t)_{t\ge 0}$.
\item \textbf{Maximum et barrières (intuition)} : la loi de $\sup_{0\le u\le T} W_u$ et les probabilités de franchissement s'obtiennent via le principe de réflexion ; c'est la base de formules fermées pour certaines options barrières en Black--Scholes.
\item \textbf{Temps d'arrêt} : les hitting times $\tau_a=\inf\{t:W_t=a\}$ interviennent naturellement dès qu'une payoff dépend du premier passage d'un niveau.
\end{itemize}

% =========================================================
% Compléments M1 — Brownien (B. Théorie)
% =========================================================

\subsection*{(1) Ce que la définition implique vraiment : structure gaussienne et covariance}
La définition d'un mouvement brownien standard $(W_t)_{t\ge 0}$ peut se résumer en une phrase :
\emph{processus gaussien centré, à trajectoires continues, de covariance $\Cov(W_s,W_t)=\min(s,t)$}.
En effet :
\begin{itemize}[leftmargin=*]
\item Les incréments stationnaires gaussiens impliquent que, pour tout $t$, $W_t\sim\Normal(0,t)$.
\item Pour $0\le s\le t$, on écrit $W_t=W_s+(W_t-W_s)$ avec $(W_t-W_s)$ indépendant de $\F_s$ et centré. Alors
\[
\Cov(W_s,W_t)=\E[W_sW_t]=\E\!\big[W_s(W_s+(W_t-W_s))\big]=\E[W_s^2]=s.
\]
\end{itemize}
La covariance $\min(s,t)$ est la façon la plus compacte de retenir la dépendance temporelle : \emph{le passé jusqu'à $\min(s,t)$ est commun, puis les incréments au-delà sont indépendants}.

\subsection*{(2) Propriété de Markov et propriété de martingale (formalisées)}
\paragraph{Markov.}
Le Brownien est un processus de Markov : conditionnellement à $W_t$, le futur $(W_{t+u}-W_t)_{u\ge 0}$ est indépendant du passé et suit une loi brownienne.
Concrètement, pour une fonction mesurable bornée $g$,
\[
\E[g(W_{t+u})\mid \F_t]=\E[g(W_t+(W_{t+u}-W_t))\mid \F_t]
=h_u(W_t),
\]
où $h_u(x)=\E[g(x+Z)]$ avec $Z\sim\Normal(0,u)$.

\paragraph{Martingale.}
$(W_t)$ est une martingale sous $\Pmeas$ :
\[
\E[W_t\mid \F_s]=W_s,\qquad 0\le s\le t.
\]
\emph{Preuve courte.} Écrire $W_t=W_s+(W_t-W_s)$ et utiliser $\E[W_t-W_s\mid \F_s]=0$ par indépendance et centrage des incréments.

\paragraph{Pourquoi ça compte en finance ?}
Deux raisons structurantes :
\begin{itemize}[leftmargin=*]
\item La martingale capture l'idée «pas de dérive prédictible une fois l'information prise en compte». Sous $\Qmeas$, ce seront les \emph{prix actualisés} qui deviennent des martingales.
\item Le Markov est le pont vers les PDE : si le prix $V(t,S_t)$ dépend seulement de $(t,S_t)$ (pas du chemin), alors $V$ peut souvent être caractérisé par une équation de Kolmogorov / une PDE.
\end{itemize}

\subsection*{(3) Invariance d'échelle (scaling) : d'où vient le $\sqrt{\Delta t}$}
La propriété de scaling dit :
\[
(W_{ct})_{t\ge 0}\ \overset{d}{=}\ (\sqrt{c}\,W_t)_{t\ge 0}.
\]
\emph{Lecture.} Quand on change d'unité de temps (par exemple jours $\to$ années), l'échelle du bruit change comme $\sqrt{c}$.
C'est la justification mathématique de la règle opérationnelle :
\[
\Delta W \approx \sqrt{\Delta t}\,Z,\qquad Z\sim\Normal(0,1).
\]
\emph{Preuve (sur les lois finies-dimensionnelles).} Pour des temps $0\le t_1<\dots<t_n$, le vecteur $(W_{ct_1},\dots,W_{ct_n})$ est gaussien centré, de covariance
\[
\Cov(W_{ct_i},W_{ct_j})=\min(ct_i,ct_j)=c\,\min(t_i,t_j),
\]
qui est exactement la covariance de $(\sqrt{c}\,W_{t_1},\dots,\sqrt{c}\,W_{t_n})$.

\subsection*{(4) Propriétés de trajectoire : continuité oui, dérivabilité non}
Deux faits contre-intuitifs mais essentiels :
\begin{itemize}[leftmargin=*]
\item Les trajectoires de $W$ sont continues, mais \textbf{presque sûrement nulle part dérivables}.
\item La \textbf{variation totale} de $W$ sur tout intervalle est infinie (presque sûrement).
\end{itemize}
\paragraph{Pourquoi c'est important ?}
Parce que cela explique pourquoi le calcul différentiel «classique» échoue : on ne peut pas traiter $\dd W_t$ comme une différentielle ordinaire, et les termes d'ordre 2 (variation quadratique) deviennent déterminants.

\paragraph{Hölder (ordre de grandeur précis).}
On peut montrer que pour tout $\alpha<\tfrac12$, $W$ est Hölder-$\alpha$ sur $[0,T]$ p.s., et pour $\alpha\ge\tfrac12$ ce n'est plus vrai.
Intuition : $\Delta W=O(\sqrt{\Delta t})$ est exactement la frontière entre «variation finie» et «variation infinie».

\subsection*{(5) Principe de réflexion : maximum, barrières, temps de premier passage}
Le principe de réflexion est l'outil classique pour obtenir des lois de maximum/minimum du Brownien.
Soit $M_T=\sup_{0\le t\le T} W_t$.
Alors, pour $a>0$ :
\[
\Prob(M_T\ge a)=2\Prob(W_T\ge a)=2\bigl(1-N(a/\sqrt{T})\bigr).
\]
\paragraph{Idée de preuve (niveau M1).}
On considère l'événement où le chemin franchit $a$ avant $T$, et on «réfléchit» la trajectoire après le premier passage $\tau_a=\inf\{t:W_t=a\}$.
Cette transformation envoie bijectivement les trajectoires qui franchissent $a$ et finissent \emph{sous} $a$ vers des trajectoires qui ne franchissent pas mais finissent \emph{au-dessus} de $a$.
On obtient alors une relation de symétrie sur les probabilités, qui se traduit en gaussien par la formule ci-dessus.

\paragraph{Lien avec finance.}
Pour une barrière (knock-out/knock-in), la probabilité de franchissement d'un niveau intervient directement dans le prix.
En Black--Scholes, on transporte ce type de résultat via la transformation logarithmique du GBM.

\subsection*{(6) Intégrale d'Itô : construction rigoureuse en $L^2$ (version cours)}
On veut définir $\int_0^T H_t\,\dd W_t$ pour des processus adaptés $H$.
La construction standard (niveau M1) suit trois étapes.

\paragraph{Étape 1 : processus simples (prévisibles en escalier).}
On appelle \emph{processus simple prévisible} sur $[0,T]$ toute forme
\[
H_t=\sum_{k=0}^{n-1} H_{t_k}\,\1_{(t_k,t_{k+1}]}(t),
\qquad 0=t_0<\dots<t_n=T,
\]
où chaque $H_{t_k}$ est $\F_{t_k}$-mesurable (donc «décidé» au début de l'intervalle).
On définit alors
\[
\int_0^T H_t\,\dd W_t\ \triangleq\ \sum_{k=0}^{n-1} H_{t_k}\bigl(W_{t_{k+1}}-W_{t_k}\bigr).
\]

\paragraph{Étape 2 : isométrie d'Itô (preuve sur les simples).}
Pour un tel $H$, on a
\[
\E\!\left[\left(\int_0^T H_t\,\dd W_t\right)^2\right]
=\E\!\left[\int_0^T H_t^2\,\dd t\right].
\]
\emph{Preuve (calcul explicite).}
Écrire $I=\sum_k H_{t_k}\Delta W_k$ avec $\Delta W_k=W_{t_{k+1}}-W_{t_k}$.
Alors
\[
\E[I^2]=\E\!\left[\sum_{k,\ell} H_{t_k}H_{t_\ell}\Delta W_k\Delta W_\ell\right].
\]
Pour $k\neq \ell$, les incréments sont indépendants, centrés, et l'un des deux est indépendant de l'information contenant l'autre coefficient : les termes croisés ont espérance nulle.
Donc
\[
\E[I^2]=\sum_k \E\!\left[H_{t_k}^2\,\E[(\Delta W_k)^2\mid \F_{t_k}]\right]
=\sum_k \E[H_{t_k}^2]\,(t_{k+1}-t_k)
=\E\!\left[\int_0^T H_t^2\,\dd t\right].
\]

\paragraph{Étape 3 : complétion en $L^2$.}
On définit l'espace
\[
\mathcal{H}^2=\left\{H\ \text{prévisible}:\ \E\left[\int_0^T H_t^2\,\dd t\right]<\infty\right\}.
\]
Les processus simples sont denses dans $\mathcal{H}^2$ pour la norme $\|H\|_{2}^2=\E[\int_0^T H_t^2\dd t]$.
On choisit une suite $(H^n)$ simple qui converge vers $H$ dans cette norme, et on pose
\[
\int_0^T H_t\,\dd W_t \ \triangleq\ \lim_{n\to\infty} \int_0^T H^n_t\,\dd W_t
\quad \text{dans } L^2.
\]
L'isométrie garantit que la limite existe et ne dépend pas de l'approximation choisie.

\paragraph{Conséquence : martingale et variance.}
Pour $H\in\mathcal{H}^2$, le processus
\[
M_t=\int_0^t H_s\,\dd W_s
\]
est une martingale, et son «énergie» est
\[
\E[M_T^2]=\E\!\left[\int_0^T H_t^2\,\dd t\right].
\]
En finance, cette identité est la brique qui relie un contrôle quadratique (risque) aux expositions instantanées $H_t$.

\subsection*{(7) Itô vs Stratonovich (en une page) : convention et sens économique}
Deux conventions d'intégration apparaissent souvent :
\begin{itemize}[leftmargin=*]
\item Itô : on gèle l'intégrande au début de l'intervalle $(t_k,t_{k+1}]$.
\item Stratonovich : on gèle «au milieu» (ou plus exactement on prend une limite symétrique).
\end{itemize}
En finance, Itô est naturel pour un raisonnement \emph{non anticipatif} : une stratégie de trading ne peut pas dépendre de l'incrément futur $\Delta W_k$ sur l'intervalle.
Stratonovich est plus fréquent en physique (modélisation de bruits «lisses» approchés).
La conversion (en dimension 1) est :
\[
\int_0^T H_t\circ \dd W_t
=\int_0^T H_t\,\dd W_t+\frac12\int_0^T \frac{\partial H_t}{\partial x}\,\dd t
\quad \text{(si }H_t=H(t,X_t)\text{)}.
\]
Le terme $\tfrac12$ est la «correction» qui renvoie au chapitre Variation quadratique.

\subsection*{(8) Brownien multi-dimensionnel et corrélations (utile pour Heston/SABR)}
Un brownien $d$-dimensionnel $\bm{W}_t=(W_t^{(1)},\dots,W_t^{(d)})$ a des composantes browniens et des covariations
\[
[W^{(i)},W^{(j)}]_t=\rho_{ij}t,
\]
où la matrice $\rho$ (symétrique, définie positive) code les corrélations instantanées.
Une représentation standard est :
\[
\bm{W}_t = L\bm{B}_t,
\]
où $\bm{B}$ est un brownien à composantes indépendantes et $L$ une matrice telle que $LL^\top=\rho$ (décomposition de Cholesky).
Cette écriture est cruciale dès qu'on couple prix et variance (Heston) ou prix et volatilité (SABR).


\section{C. Dérivations (ultra détaillées)}
\subsection*{De $W_t$ à la dynamique de prix (GBM)}
Le modèle de Black--Scholes suppose une dynamique (sous la mesure risque-neutre $\Qmeas$) :
\[
\frac{\dd S_t}{S_t}=(r-q),\dd t + \sigma,\dd W_t^{\Qmeas}.
\]
On le réécrit :
\[
\dd S_t = (r-q)S_t,\dd t + \sigma S_t,\dd W_t.
\]
L'objet \emph{non classique} est $\dd W_t$ : ce n'est pas une dérivée ordinaire, mais un incrément brownien.

\subsection*{Solution exacte discrétisée}
On sait (via Itô, démontré au chapitre Itô) que la solution est :
\[
S_t = S_0\exp\Big((r-q-\tfrac12\sigma^2)t+\sigma W_t\Big).
\]
Pour une discrétisation $t_{n+1}=t_n+\Delta t$,
\[
\ln S_{t_{n+1}}-\ln S_{t_n}=(r-q-\tfrac12\sigma^2)\Delta t+\sigma (W_{t_{n+1}}-W_{t_n}),
\]
et comme $W_{t_{n+1}}-W_{t_n}\sim\Normal(0,\Delta t)$, on simule :
\[
W_{t_{n+1}}-W_{t_n}=\sqrt{\Delta t},Z_n,\quad Z_n\sim\Normal(0,1),
\]
d'où
\[
S_{t_{n+1}} = S_{t_n}\exp\Big((r-q-\tfrac12\sigma^2)\Delta t+\sigma\sqrt{\Delta t},Z_n\Big).
\]

\subsection*{Loi de $S_T$ (log-normale) et moments}
La formule fermée donne directement la loi :
\[
\ln S_T = \ln S_0 + \Big(r-q-\tfrac12\sigma^2\Big)T + \sigma W_T.
\]
Comme $W_T\sim\Normal(0,T)$, on a
\[
\ln S_T \sim \Normal\Big(\ln S_0 + (r-q-\tfrac12\sigma^2)T,\ \sigma^2 T\Big),
\]
donc $S_T$ est \textbf{log-normal}.

\paragraph{Moments (taux constants).}
Sous la mesure risque-neutre (drift $r-q$), on retrouve le forward :
\[
\E^{\Qmeas}[S_T]=S_0e^{(r-q)T}=F_{0,T}.
\]
Et la variance est
\[
\Var^{\Qmeas}(S_T)=S_0^2e^{2(r-q)T}\big(e^{\sigma^2T}-1\big),
\]
ce qui met en évidence l'effet \emph{exponentiel} de $\sigma^2T$.

\subsection*{Exact vs Euler : pourquoi l'exponentielle est préférable}
Une discrétisation Euler directe sur $S$,
\[
S_{t_{n+1}}=S_{t_n} + (r-q)S_{t_n}\Delta t + \sigma S_{t_n}\Delta W_n,
\]
peut produire des valeurs négatives (car l'approximation est additive).
À l'inverse, la discrétisation \emph{exacte} via $\ln S$ garantit $S_{t_{n+1}}>0$ et reproduit la loi marginale correcte à chaque pas.

% =========================================================
% Compléments M1 — Brownien (C. Dérivations)
% =========================================================

\subsection*{Dérivation guidée : résoudre le GBM sans «recette»}
On repart de la SDE (sous $\Qmeas$) :
\[
\frac{\dd S_t}{S_t}=(r-q)\,\dd t+\sigma\,\dd W_t.
\]
L'idée est de transformer une dynamique multiplicative en dynamique additive via le logarithme.
Posons $X_t=\ln S_t$.
Au chapitre Itô, on montre la formule (ici on l'utilise) :
\[
\dd \ln S_t = \frac{1}{S_t}\dd S_t-\frac12\frac{1}{S_t^2}(\dd S_t)^2.
\]
Or $\dd S_t=(r-q)S_t\dd t+\sigma S_t\dd W_t$, donc
\[
\frac{1}{S_t}\dd S_t=(r-q)\dd t+\sigma \dd W_t,
\qquad
(\dd S_t)^2=(\sigma S_t)^2(\dd W_t)^2=\sigma^2 S_t^2\,\dd t,
\]
et ainsi
\[
\dd \ln S_t
=(r-q)\dd t+\sigma \dd W_t-\tfrac12\sigma^2\dd t
=\bigl(r-q-\tfrac12\sigma^2\bigr)\dd t+\sigma \dd W_t.
\]
On a donc une dynamique \textbf{linéaire} pour $X_t=\ln S_t$ :
\[
X_t=X_0+\int_0^t \bigl(r-q-\tfrac12\sigma^2\bigr)\dd u+\sigma\int_0^t \dd W_u
=\ln S_0+\bigl(r-q-\tfrac12\sigma^2\bigr)t+\sigma W_t.
\]
En exponentiant :
\[
\boxed{
S_t=S_0\exp\Big(\bigl(r-q-\tfrac12\sigma^2\bigr)t+\sigma W_t\Big).
}
\]
\paragraph{Point pédagogique.}
La présence de $-\tfrac12\sigma^2 t$ n'est pas un «détail» : c'est la correction d'Itô (convexité) qui rend $S_t$ log-normal et fait que $\E[S_t]=S_0e^{(r-q)t}$.

\subsection*{Du modèle continu aux rendements discrets (annualisation de la volatilité)}
Sur un pas $\Delta t$, on obtient
\[
\ln\frac{S_{t+\Delta t}}{S_t}
=\bigl(r-q-\tfrac12\sigma^2\bigr)\Delta t+\sigma (W_{t+\Delta t}-W_t).
\]
Or $W_{t+\Delta t}-W_t\sim\Normal(0,\Delta t)$, donc
\[
\ln\frac{S_{t+\Delta t}}{S_t}\sim
\Normal\Big(\bigl(r-q-\tfrac12\sigma^2\bigr)\Delta t,\ \sigma^2\Delta t\Big).
\]
\paragraph{Conséquence.}
L'écart-type d'un log-return sur $\Delta t$ est $\sigma\sqrt{\Delta t}$.
Si $\Delta t$ est en \emph{années} (par exemple $\Delta t=1/252$), on retrouve l'heuristique «vol journalière $\approx \sigma/\sqrt{252}$».

\subsection*{Exponential martingale (outil clé pour Girsanov)}
Un objet qui revient en changement de mesure est la martingale exponentielle :
\[
Z_t=\exp\Big(\theta W_t-\tfrac12\theta^2 t\Big),\qquad \theta\in\R.
\]
\paragraph{Pourquoi c'est une martingale ?}
On sait que $W_t-W_s\sim\Normal(0,t-s)$ indépendant de $\F_s$.
Donc, pour $s<t$,
\[
\E[Z_t\mid \F_s]
=Z_s\,\E\!\left[\exp\Big(\theta(W_t-W_s)-\tfrac12\theta^2(t-s)\Big)\right]
=Z_s,
\]
car si $X\sim\Normal(0,t-s)$ alors $\E[e^{\theta X}]=\exp(\tfrac12\theta^2(t-s))$.
Cette identité est l'ombre portée du fait que «gaussien + exponentielle = calcul exact», ce qui rend le changement de mesure tractable en BS.

\subsection*{Exercices corrigés (Brownien, intégrale d'Itô, GBM)}
\paragraph{Exercice 1 --- loi d'une intégrale d'Itô à intégrande déterministe.}
Soit $h:[0,T]\to\R$ déterministe, mesurable, et $\int_0^T h(t)^2\,\dd t<\infty$.
Montrer que
\[
I=\int_0^T h(t)\,\dd W_t
\]
est une variable gaussienne centrée, et calculer sa variance.

\emph{Solution.}
On approxime $h$ par des fonctions en escalier $h^n$.
Pour $h^n$, on a
\[
I_n=\int_0^T h^n(t)\,\dd W_t=\sum_k h^n(t_k)\Delta W_k,
\]
combinaison linéaire d'incréments gaussiens indépendants $\Rightarrow I_n$ est gaussien centré.
Par isométrie d'Itô,
\[
\Var(I_n)=\E[I_n^2]=\int_0^T (h^n(t))^2\,\dd t.
\]
Comme $h^n\to h$ dans $L^2([0,T])$, on a $I_n\to I$ dans $L^2(\Omega)$ et la variance converge vers $\int_0^T h(t)^2\dd t$.
Une limite $L^2$ de gaussiennes centrées (en fait, convergence en loi) est gaussienne, donc
\[
I\sim \Normal\Big(0,\ \int_0^T h(t)^2\,\dd t\Big).
\]

\paragraph{Exercice 2 --- moments de $S_T$ sous GBM.}
Sous $\Qmeas$, $S_T=S_0\exp((r-q-\tfrac12\sigma^2)T+\sigma W_T)$.
\begin{enumerate}[leftmargin=*]
\item Calculer $\E^{\Qmeas}[S_T]$.
\item Calculer $\Var^{\Qmeas}(S_T)$.
\end{enumerate}

\emph{Solution.}
On utilise que $W_T\sim\Normal(0,T)$, donc $\E[e^{\sigma W_T}]=e^{\frac12\sigma^2T}$.
\[
\E[S_T]
=S_0e^{(r-q-\frac12\sigma^2)T}\E[e^{\sigma W_T}]
=S_0e^{(r-q-\frac12\sigma^2)T}e^{\frac12\sigma^2T}
=S_0e^{(r-q)T}.
\]
Pour la variance, on calcule $\E[S_T^2]$ :
\[
S_T^2=S_0^2\exp\Big(2(r-q-\tfrac12\sigma^2)T+2\sigma W_T\Big),
\]
donc $\E[e^{2\sigma W_T}]=e^{2^2\sigma^2T/2}=e^{2\sigma^2T}$ et
\[
\E[S_T^2]=S_0^2e^{2(r-q-\frac12\sigma^2)T}e^{2\sigma^2T}
=S_0^2e^{2(r-q)T}e^{\sigma^2T}.
\]
Ainsi
\[
\Var(S_T)=\E[S_T^2]-\E[S_T]^2
=S_0^2e^{2(r-q)T}\bigl(e^{\sigma^2T}-1\bigr).
\]

\paragraph{Exercice 3 --- probabilité de franchissement (principe de réflexion).}
Soit $a>0$. Montrer que
\[
\Prob\Big(\sup_{0\le t\le T} W_t\ge a\Big)=2\Prob(W_T\ge a).
\]
\emph{Solution (idée).}
On note $\tau_a=\inf\{t:W_t=a\}$ et on considère l'application qui, pour une trajectoire franchissant $a$ avant $T$ et finissant en dessous de $a$, réfléchit la partie après $\tau_a$.
On construit une bijection entre cet ensemble et l'ensemble des trajectoires finissant au-dessus de $a$.
On obtient
\[
\Prob(M_T\ge a)=\Prob(W_T\ge a)+\Prob(M_T\ge a,\ W_T<a)
=\Prob(W_T\ge a)+\Prob(W_T\ge a)=2\Prob(W_T\ge a).
\]


\begin{comment}
\section{D. Numérique / algorithmes (stabilité, complexité, choix)}
\begin{itemize}[leftmargin=*]
\item La simulation via exponentielle est \textbf{stable} (garantit $S_t>0$).
\item Complexité : $O(N_{\text{paths}}\times N_{\text{steps}})$.
\item Le choix $Z\sim\Normal(0,1)$ est le pont concret entre théorie (brownien) et code (numpy RNG).
\end{itemize}

\end{comment}
\begin{comment}
\section{E. Implémentation dans le repo (fichiers, fonctions, pipeline, paramètres)}
\subsection*{Traduction directe de la formule GBM}
Le simulateur le plus lisible est \codeid{simulate_gbm_paths} dans \filepath{app/model/options/core/trees.py}. Extrait (raccourci) :

\begin{lstlisting}[caption={Simulation GBM (extrait) : \filepath{app/model/options/core/trees.py}, fonction \codeid{simulate_gbm_paths}.},label={lst:gbm}]
def simulate_gbm_paths(S0, r, q, sigma, T, steps, n_paths, seed=None):
dt = float(T) / max(int(steps), 1)
drift = (float(r) - float(q) - 0.5 * float(sigma) ** 2) * dt
vol = float(sigma) * math.sqrt(dt)
rng = np.random.default_rng(seed)
Z = rng.standard_normal((int(n_paths), max(int(steps), 1)))
increments = drift + vol * Z
log_paths = np.cumsum(increments, axis=1)
S = float(S0) * np.exp(log_paths)
S = np.concatenate([np.full((int(n_paths), 1), float(S0)), S], axis=1)
return S
\end{lstlisting}

\paragraph{Lecture ligne à ligne (lien formule $\leftrightarrow$ code).}
\begin{itemize}[leftmargin=*]
\item \codeid{dt} implémente $\Delta t=T/N$.
\item \codeid{drift} implémente $(r-q-\frac12\sigma^2)\Delta t$.
\item \codeid{Z} est la matrice des $Z_n\sim\Normal(0,1)$.
\item \codeid{increments = drift + vol * Z} implémente $(r-q-\frac12\sigma^2)\Delta t+\sigma\sqrt{\Delta t},Z_n$.
\item \codeid{log_paths = cumsum(increments)} somme les incréments de $\ln S$.
\item \codeid{S = S0 * exp(log_paths)} revient à $S=\exp(\ln S)$.
\end{itemize}

\subsection*{Où l'app l'utilise}
\begin{itemize}[leftmargin=*]
\item Options exotiques MC (Asian, lookback, barrières, cliquet) : \filepath{app/model/options/core/trees.py} (\codeid{price_asian_mc}, \codeid{price_lookback_mc}, \ldots).
\item MC unifié (European/Bermudan/American via LSMC) : \filepath{app/model/options/service.py} (\codeid{price_option_mc_unified}) appelle \codeid{price_mc_lsmc} dans \filepath{app/model/options/mc_engine.py}.
\item Calibration rough/Volterra : simulateurs dédiés (non-GBM) dans \filepath{app/model/volatility_models/rbergomi/simulator.py} et \filepath{app/model/volatility_models/volterra/simulator.py} utilisent aussi des \codeid{standard_normal}.
\end{itemize}

\section{F. Exemple end-to-end (utilisateur $\to$ calcul)}
\begin{example}[Asian arithmetic (MC) dans l'app]
Dans l'UI Options, un panneau «Asian (arith)» est disponible via \filepath{app/controller/options_controller.py} qui exporte \codeid{price_asian_mc} depuis \filepath{app/model/options/core/trees.py}. Quand l'utilisateur clique «Compute», le flux conceptuel est :
\[
(S_0,K,T,\sigma,r,q)\xrightarrow{\text{simulate}} (S^{(i)}_{t_j}) \xrightarrow{\text{moyenne}} \bar S^{(i)} \xrightarrow{\text{payoff}} (\bar S^{(i)}-K)^+ \xrightarrow{\E,\text{disc}} \hat V.
\]
\end{example}

\section{G. Limites \& extensions crédibles}
\begin{itemize}[leftmargin=*]
\item Le GBM ignore smile/skew : la volatilité implicite observée varie avec $K$ et $T$. D'où la partie calibration/surface.
\item Le MC naïf converge en $O(1/\sqrt{N})$ : variance reduction (antithétiques, control variates) est une extension naturelle (\NonImplem\ explicitement : aucune implémentation dédiée repérée).
\end{itemize}

\end{comment}
\begin{jurybox}
Le brownien est une brique de base : il explique l'échelle $\sqrt{\Delta t}$ du bruit (donc la forme $\Delta W\approx \sqrt{\Delta t}\,Z$) et structure toutes les diffusions utilisées en pricing et en simulation Monte Carlo.
\end{jurybox}

% ---------------------------------------------------------
\chapter{Variation quadratique}

\section{A. Problème financier}
Pourquoi une dérivation (Itô) diffère-t-elle du calcul classique ? Parce qu'un brownien a des variations «trop rugueuses» : la somme des carrés des incréments ne tend pas vers 0 mais vers le temps.

\section{B. Théorie}
\begin{definition}[Variation quadratique]
Soit une partition $\pi={0=t_0<\cdots<t_n=T}$. La variation quadratique d'un processus $X$ sur $[0,T]$ est (si la limite existe) :
\[
[X]T=\lim{|\pi|\to 0}\sum_{k=0}^{n-1}(X_{t_{k+1}}-X_{t_k})^2.
\]
\end{definition}

\begin{proposition}[Cas du brownien]
Pour un brownien $W$, on a $[W]_T=T$ (p.s.).
\end{proposition}

\subsection*{Intuition}
Sur un intervalle $\Delta t$, l'incrément $W_{t+\Delta t}-W_t$ est de taille typique $\sqrt{\Delta t}$. Son carré est donc typiquement $\Delta t$. En sommant $n\simeq T/\Delta t$ termes, on obtient un total $\simeq T$.

\qa
{C'est la raison mathématique de la règle \emph{clé} : $(\dd W)^2=\dd t$ en calcul d'Itô, utilisée pour obtenir les solutions exactes et la PDE de Black--Scholes.}
{Hypothèse : bruit brownien (gaussien, continu).}
{Avec sauts ou volatilité rough, la structure de variation change ; les règles d'Itô se généralisent mais ne sont plus «$(\dd W)^2=\dd t$» au sens classique.}
{Justifie le terme correctif d'Itô et, très concrètement, l'usage des schémas log-normaux (exponentiels) pour préserver la positivité et la bonne variance en diffusion.}

\subsection*{Variation totale vs variation quadratique}
Deux notions sont souvent confondues :
\begin{itemize}[leftmargin=*]
\item la \textbf{variation totale} $\mathrm{Var}(X;[0,T])=\sup_\pi \sum_k |\Delta X_k|$ mesure l'«accumulation en valeur absolue» ;
\item la \textbf{variation quadratique} $[X]_T=\lim_{|\pi|\to 0}\sum_k (\Delta X_k)^2$ mesure l'«accumulation d'énergie».
\end{itemize}
Un processus peut avoir une variation totale infinie mais une variation quadratique finie (c'est le cas du brownien).

\begin{proposition}[Cas de variation finie]
Si $X$ est de variation finie sur $[0,T]$ (par exemple $X_t=\int_0^t a_u\,\dd u$), alors $[X]_T=0$.
\end{proposition}

\begin{proposition}[Cas des processus d'Itô]
Si $X_t=X_0+\int_0^t a_u\,\dd u + \int_0^t b_u\,\dd W_u$ avec $a,b$ adaptés et carrés-intégrables, alors
\[
[X]_t=\int_0^t b_u^2\,\dd u.
\]
\end{proposition}

\subsection*{Covariation (utile en dimension $d$)}
Pour deux processus $X,Y$, on définit la covariation
\[
[X,Y]_T=\lim_{|\pi|\to 0}\sum_k (X_{t_{k+1}}-X_{t_k})(Y_{t_{k+1}}-Y_{t_k}).
\]
En particulier, si $(W^{(1)},W^{(2)})$ est un brownien corrélé avec corrélation $\rho$, alors $[W^{(1)},W^{(2)}]_t=\rho t$.

% =========================================================
% Compléments M1 — Variation quadratique (B. Théorie)
% =========================================================

\subsection*{(1) Propriétés algébriques : pourquoi on parle de «crochet»}
La covariation $[X,Y]_t$ (quand elle existe) se comporte comme une forme bilinéaire :
\begin{itemize}[leftmargin=*]
\item \textbf{Symétrie :} $[X,Y]_t=[Y,X]_t$.
\item \textbf{Bilinéarité :} $[aX+bY,Z]_t=a[X,Z]_t+b[Y,Z]_t$.
\item \textbf{Lien avec la variation quadratique :} $[X]_t=[X,X]_t$.
\end{itemize}
Une identité à retenir (utile en multi-facteurs) est :
\[
[X+Y]_t=[X]_t+[Y]_t+2[X,Y]_t.
\]
Elle est l'analogue stochastique de $(x+y)^2=x^2+y^2+2xy$, mais ici c'est un résultat \emph{au niveau des limites de partitions}.

\subsection*{(2) Variation finie vs bruit : séparation structurelle}
Un message fondamental (et exploité partout en calcul d'Itô) est :
\begin{center}
\emph{La variation quadratique «voit» la partie bruitée (martingale) mais pas la partie de variation finie.}
\end{center}
Plus précisément :
\begin{itemize}[leftmargin=*]
\item Si $A$ est de variation finie, alors $[A]_t=0$.
\item Si $M$ est une martingale continue de type Itô, alors $[M]_t$ est typiquement non nul.
\end{itemize}
Pour un processus d'Itô
\[
X_t=X_0+\int_0^t a_s\,\dd s+\int_0^t b_s\,\dd W_s,
\]
on obtient (chapitre courant) :
\[
[X]_t=\int_0^t b_s^2\,\dd s,
\]
ce qui formalise l'idée : \emph{seul le coefficient de diffusion $b_s$ crée de la variation quadratique}.

\subsection*{(3) Crochet d'une intégrale d'Itô : règle de calcul à retenir}
Soit $M_t=\int_0^t H_s\,\dd W_s$ avec $H\in\mathcal{H}^2$.
Alors :
\[
[M]_t=\int_0^t H_s^2\,\dd s,
\qquad
[M,W]_t=\int_0^t H_s\,\dd s.
\]
\paragraph{Intuition.}
Sur un petit pas $\Delta t$, $\Delta M\approx H\,\Delta W$ donc
$(\Delta M)^2\approx H^2(\Delta W)^2\approx H^2\Delta t$,
et $\Delta M\,\Delta W\approx H(\Delta W)^2\approx H\Delta t$.
La théorie (limite) rend cette intuition exacte.

\subsection*{(4) Table d'Itô : comment manipuler les différentiels}
Le calcul stochastique se fait souvent avec une notation différentielle «comme si» :
\[
\dd X_t = a_t\,\dd t + b_t\,\dd W_t.
\]
La règle opérationnelle est la table d'Itô :
\[
(\dd t)^2=0,\qquad \dd t\,\dd W_t=0,\qquad (\dd W_t)^2=\dd t.
\]
En dimension $d$ avec corrélation instantanée $\rho_{ij}$ :
\[
\dd W_t^{(i)}\,\dd W_t^{(j)}=\rho_{ij}\,\dd t.
\]
\paragraph{Pourquoi c'est cohérent ?}
Parce que l'échelle est $\dd W=O(\sqrt{\dd t})$. Donc $(\dd W)^2$ est de l'ordre $\dd t$ (à conserver), alors que $\dd t\,\dd W$ est de l'ordre $\dd t^{3/2}$ (à négliger) et $(\dd t)^2$ de l'ordre $\dd t^2$ (négligeable).

\subsection*{(5) Au-delà du Brownien : semimartingales (message M1)}
Dans un cadre général, le calcul d'Itô s'applique aux \emph{semimartingales} :
\[
X_t=X_0+A_t+M_t,
\]
où $A$ est de variation finie (drift) et $M$ une martingale locale continue.
La variation quadratique s'identifie alors au crochet de $M$ :
\[
[X]_t=[M]_t,
\]
ce qui permet de réutiliser la table d'Itô sur une classe très large de modèles continus.
Les modèles à sauts (Merton/Bates) demandent une extension (Itô avec sauts), mais l'idée «variation quadratique = bruit» reste centrale.


\section{C. Dérivations (ultra détaillées)}
\subsection*{Pourquoi $\sum (\Delta W)^2 \to T$}
Pour une partition uniforme $t_k=k\Delta t$, $\Delta W_k=W_{t_{k+1}}-W_{t_k}\sim\Normal(0,\Delta t)$ i.i.d. Donc
\[
\E[(\Delta W_k)^2]=\Delta t,\qquad \Var((\Delta W_k)^2)=2(\Delta t)^2.
\]
La somme $S_n=\sum_{k=0}^{n-1}(\Delta W_k)^2$ vérifie :
\[
\E[S_n]=n\Delta t=T,\qquad \Var(S_n)=n\cdot 2(\Delta t)^2=2T\Delta t\to 0.
\]
Donc $S_n\to T$ en $L^2$, donc en probabilité, et (avec un argument plus fin) p.s.

\subsection*{Lecture en termes de loi $\chi^2$ (intuition forte)}
Sur une partition uniforme, on peut écrire $\Delta W_k=\sqrt{\Delta t}\,Z_k$ avec $Z_k\sim\Normal(0,1)$ i.i.d. Alors
\[
S_n=\sum_{k=0}^{n-1}(\Delta W_k)^2 = \Delta t \sum_{k=0}^{n-1} Z_k^2.
\]
La somme $\sum Z_k^2$ suit une loi $\chi^2(n)$, de moyenne $n$ et variance $2n$.
Ainsi $S_n = \Delta t\,\chi^2(n)$ est concentré autour de $n\Delta t=T$ avec des fluctuations typiques de taille $\sqrt{2n}\Delta t = \sqrt{2T\Delta t}$, qui tendent vers 0.

\subsection*{Conséquence opératoire : la règle $(\dd W)^2=\dd t$}
Dans un calcul «au premier ordre», on retient les ordres de grandeur :
\[
\Delta W = O(\sqrt{\Delta t}),\qquad (\Delta W)^2 = O(\Delta t).
\]
Donc quand on développe une fonction non linéaire (Taylor), les termes en $(\Delta W)^2$ \emph{ne sont pas négligeables} : ils contribuent au terme en $\Delta t$, ce qui produit la correction d'Itô.

% =========================================================
% Compléments M1 — Variation quadratique (C. Dérivations)
% =========================================================

\subsection*{(1) De la convergence en $L^2$ à la convergence p.s. : une esquisse propre}
Dans le texte principal, on a montré sur une partition uniforme que
\[
Q_n=\sum_{k=0}^{n-1}(\Delta W_k)^2 \to T \quad \text{en }L^2.
\]
Pour obtenir une convergence \emph{presque sûre}, on peut utiliser une sous-suite dyadique.
Posons $n_m=2^m$ et $\Delta t_m=T/2^m$, et $Q_{n_m}$ la somme correspondante.
On a (calcul déjà fait) :
\[
\Var(Q_{n_m})=2T\Delta t_m = 2T\frac{T}{2^m}=\frac{2T^2}{2^m}.
\]
Par l'inégalité de Tchebychev, pour tout $\varepsilon>0$,
\[
\Prob(|Q_{n_m}-T|>\varepsilon)\le \frac{\Var(Q_{n_m})}{\varepsilon^2}
=\frac{2T^2}{\varepsilon^2}\,2^{-m}.
\]
La série $\sum_m \Prob(|Q_{n_m}-T|>\varepsilon)$ converge, donc par Borel--Cantelli :
\[
Q_{n_m}\to T \quad \text{p.s.}
\]
\paragraph{Remarque (niveau M1).}
Passer ensuite de la sous-suite dyadique à une suite générale demande un argument de stabilité (contrôle des raffinements). L'idée est que raffiner une partition n'ajoute qu'un terme dont l'espérance est contrôlée et dont la variance décroît avec le pas maximal.

\subsection*{(2) Covariation : calcul explicite en dimension 2}
Soient $W^{(1)}$ et $W^{(2)}$ deux browniens corrélés tels que
\[
\dd W_t^{(2)}=\rho\,\dd W_t^{(1)}+\sqrt{1-\rho^2}\,\dd B_t,
\]
où $B$ est un brownien indépendant de $W^{(1)}$.
Alors, en utilisant la table d'Itô :
\[
\dd[W^{(1)},W^{(2)}]_t=\dd W_t^{(1)}\,\dd W_t^{(2)}
=\rho(\dd W_t^{(1)})^2+\sqrt{1-\rho^2}\,\dd W_t^{(1)}\dd B_t
=\rho\,\dd t.
\]
En intégrant :
\[
\boxed{[W^{(1)},W^{(2)}]_t=\rho t.}
\]
Ce calcul est exactement ce qui se cache derrière la phrase : «corrélation instantanée $\rho$».

\subsection*{(3) Crochet d'un processus d'Itô : dérivation pas à pas}
Soit
\[
X_t=X_0+\int_0^t a_s\,\dd s+\int_0^t b_s\,\dd W_s.
\]
On note $A_t=\int_0^t a_s\,\dd s$ (variation finie) et $M_t=\int_0^t b_s\,\dd W_s$ (martingale continue).
Alors $X=A+M$ et
\[
[X]_t=[A+M]_t=[A]_t+[M]_t+2[A,M]_t.
\]
Or $[A]_t=0$ (variation finie), et $[A,M]_t=0$ (produit d'une variation finie et d'une martingale continue), donc
\[
[X]_t=[M]_t.
\]
Il reste à calculer $[M]_t$.
Sur une partition, $\Delta M_k\approx b_{t_k}\Delta W_k$, donc heuristiquement
$\sum (\Delta M_k)^2\approx \sum b_{t_k}^2(\Delta W_k)^2\to \int_0^t b_s^2\,\dd s$.
La preuve rigoureuse passe par approximation de $b$ par des processus simples et utilisation de l'isométrie d'Itô.
On obtient :
\[
\boxed{[X]_t=\int_0^t b_s^2\,\dd s.}
\]

\subsection*{(4) Exemple : variation quadratique d'un GBM et variance réalisée}
Sous GBM, $\dd S_t=(r-q)S_t\,\dd t+\sigma S_t\,\dd W_t$.
La partie martingale est $\int_0^t \sigma S_s\,\dd W_s$, donc
\[
[S]_t=\int_0^t \sigma^2 S_s^2\,\dd s.
\]
\paragraph{Lien finance.}
Si l'on observe des prix à haute fréquence et qu'on somme les carrés des rendements (réalisée),
on approxime une quantité de la forme $\int_0^t \sigma^2 S_s^2\,\dd s$ (ou, en log, $\int_0^t \sigma^2\,\dd s$).
C'est l'origine théorique des estimateurs de volatilité réalisée.

\subsection*{(5) Exercices corrigés (variation quadratique)}
\paragraph{Exercice 1 --- $W_t^2-t$ est une martingale.}
Montrer que $M_t=W_t^2-t$ est une martingale et en déduire que $[W]_t=t$.

\emph{Solution.}
Appliquer Itô à $f(x)=x^2$ :
\[
\dd(W_t^2)=2W_t\,\dd W_t+(\dd W_t)^2=2W_t\,\dd W_t+\dd t.
\]
Donc $\dd(W_t^2-t)=2W_t\,\dd W_t$ : c'est une intégrale d'Itô, donc une martingale.
On lit aussi que le terme correctif est précisément $(\dd W_t)^2=\dd t$, ce qui correspond à $[W]_t=t$.

\paragraph{Exercice 2 --- covariation linéaire.}
Soit $X_t=a t+b W_t$ avec $a,b\in\R$. Calculer $[X]_t$.

\emph{Solution.}
La partie variation finie $at$ a crochet nul, donc
\[
[X]_t=[bW]_t=b^2[W]_t=b^2 t.
\]

\paragraph{Exercice 3 --- crochet d'une intégrale.}
Soit $M_t=\int_0^t H_s\,\dd W_s$ avec $H\in\mathcal{H}^2$. Montrer que $[M]_t=\int_0^t H_s^2\,\dd s$.

\emph{Solution (esquisse).}
Pour $H$ simple, $M$ est une somme $\sum H_{t_k}\Delta W_k$, donc la somme des carrés $\sum (\Delta M_k)^2$ est une somme de $H_{t_k}^2(\Delta W_k)^2$ qui converge vers $\int H^2$.
On étend par densité aux $H$ généraux via l'isométrie d'Itô.


\begin{comment}
\section{D. Numérique / algorithmes}
En simulation, cela justifie :
\begin{itemize}[leftmargin=*]
\item que l'échelle correcte de l'incrément est \codeid{sqrt(dt)} ;
\item qu'un schéma d'Euler pour $S$ doit inclure les termes d'Itô si on manipule des fonctions non linéaires.
\end{itemize}

\end{comment}
\begin{comment}
\section{E. Implémentation dans le repo}
Même si le code n'écrit pas «$(\dd W)^2=\dd t$», c'est encodé implicitement par :
\begin{itemize}[leftmargin=*]
\item \codeid{Z = rng.standard_normal(...)} et \codeid{vol = sigmasqrt(dt)} dans \filepath{app/model/options/core/trees.py}.
\item La construction de processus rough/Volterra où la convolution de noyau est appliquée sur des incréments \codeid{dB = sqrt(dt)*z1} (ex. \filepath{app/model/volatility_models/rbergomi/simulator.py} et \filepath{app/model/volatility_models/volterra/simulator.py}).
\end{itemize}

\section{F. Exemple concret (chiffré)}
\begin{example}[Ordre de grandeur]
Prenons $T=1$ an et $\Delta t=1/252$. Un incrément brownien typique a un écart-type $\sqrt{1/252}\approx 0.063$. Son carré typique vaut $\approx 0.004$. En multipliant par 252 pas, on retombe sur $\approx 1$.
\end{example}

\section{G. Limites \& extensions}
\begin{itemize}[leftmargin=*]
\item Si l'on veut une dynamique avec sauts (Merton/Bates) ou rough (rBergomi), la variation quadratique de certains objets change : le code traite ces cas via d'autres modèles (voir Partie II).
\end{itemize}

\end{comment}
\begin{jurybox}
La variation quadratique est la justification rigoureuse de la règle de calcul d'Itô. Sans elle, on ne peut pas dériver correctement ni la solution GBM, ni la PDE de Black--Scholes.
\end{jurybox}

% ---------------------------------------------------------
\chapter{Lemme d'Itô (1D et multi-dim)}

\section{A. Problème financier}
Les produits dérivés sont des fonctions $V(t,S_t,\ldots)$ d'états aléatoires. Pour relier :
\begin{itemize}[leftmargin=*]
\item la dynamique du sous-jacent $S_t$,
\item la dynamique du prix de l'option $V_t=V(t,S_t)$,
\end{itemize}
il faut une règle de dérivation adaptée au bruit brownien.

\section{B. Théorie}
\begin{theorem}[Itô 1D]
Soit $X_t$ solution de $\dd X_t=\mu(t,X_t)\dd t+\sigma(t,X_t)\dd W_t$, et $f\in C^{1,2}$. Alors
\[
\dd f(t,X_t)=\partial_t f,\dd t+\partial_x f,\dd X_t+\tfrac12,\partial_{xx}f,(\dd X_t)^2,
\]
et comme $(\dd X_t)^2=\sigma^2(t,X_t)\dd t$ (via $(\dd W)^2=\dd t$), on obtient
\[
\dd f(t,X_t)=\Big(\partial_t f+\mu,\partial_x f+\tfrac12\sigma^2,\partial_{xx}f\Big)\dd t+\sigma,\partial_x f,\dd W_t.
\]
\end{theorem}

\begin{theorem}[Itô multi-dim (forme utile)]
Si $\bm{X}_t\in\R^d$ suit $\dd \bm{X}_t=\bm{\mu}\dd t+\Sigma,\dd \bm{W}_t$ et $f\in C^{1,2}$, alors
\[
\dd f(t,\bm{X}_t)=\partial_t f,\dd t+\nabla f^\top \dd \bm{X}_t+\tfrac12 \mathrm{Tr}!\big(\Sigma\Sigma^\top \nabla^2 f\big)\dd t.
\]
\end{theorem}

\subsection*{Intuition}
Le terme $\tfrac12\sigma^2 f_{xx}\dd t$ est la «correction de convexité» : une fonction convexe appliquée à un bruit a un drift positif (Jensen), et Itô quantifie exactement cet effet à l'ordre $\dd t$.

\qa
{Transformer une dynamique de prix (SDE) en dynamique de prix d'option ; dériver la PDE Black--Scholes ; obtenir des solutions fermées (GBM) ; justifier des schémas numériques.}
{Hypothèses de régularité ($C^{1,2}$) et bruit brownien continu.}
{Les payoffs usuels ne sont pas lisses (kink) : on travaille avec des solutions faibles (PDE de viscosité), des approximations lisses, ou des extensions (Tanaka et temps local) pour justifier rigoureusement les passages à la limite.}
{Outil omniprésent : Itô relie SDE $\leftrightarrow$ PDE (Feynman--Kac) et donne des contrôles de cohérence (drift/variance) utilisés dans les formules fermées et la conception de schémas numériques.}

\subsection*{Formule produit : l'«intégration par parties» stochastique}
Si $X_t$ et $Y_t$ sont deux processus d'Itô, la règle de dérivation n'est plus $\,\dd(XY)=X\,\dd Y + Y\,\dd X\,$ mais
\[
\dd(X_tY_t)=X_t\,\dd Y_t + Y_t\,\dd X_t + \dd[X,Y]_t,
\]
où $[X,Y]$ est la covariation (et $[W,W]_t=t$).
C'est la formule centrale derrière :
\begin{itemize}[leftmargin=*]
\item la dynamique d'un portefeuille auto-financé ;
\item la construction de stratégies de couverture (on élimine la partie martingale en choisissant $\Delta$).
\end{itemize}

\subsection*{Payoffs non lisses : pourquoi ça marche quand même (idée)}
Le payoff d'un call $(S-K)^+$ n'est pas $C^2$ en $S=K$. En pratique on peut :
\begin{itemize}[leftmargin=*]
\item soit raisonner via la PDE et des solutions faibles (viscosité) ;
\item soit approximer le payoff par une famille lisse et passer à la limite ;
\item soit utiliser une extension d'Itô (formule de Tanaka) qui introduit un terme de \emph{temps local} au point $K$.
\end{itemize}
Intuition : le «kink» du payoff concentre de la convexité au strike ; cette convexité se manifeste dans la couverture via $\Gamma$.

% =========================================================
% Compléments M1 — Lemme d'Itô (B. Théorie)
% =========================================================

\subsection*{(1) Forme intégrale : ce que la formule dit vraiment}
La version différentielle est très pratique, mais pour raisonner sans ambiguïté il est utile de retenir la \textbf{forme intégrale}.
Si $X$ suit
\[
X_t=X_0+\int_0^t \mu(s,X_s)\,\dd s+\int_0^t \sigma(s,X_s)\,\dd W_s
\]
et si $f\in C^{1,2}$, alors
\[
\boxed{
\begin{aligned}
f(t,X_t)=\;&f(0,X_0)
 +\int_0^t \partial_t f(s,X_s)\,\dd s
 +\int_0^t \partial_x f(s,X_s)\,\dd X_s\\
&+\frac12\int_0^t \partial_{xx} f(s,X_s)\,\dd [X]_s,
\end{aligned}}
\]
avec $[X]_s=\int_0^s \sigma(u,X_u)^2\,\dd u$.
En remplaçant $\dd X_s=\mu\,\dd s+\sigma\,\dd W_s$ et $\dd[X]_s=\sigma^2\,\dd s$, on obtient la forme classique :
\[
f(t,X_t)=f(0,X_0)
\!+\!\int_0^t\!\Big(\partial_t f+\mu\,\partial_x f+\tfrac12\sigma^2\partial_{xx} f\Big)(s,X_s)\,\dd s
\!+\!\int_0^t\!\sigma\,\partial_x f(s,X_s)\,\dd W_s.
\]
\paragraph{Lecture martingale.}
Le dernier terme est une intégrale d'Itô, donc une martingale (sous conditions). Le reste est un terme de drift (déterministe conditionnellement à la trajectoire).

\subsection*{(2) Pourquoi $C^{1,2}$ ? et que faire quand ce n'est pas vrai}
\paragraph{Rôle de $C^{1,2}$.}
Le développement de Taylor utilisé dans la preuve requiert une dérivée en temps et deux dérivées en espace pour contrôler les restes.

\paragraph{Payoffs non lisses en finance.}
Les payoffs vanilles ont des singularités :
\begin{itemize}[leftmargin=*]
\item Call : $\Phi(S)=(S-K)^+$ non $C^1$ en $S=K$.
\item Digital : $\Phi(S)=\1_{\{S>K\}}$ discontinu.
\end{itemize}
On contourne en pratique via :
\begin{itemize}[leftmargin=*]
\item PDE en sens faible / viscosité ;
\item approximation lisse (mollification) + passage à la limite ;
\item extensions (Tanaka, temps local) pour traiter des convexités concentrées.
\end{itemize}

\subsection*{(3) Itô multi-dimensionnel : notation propre et cas corrélé}
Soit $\bm{X}_t\in\R^d$ avec
\[
\dd \bm{X}_t=\bm{\mu}(t,\bm{X}_t)\,\dd t+\Sigma(t,\bm{X}_t)\,\dd \bm{W}_t,
\]
où $\bm{W}$ est un brownien $d$-dimensionnel à composantes indépendantes.
Alors, pour $f\in C^{1,2}$,
\[
\dd f(t,\bm{X}_t)
=\partial_t f\,\dd t+\nabla f^\top \dd \bm{X}_t+\frac12\mathrm{Tr}\!\big(\Sigma\Sigma^\top \nabla^2 f\big)\dd t.
\]
\paragraph{Cas corrélé.}
Si l'on part d'un brownien corrélé $\bm{\widetilde W}$ tel que $\dd \bm{\widetilde W}=L\,\dd \bm{W}$ avec $LL^\top=\rho$, on peut absorber la corrélation dans $\Sigma$.
Dans les modèles multi-facteurs, l'objet important est toujours $\Sigma\Sigma^\top$ : c'est lui qui porte les covariations.

\subsection*{(4) Itô--Tanaka (optionnel mais très utile pour les payoffs à kink)}
Pour une fonction convexe mais non $C^2$ (comme $x\mapsto |x|$), la formule d'Itô se prolonge via la \textbf{formule de Tanaka}.
Pour un semimartingale continu $X$ :
\[
|X_t|
=|X_0|+\int_0^t \mathrm{sgn}(X_s)\,\dd X_s + L_t^0(X),
\]
où $L_t^0(X)$ est le \emph{temps local} en 0.
\paragraph{Interprétation.}
Le temps local mesure combien de temps (au sens «densité de passage») la trajectoire passe près d'un niveau.
En finance, il apparaît naturellement quand on manipule des payoffs avec un kink (call/put) ou des barrières (densité de premier passage).
\paragraph{Message pédagogique.}
Même si l'on ne l'utilise pas directement pour pricer une vanilla, Tanaka explique pourquoi la convexité concentrée en $K$ fait apparaître des effets de couverture (Gamma) «localisés».

\subsection*{(5) Itô pour semimartingales : la version «générique» (message M1)}
Dans un cours de calcul stochastique, on retient que l'objet central est la semimartingale continue
\[
X_t=X_0+A_t+M_t,
\]
où $A$ est de variation finie et $M$ une martingale locale continue.
Alors, pour $f\in C^2$ (sans dépendance en $t$ pour simplifier),
\[
f(X_t)=f(X_0)+\int_0^t f'(X_s)\,\dd X_s+\frac12\int_0^t f''(X_s)\,\dd [X]_s.
\]
Cette forme est souvent la plus robuste conceptuellement : elle fait apparaître explicitement $\dd[X]_s$, donc l'endroit exact où la rugosité (variation quadratique) intervient.


\section{C. Dérivations (ultra détaillées)}
\subsection*{Dérivation (Taylor) et apparition de $(\dd W)^2$}
On écrit un développement de Taylor à l'ordre 2 :
\[
f(t+\dd t, X_t+\dd X)-f(t,X_t)
= f_t\dd t + f_x\dd X + \tfrac12 f_{xx}(\dd X)^2 + o(\dd t).
\]
Or $\dd X=\mu \dd t+\sigma \dd W$. En développant :
\[
(\dd X)^2 = (\mu \dd t)^2 + 2\mu\sigma,\dd t,\dd W + \sigma^2 (\dd W)^2.
\]
On utilise les ordres de grandeur :
\[
\dd t = O(\dd t),\quad \dd W = O(\sqrt{\dd t}),\quad (\dd W)^2=O(\dd t),\quad \dd t,\dd W=O(\dd t^{3/2}),\quad (\dd t)^2=O(\dd t^2).
\]
Donc à l'ordre $\dd t$, on néglige $(\dd t)^2$ et $\dd t,\dd W$, mais on \emph{garde} $\sigma^2(\dd W)^2$, et avec $[W]t=t$, on remplace $(\dd W)^2$ par $\dd t$. D'où le terme $\tfrac12\sigma^2 f{xx}\dd t$.

\subsection*{Exemple Itô : $f(x)=\ln x$ et GBM}
Sous GBM : $\dd S=(r-q)S,\dd t+\sigma S,\dd W$.
Pour $f(S)=\ln S$ :
\[
f_S=\frac1S,\qquad f_{SS}=-\frac1{S^2}.
\]
Itô :
\[
\dd \ln S = \frac1S\dd S +\tfrac12\Big(-\frac1{S^2}\Big)(\sigma^2 S^2)\dd t
= (r-q)\dd t+\sigma \dd W - \tfrac12\sigma^2 \dd t
\]
soit
\[
\dd \ln S = \big(r-q-\tfrac12\sigma^2\big)\dd t + \sigma,\dd W,
\]
ce qui justifie la solution exponentielle utilisée en simulation.

\subsection*{Exemple : formule produit à partir d'Itô}
Appliquons Itô à $f(x,y)=xy$ avec $(X_t,Y_t)$ un couple d'Itô :
\[
\dd X_t=a_t\,\dd t+b_t\,\dd W_t,\qquad \dd Y_t=c_t\,\dd t+d_t\,\dd W_t.
\]
Comme $f_x=y$, $f_y=x$, $f_{xy}=1$, et que $\dd X\,\dd Y=b_td_t\,\dd t$, on obtient
\[
\dd(X_tY_t)=X_t\,\dd Y_t + Y_t\,\dd X_t + b_td_t\,\dd t
          =X_t\,\dd Y_t + Y_t\,\dd X_t + \dd[X,Y]_t.
\]
Ce calcul explique pourquoi des termes «d'ordre 2» apparaissent dans les portefeuilles continus.

\subsection*{Exemple : puissance et moments de GBM}
Sous GBM, $S_t^p$ a une dynamique explicite.
En posant $f(s)=s^p$, on obtient
\[
\dd(S_t^p)=pS_t^{p}\Big((r-q)\dd t+\sigma\dd W_t\Big)+\tfrac12 p(p-1)\sigma^2 S_t^{p}\dd t,
\]
soit
\[
\dd(S_t^p)=S_t^p\Big(p(r-q)+\tfrac12 p(p-1)\sigma^2\Big)\dd t + p\sigma S_t^p\,\dd W_t.
\]
En prenant l'espérance sous $\Qmeas$, on retrouve des formules fermées pour $\E[S_T^p]$ (utile pour comprendre les moments et les contraintes de Carr--Madan/FFT).

% =========================================================
% Compléments M1 — Lemme d'Itô (C. Dérivations)
% =========================================================

\subsection*{(1) Preuve guidée (dimension 1) : où apparaît exactement le terme $\tfrac12 f_{xx}\sigma^2$}
On donne ici une preuve structurée (niveau M1) de la formule d'Itô en dimension 1, en suivant l'idée «partition + Taylor + variation quadratique».
Soit $X$ un processus d'Itô :
\[
X_t=X_0+\int_0^t a_s\,\dd s+\int_0^t b_s\,\dd W_s.
\]
Fixons une partition $\pi:0=t_0<\dots<t_n=t$ et écrivons, par télescopage,
\[
f(t,X_t)-f(0,X_0)=\sum_{k=0}^{n-1}\Big(f(t_{k+1},X_{t_{k+1}})-f(t_k,X_{t_k})\Big).
\]
On ajoute et retranche $f(t_{k+1},X_{t_k})$ :
\[
\Delta f_k
=\underbrace{f(t_{k+1},X_{t_k})-f(t_k,X_{t_k})}_{(\mathrm{I})}
\;+\;\underbrace{f(t_{k+1},X_{t_{k+1}})-f(t_{k+1},X_{t_k})}_{(\mathrm{II})}.
\]

\paragraph{(I) Variation en temps.}
Par Taylor en temps :
\[
(\mathrm{I})=\partial_t f(t_k,X_{t_k})\Delta t_k + r_k^{(t)},
\]
où $\sum_k r_k^{(t)}\to 0$ quand $|\pi|\to 0$ (régularité $C^{1,2}$).

\paragraph{(II) Variation en espace.}
Par Taylor en espace à l'ordre 2 :
\[
(\mathrm{II})
=\partial_x f(t_{k+1},X_{t_k})\Delta X_k
\frac12\partial_{xx}f(t_{k+1},X_{t_k})(\Delta X_k)^2+r_k^{(x)}.
\]

\paragraph{Contrôle des ordres de grandeur.}
On a
\[
\Delta X_k = \int_{t_k}^{t_{k+1}} a_s\,\dd s + \int_{t_k}^{t_{k+1}} b_s\,\dd W_s.
\]
Le premier terme est $O(\Delta t_k)$, le second est $O(\sqrt{\Delta t_k})$.
Donc $(\Delta X_k)^2$ est dominé par la partie brownienne et contribue à l'ordre $\Delta t_k$ (c'est le c\oe{}ur du raisonnement).

\paragraph{Passage à la limite : identification des trois termes.}
En sommant :
\begin{itemize}[leftmargin=*]
\item Les sommes des termes $\partial_t f\,\Delta t_k$ convergent vers $\int_0^t \partial_t f(s,X_s)\,\dd s$.
\item Les sommes des termes $\partial_x f\,\Delta X_k$ convergent vers $\int_0^t \partial_x f(s,X_s)\,\dd X_s$ (qui se décompose ensuite en drift + intégrale d'Itô).
\item La somme des termes $\frac12\partial_{xx}f\,(\Delta X_k)^2$ converge vers $\frac12\int_0^t \partial_{xx} f(s,X_s)\,\dd[X]_s$.
\end{itemize}
Le point crucial est le troisième item : comme $[X]_s=\int_0^s b_u^2\,\dd u$, on obtient le terme
\[
\frac12\int_0^t \partial_{xx}f(s,X_s)b_s^2\,\dd s,
\]
qui est \emph{absent} en calcul classique.
Les restes $\sum_k r_k^{(t)}$ et $\sum_k r_k^{(x)}$ tendent vers 0 par régularité.
On obtient la formule d'Itô.

\subsection*{(2) Exemple canonique : martingale exponentielle via Itô}
Reprenons
\[
Z_t=\exp\Big(\theta W_t-\tfrac12\theta^2 t\Big).
\]
Posons $f(t,w)=\exp(\theta w-\tfrac12\theta^2 t)$.
Alors
\[
\partial_t f=-\tfrac12\theta^2 f,\qquad \partial_w f=\theta f,\qquad \partial_{ww}f=\theta^2 f.
\]
Par Itô :
\[
\dd Z_t = \Big(\partial_t f+\tfrac12\partial_{ww}f\Big)\dd t + \partial_w f\,\dd W_t
=\Big(-\tfrac12\theta^2 f+\tfrac12\theta^2 f\Big)\dd t + \theta f\,\dd W_t
=\theta Z_t\,\dd W_t.
\]
Le drift est nul : $Z$ est une martingale. C'est exactement la forme utilisée dans Girsanov.

\subsection*{(3) Exemple finance : actualisation et produit stochastique}
Soit $B_t=e^{rt}$ le compte monétaire.
La variable actualisée $\widetilde S_t=e^{-(r-q)t}S_t$ doit être une martingale sous $\Qmeas$ en BS.
Vérifions-le par calcul :
\[
\dd S_t=(r-q)S_t\,\dd t+\sigma S_t\,\dd W_t^{\Qmeas},
\qquad
\dd\big(e^{-(r-q)t}\big)=-(r-q)e^{-(r-q)t}\,\dd t.
\]
En appliquant la formule produit $\,\dd(UV)=U\,\dd V+V\,\dd U+\dd[U,V]\,$ avec $U=e^{-(r-q)t}$ (variation finie, donc crochet nul), on obtient
\[
\dd\widetilde S_t
=e^{-(r-q)t}\dd S_t+S_t\,\dd(e^{-(r-q)t})
=e^{-(r-q)t}\sigma S_t\,\dd W_t^{\Qmeas},
\]
donc $\widetilde S$ est bien une martingale.

\subsection*{(4) Exercices corrigés (Itô)}
\paragraph{Exercice 1 --- SDE de $1/S_t$ sous GBM.}
Sous $\dd S_t=(r-q)S_t\,\dd t+\sigma S_t\,\dd W_t$, calculer la dynamique de $Y_t=1/S_t$.

\emph{Solution.}
Poser $f(s)=s^{-1}$, alors $f'(s)=-s^{-2}$ et $f''(s)=2s^{-3}$.
Itô :
\[
\dd Y_t=f'(S_t)\dd S_t+\tfrac12 f''(S_t)(\dd S_t)^2.
\]
Or $(\dd S_t)^2=\sigma^2 S_t^2\,\dd t$.
Donc
\[
\dd Y_t=-\frac{1}{S_t^2}\big((r-q)S_t\,\dd t+\sigma S_t\,\dd W_t\big)
\frac12\cdot 2\frac{1}{S_t^3}\sigma^2 S_t^2\,\dd t
=-(r-q)Y_t\,\dd t-\sigma Y_t\,\dd W_t+\sigma^2 Y_t\,\dd t,
\]
soit
\[
\boxed{\dd Y_t=(\sigma^2-(r-q))Y_t\,\dd t-\sigma Y_t\,\dd W_t.}
\]

\paragraph{Exercice 2 --- calcul de $\E[S_T^p]$ (rappel).}
Montrer à partir de la dynamique de $S_t^p$ que
\[
\E^{\Qmeas}[S_T^p]=S_0^p\exp\Big(p(r-q)T+\tfrac12 p(p-1)\sigma^2T\Big).
\]
\emph{Solution.}
On a déjà obtenu
\[
\dd(S_t^p)=S_t^p\Big(p(r-q)+\tfrac12 p(p-1)\sigma^2\Big)\dd t + p\sigma S_t^p\,\dd W_t.
\]
En prenant l'espérance, le terme martingale disparaît, donc
\[
\frac{\dd}{\dd t}\E[S_t^p]=\Big(p(r-q)+\tfrac12 p(p-1)\sigma^2\Big)\E[S_t^p],
\]
ODE résolue par exponentielle, avec condition initiale $\E[S_0^p]=S_0^p$.


\begin{comment}
\section{D. Numérique / algorithmes}
Deux traductions numériques :
\begin{itemize}[leftmargin=*]
\item \textbf{Solution exacte} (GBM) : exponentielle (stable, $S>0$) $\Rightarrow$ utilisée dans \filepath{app/model/options/core/trees.py}.
\item \textbf{PDE} : si on préfère résoudre la PDE plutôt que simuler, on utilise Crank--Nicolson (\filepath{app/model/options/engines/pricing.py}).
\end{itemize}

\end{comment}
\begin{comment}
\section{E. Implémentation dans le repo}
\begin{itemize}[leftmargin=*]
\item GBM exact : \filepath{app/model/options/core/trees.py}, \codeid{simulate_gbm_paths}.
\item Itô multi-dim (corrélations) : dans rBergomi, construction de $dW=\rho dB + \sqrt{1-\rho^2} dB^\perp$ (voir \filepath{app/model/volatility_models/rbergomi/simulator.py}, variables \codeid{dB}, \codeid{dW}).
\end{itemize}

\section{F. Exemple concret}
\begin{example}[Corrélation en code (rBergomi)]
Dans \filepath{app/model/volatility_models/rbergomi/simulator.py}, on construit deux bruits gaussiens \codeid{z1}, \codeid{z2} puis on corrèle via
\[
\dd W = \rho,\dd B + \sqrt{1-\rho^2},\dd B^\perp,
\]
ce qui est la traduction directe du fait que $\Corr(\dd W,\dd B)=\rho$.
\end{example}

\section{G. Limites \& extensions}
\begin{itemize}[leftmargin=*]
\item Itô suppose continuité : pour les sauts (Merton/Bates) on utilise des formules via fonctions caractéristiques (Partie II) plutôt qu'Itô pur.
\end{itemize}

\end{comment}
\begin{jurybox}
Itô est la «règle de dérivation» du monde stochastique. La correction $\tfrac12\sigma^2 f_{xx}$ est \emph{le} point conceptuel à comprendre : c'est elle qui rend GBM log-normal et qui mène à la PDE de pricing.
\end{jurybox}

% ---------------------------------------------------------
\chapter{Risk-neutral pricing, hedging et PDE (Black--Scholes \& Feynman--Kac)}

\section{A. Problème financier}
Une option européenne a un payoff $\Phi(S_T)$ à maturité. Objectif : déterminer un prix «sans arbitrage» $V_0$ aujourd'hui, cohérent avec la possibilité de couvrir (hedger) dynamiquement.

\section{B. Théorie}
\subsection*{Définition formelle : absence d'arbitrage et mesure risque-neutre}
Dans un marché frictionless complet (hypothèse BS), absence d'arbitrage $\Rightarrow$ existence d'une mesure $\Qmeas$ telle que le prix actualisé est une martingale :
\[
\tilde S_t = e^{-(r-q)t}S_t \text{ est une $\Qmeas$-martingale}.
\]
Pour un payoff $\Phi(S_T)$, le prix est :
\[
V_0 = e^{-rT}\E^{\Qmeas}[\Phi(S_T)\mid \F_0].
\]

\subsection*{Intuition financière}
La mesure $\Qmeas$ n'est pas «la vraie probabilité», c'est une transformation qui fait disparaître les primes de risque dans le drift. Le prix devient une espérance actualisée car on raisonne en \emph{numéraire} (monnaie à taux $r$).

\qa
{Relier couverture (réplication) $\leftrightarrow$ prix ; produire une PDE solvable numériquement ; justifier la formule BS.}
{Marché complet, trading continu, pas de coûts, volatilité constante, taux $r$ et dividend yield $q$ constants.}
{En rebalancement discret, avec coûts, sauts ou volatilité stochastique, la réplication parfaite échoue : on bascule vers marché incomplet (mesure risque-neutre non unique) et/ou vers des couvertures optimales (min-variance, utilité).}
{Cadre de référence : PDE de Black--Scholes, lien Feynman--Kac, et conception de schémas PDE (implicites/Crank--Nicolson) comme outils numériques de pricing et de stress.}

\subsection*{Portefeuilles auto-financés : l'hypothèse cachée des dérivations}
On considère un portefeuille composé du sous-jacent et du compte monétaire :
\[
X_t=\Delta_t S_t + \beta_t B_t.
\]
Dire que le portefeuille est \emph{auto-financé} signifie que toute variation de $X_t$ vient uniquement des variations de prix des actifs (on n'injecte pas de cash exogène).
Dans un modèle à dividendes continus ($q$), il faut aussi préciser si $S_t$ est un prix \emph{ex-dividend} : c'est ce point qui explique l'apparition de $(r-q)$ dans la PDE.

\subsection*{Théorème fondamental de l'évaluation des actifs (énoncé, intuition)}
Sans entrer dans toutes les subtilités d'intégrabilité, l'idée est :
\begin{itemize}[leftmargin=*]
\item \textbf{Absence d'arbitrage} $\Longleftrightarrow$ il existe au moins une mesure $\Qmeas$ équivalente à $\Pmeas$ sous laquelle les prix actualisés sont des martingales.
\item \textbf{Complétude} $\Longleftrightarrow$ cette mesure est unique (et alors tout payoff intégrable est réplicable).
\end{itemize}
Le modèle Black--Scholes est complet (une source de risque, un actif risqué), d'où l'unicité de $\Qmeas$ et la possibilité de couvrir parfaitement une option européenne.

\subsection*{Girsanov : comment le drift «réel» disparaît}
Sous $\Pmeas$, on écrit souvent
\[
\frac{\dd S_t}{S_t}=\mu\,\dd t+\sigma\,\dd W_t^{\Pmeas}.
\]
On définit le \emph{prix du risque de marché} $\lambda=\frac{\mu-(r-q)}{\sigma}$ et un nouveau brownien
\[
\dd W_t^{\Qmeas}=\dd W_t^{\Pmeas}+\lambda\,\dd t.
\]
Le théorème de Girsanov garantit l'existence d'une mesure $\Qmeas$ sous laquelle $W^{\Qmeas}$ est brownien, et la dynamique devient
\[
\frac{\dd S_t}{S_t}=(r-q)\,\dd t+\sigma\,\dd W_t^{\Qmeas}.
\]
Message pratique : \textbf{le pricing sans arbitrage dépend de $\sigma$ et de $(r-q)$, pas de $\mu$}.

% =========================================================
% Compléments M1 — Risk-neutral pricing (B. Théorie)
% =========================================================

\subsection*{(1) Numéraire : pourquoi «prix = espérance» n'est pas une magie}
Le point conceptuel est le suivant :
\begin{quote}
\emph{Un prix n'est une espérance que relativement à un choix de numéraire.}
\end{quote}
Soit $B_t$ un actif strictement positif (le numéraire). Typiquement $B_t=e^{rt}$ (compte monétaire).
On définit le prix \emph{actualisé} (ou numéraire-quotient) :
\[
\widetilde S_t \triangleq \frac{S_t}{B_t}.
\]
Le théorème fondamental de l'évaluation des actifs dit (version intuitive) :
\[
\text{NA} \Longrightarrow \exists\,\Qmeas\sim\Pmeas \text{ telle que } (\widetilde S_t) \text{ est une martingale sous }\Qmeas.
\]
Alors pour tout payoff intégrable $X_T$ payable à $T$, son prix à $t$ (dans un marché complet) est
\[
\boxed{
V_t = B_t\,\E^{\Qmeas}\!\left[\left.\frac{X_T}{B_T}\right|\F_t\right].
}
\]
\paragraph{Lecture.}
L'espérance porte sur un payoff \emph{en unité de numéraire}. On re-multiplie ensuite par $B_t$ pour revenir en monnaie.

\subsection*{(2) Mesure équivalente et dérivée de Radon--Nikodym : densité de changement de mesure}
Dire $\Qmeas\sim\Pmeas$ signifie : mêmes événements nuls (pas d'événement «impossible» sous l'une et possible sous l'autre).
On introduit la densité
\[
Z_T=\frac{\dd \Qmeas}{\dd \Pmeas}\Big|_{\F_T},
\qquad Z_T>0,\ \E^{\Pmeas}[Z_T]=1,
\]
et la formule de Bayes (conditionnelle) :
\[
\E^{\Qmeas}[Y\mid \F_t]
=\frac{\E^{\Pmeas}[Z_T Y\mid \F_t]}{\E^{\Pmeas}[Z_T\mid \F_t]}.
\]
\paragraph{Interprétation finance.}
$Z_T$ est proche d'une \emph{densité d'état} (state-price density) : elle repondère les scénarios pour refléter aversion au risque / prix du risque.
En modèle BS, $Z_T$ est une exponentielle de Brownien (cf. martingale exponentielle).

\subsection*{(3) Prix du risque de marché : séparer drift réel et drift de pricing}
Sous $\Pmeas$, on écrit
\[
\frac{\dd S_t}{S_t}=\mu\,\dd t+\sigma\,\dd W_t^{\Pmeas}.
\]
Sous $\Qmeas$ (numéraire $B_t=e^{rt}$), on veut
\[
\frac{\dd S_t}{S_t}=(r-q)\,\dd t+\sigma\,\dd W_t^{\Qmeas}.
\]
La différence $\mu-(r-q)$ est la \emph{prime de risque}. On définit
\[
\lambda \triangleq \frac{\mu-(r-q)}{\sigma},
\]
appelé \emph{market price of risk}.
Girsanov dit précisément comment fabriquer $\Qmeas$ telle que
\[
W_t^{\Qmeas}=W_t^{\Pmeas}+\int_0^t \lambda\,\dd s
\]
soit un brownien sous $\Qmeas$.
Message clé : $\mu$ est un paramètre «historique», alors que $(r-q)$ et $\sigma$ sont les paramètres de pricing en BS.

\subsection*{(4) Complétude et unicité de $\Qmeas$ : que signifie «réplication»}
Dans un marché frictionless :
\begin{itemize}[leftmargin=*]
\item \textbf{NA} (no arbitrage) $\Rightarrow$ existence d'au moins une mesure martingale équivalente (EMM).
\item \textbf{Complétude} $\Rightarrow$ unicité de l'EMM et réplication de tout payoff intégrable.
\end{itemize}
En Black--Scholes (un brownien, un actif risqué), le marché est complet : il y a exactement «autant de sources de risque que d'actifs risqués indépendants», ce qui permet de couvrir.
\paragraph{Message M1.}
La complétude est intimement liée au \emph{théorème de représentation des martingales} : toute martingale continue (adaptée à la filtration brownienne) peut s'écrire comme une intégrale d'Itô. C'est la forme mathématique de «toute incertitude peut être répliquée par trading continu».

\subsection*{(5) Marché incomplet : non-unicité et choix de prix}
Dès qu'on ajoute :
\begin{itemize}[leftmargin=*]
\item sauts (Merton/Bates),
\item volatilité stochastique non entièrement tradable (Heston),
\item frictions (coûts, contraintes, hedging discret),
\end{itemize}
le marché devient (souvent) incomplet : plusieurs mesures $\Qmeas$ sont possibles.
On doit alors choisir un critère (exemples) :
\begin{itemize}[leftmargin=*]
\item \emph{prix par utilité} (maximisation d'espérance d'utilité),
\item \emph{minimal martingale measure} (minimiser un risque résiduel),
\item \emph{calibration} : choisir une $\Qmeas$ (ou des paramètres) qui colle à la surface d'options observée.
\end{itemize}
Cela explique pourquoi, en pratique, «le modèle» est souvent \emph{identifié} par calibration plutôt que par estimation historique.

\subsection*{(6) Changement de numéraire : forward measure (utile en taux, et même en equity)}
Le choix du numéraire n'est pas unique.
Par exemple, en taux, on choisit souvent comme numéraire un zéro-coupon $P(t,T)$, ce qui définit la \emph{mesure $T$-forward}.
En equity avec dividendes, il est parfois naturel de travailler avec le numéraire «forward» $B_t e^{-qt}$.
La règle générale est :
\[
V_t = N_t\,\E^{\Qmeas^{N}}\!\left[\left.\frac{X_T}{N_T}\right|\F_t\right],
\]
où $\Qmeas^{N}$ est la mesure associée au numéraire $N$ (sous laquelle les prix divisés par $N$ sont des martingales).
Ce point clarifie beaucoup de dérivations (et évite des erreurs de drift).


\section{C. Dérivations (ultra détaillées)}
\subsection*{Dérivation complète de Black--Scholes via portefeuille couvert}
On suppose sous $\Pmeas$ : $\dd S=\mu S\dd t+\sigma S\dd W$ (le drift réel $\mu$ est inconnu/irrelevant pour le pricing BS).

Soit $V(t,S)$ le prix de l'option. Par Itô :
\[
\dd V = V_t \dd t + V_S \dd S + \tfrac12 V_{SS}(\dd S)^2
= \Big(V_t+\mu S V_S+\tfrac12\sigma^2 S^2 V_{SS}\Big)\dd t + \sigma S V_S \dd W.
\]
Construisons un portefeuille auto-financé :
\[
\Pi = V - \Delta S,
\]
où $\Delta$ est choisi pour annuler le terme aléatoire. On calcule :
\[
\dd \Pi = \dd V - \Delta\,\dd S.
\]
En substituant :
\[
\begin{aligned}
\dd \Pi &= \dd V - \Delta\,\dd S\\
&= \Big(V_t+\mu S V_S+\tfrac12\sigma^2 S^2 V_{SS}\Big)\dd t + \sigma S V_S \dd W\\
&\quad - \Delta(\mu S\dd t+\sigma S\dd W).
\end{aligned}
\]
Choisir $\Delta=V_S$ annule le terme en $\dd W$ :
\[
\dd \Pi = \Big(V_t+\tfrac12\sigma^2 S^2 V_{SS}\Big)\dd t.
\]
Maintenant, en absence d'arbitrage, un portefeuille sans risque doit rapporter le taux sans risque \emph{sur la valeur investie}. Ici, attention : la position en $S$ doit être financée et il existe un yield $q$ (dividende continu). La forme standard conduit à :
\[
\dd \Pi = r\Pi,\dd t = r(V - V_S S)\dd t.
\]
En égalant les deux expressions :
\[
V_t+\tfrac12\sigma^2 S^2 V_{SS}=r(V - S V_S).
\]
En réarrangeant, on obtient la PDE de Black--Scholes avec dividende :
\[
\boxed{V_t + (r-q)S V_S + \tfrac12\sigma^2 S^2 V_{SS} - rV = 0.}
\]
Condition terminale : $V(T,S)=\Phi(S)$.
\subsection*{Feynman--Kac (PDE $\leftrightarrow$ espérance)}
Le théorème de Feynman--Kac dit essentiellement : la solution de la PDE précédente est l'espérance risque-neutre actualisée. Intuition : on suit la trajectoire stochastique de $S_t$ sous $\Qmeas$ (drift $r-q$) et la martingale du prix actualisé force l'égalité PDE/espérance.

\subsection*{Feynman--Kac : énoncé (forme minimale) et idée de preuve}
Considérons un processus $S_t$ sous $\Qmeas$ :
\[
\dd S_t=(r-q)S_t\,\dd t+\sigma S_t\,\dd W_t^{\Qmeas}.
\]
Soit $V(t,S)$ solution régulière de
\[
V_t+(r-q)SV_S+\tfrac12\sigma^2S^2V_{SS}-rV=0,\qquad V(T,S)=\Phi(S).
\]
Alors Feynman--Kac donne
\[
V(t,S_t)=\E^{\Qmeas}\big[e^{-r(T-t)}\Phi(S_T)\mid \F_t\big].
\]
\emph{Idée} : appliquer Itô à $e^{-rt}V(t,S_t)$ ; la PDE force le terme en $\dd t$ à disparaître, donc $e^{-rt}V(t,S_t)$ est une martingale, et on identifie sa valeur par condition terminale.

\subsection*{Changement de variables : la PDE devient une équation de la chaleur}
En posant $x=\ln S$ et $\tau=T-t$, la PDE de Black--Scholes se transforme (après centrage) en une équation de la chaleur sur $\tau$.
Ce point explique pourquoi la solution fermée (formule BS) s'exprime en fonction de la loi normale.
Il est aussi utile numériquement : beaucoup de schémas PDE s'analysent plus facilement sur la version «heat equation».

% =========================================================
% Compléments M1 — Risk-neutral pricing (C. Dérivations)
% =========================================================

\subsection*{(1) Portefeuille auto-financé avec dividendes : dérivation propre de la PDE}
On détaille ici la dérivation «delta-hedging $\Rightarrow$ PDE» en prenant au sérieux le dividende continu $q$.

\paragraph{Hypothèses (cadre Black--Scholes).}
\begin{itemize}[leftmargin=*]
\item Un actif risqué (action) de prix ex-dividende $S_t$.
\item Dividende continu au taux $q$ : détenir une action rapporte un flux $qS_t\,\dd t$.
\item Un actif sans risque $B_t=e^{rt}$.
\item Sous $\Pmeas$ (mesure historique), on suppose $\dd S_t=\mu S_t\,\dd t+\sigma S_t\,\dd W_t$.
\end{itemize}

\paragraph{Étape 1 : dynamique de l'option par Itô.}
Soit $V(t,S)$ le prix d'une option européenne (payoff $\Phi(S_T)$).
Par Itô :
\[
\dd V
=\Big(V_t+\mu S V_S+\tfrac12\sigma^2S^2V_{SS}\Big)\dd t
 +\sigma S V_S\,\dd W_t.
\]

\paragraph{Étape 2 : construire un portefeuille couvert.}
Considérons le portefeuille
\[
\Pi_t = V(t,S_t)-\Delta_t S_t.
\]
Interprétation : long l'option, short $\Delta$ actions.
\emph{Attention dividendes :} en étant short $\Delta$ actions, on \emph{paie} le dividende au prêteur du titre.
Ainsi, le gain sur la position short est $-\Delta\,\dd S_t - \Delta\,qS_t\,\dd t$.
On obtient :
\[
\dd \Pi_t=\dd V - \Delta_t\,\dd S_t - \Delta_t q S_t\,\dd t.
\]

\paragraph{Étape 3 : annuler le risque brownien.}
Substituons $\dd V$ et $\dd S$ :
\[
\begin{aligned}
\dd \Pi_t
&=\Big(V_t+\mu S V_S+\tfrac12\sigma^2S^2V_{SS}\Big)\dd t +\sigma S V_S\,\dd W\\
&\quad-\Delta\big(\mu S\,\dd t+\sigma S\,\dd W\big)-\Delta qS\,\dd t.
\end{aligned}
\]
Choisissons
\[
\boxed{\Delta_t=V_S(t,S_t).}
\]
Alors les termes en $\dd W$ se compensent :
\[
\sigma S V_S\,\dd W - \Delta\sigma S\,\dd W = 0.
\]
Il reste un portefeuille \textbf{localement sans risque} :
\[
\dd \Pi_t=\Big(V_t+\tfrac12\sigma^2S^2V_{SS}-qSV_S\Big)\dd t.
\]

\paragraph{Étape 4 : absence d'arbitrage $\Rightarrow$ rendement au taux $r$.}
Un portefeuille sans risque doit rémunérer au taux sans risque :
\[
\dd \Pi_t = r\,\Pi_t\,\dd t = r\big(V-\Delta S\big)\dd t=r\big(V-SV_S\big)\dd t.
\]
En identifiant :
\[
V_t+\tfrac12\sigma^2S^2V_{SS}-qSV_S = rV-rSV_S.
\]
On regroupe :
\[
\boxed{
V_t+\tfrac12\sigma^2S^2V_{SS}+(r-q)SV_S-rV=0,
}
\qquad V(T,S)=\Phi(S).
\]
\paragraph{Message.}
Le drift $\mu$ a disparu : l'équation de pricing est indépendante de la prime de risque.

\subsection*{(2) De la PDE à l'espérance : preuve guidée de Feynman--Kac (cas BS)}
Supposons maintenant que $S$ suit sous $\Qmeas$ :
\[
\dd S_t=(r-q)S_t\,\dd t+\sigma S_t\,\dd W_t^{\Qmeas}.
\]
Soit $V$ solution régulière de la PDE BS ci-dessus.
Considérons le processus actualisé :
\[
M_t=e^{-rt}V(t,S_t).
\]
Par Itô (produit d'une fonction de $t,S_t$ par $e^{-rt}$) :
\[
\dd M_t = e^{-rt}\Big(\dd V-rV\,\dd t\Big).
\]
Or (sous $\Qmeas$) :
\[
\dd V=\Big(V_t+(r-q)SV_S+\tfrac12\sigma^2S^2V_{SS}\Big)\dd t+\sigma SV_S\,\dd W^{\Qmeas}.
\]
Donc
\[
\dd M_t=e^{-rt}\Big(\underbrace{V_t+(r-q)SV_S+\tfrac12\sigma^2S^2V_{SS}-rV}_{=0\ \text{par PDE}}\Big)\dd t
e^{-rt}\sigma SV_S\,\dd W^{\Qmeas}.
\]
Le terme de drift est nul. Il reste :
\[
\dd M_t=e^{-rt}\sigma S_tV_S(t,S_t)\,\dd W_t^{\Qmeas},
\]
donc $M$ est une martingale.
On en déduit, pour $t\le T$,
\[
M_t=\E^{\Qmeas}[M_T\mid \F_t]
=\E^{\Qmeas}\big[e^{-rT}V(T,S_T)\mid \F_t\big]
=\E^{\Qmeas}\big[e^{-rT}\Phi(S_T)\mid \F_t\big].
\]
En multipliant par $e^{rt}$ :
\[
\boxed{
V(t,S_t)=\E^{\Qmeas}\big[e^{-r(T-t)}\Phi(S_T)\mid \F_t\big].
}
\]
\paragraph{Point pédagogique.}
La PDE et la formule d'espérance sont \emph{la même information} écrite dans deux langages :
\begin{itemize}[leftmargin=*]
\item PDE : langage «local» (infinitésimal) ;
\item Feynman--Kac : langage «global» (espérance terminale).
\end{itemize}

\subsection*{(3) Girsanov en pratique : construire $\Qmeas$ et changer le drift}
On montre ici le mécanisme, dans le cas simple où $\lambda$ est constant.

\paragraph{Sous $\Pmeas$.}
\[
\frac{\dd S_t}{S_t}=\mu\,\dd t+\sigma\,\dd W_t^{\Pmeas}.
\]
On définit
\[
\lambda=\frac{\mu-(r-q)}{\sigma}.
\]

\paragraph{Étape 1 : densité exponentielle.}
On pose
\[
Z_t=\exp\Big(-\lambda W_t^{\Pmeas}-\tfrac12\lambda^2 t\Big).
\]
Du chapitre Brownien, on sait que $Z_t$ est une martingale (exponentielle de Brownien).
La condition de Novikov est vérifiée ici automatiquement (constante), donc on peut définir une mesure $\Qmeas$ par :
\[
\frac{\dd \Qmeas}{\dd \Pmeas}\Big|_{\F_T}=Z_T.
\]

\paragraph{Étape 2 : nouveau Brownien.}
Définissons
\[
W_t^{\Qmeas}=W_t^{\Pmeas}+\lambda t.
\]
Le théorème de Girsanov garantit que $(W_t^{\Qmeas})$ est un brownien sous $\Qmeas$.

\paragraph{Étape 3 : dynamique de $S$ sous $\Qmeas$.}
Comme $\dd W_t^{\Pmeas}=\dd W_t^{\Qmeas}-\lambda\,\dd t$, on substitue :
\[
\frac{\dd S_t}{S_t}
=\mu\,\dd t+\sigma(\dd W_t^{\Qmeas}-\lambda\,\dd t)
=(\mu-\sigma\lambda)\dd t+\sigma\,\dd W_t^{\Qmeas}
=(r-q)\dd t+\sigma\,\dd W_t^{\Qmeas}.
\]
On a bien remplacé le drift historique $\mu$ par le drift de pricing $(r-q)$.

\subsection*{(4) Exercices corrigés (pricing risque-neutre)}
\paragraph{Exercice 1 --- forward et drift risque-neutre.}
Sous $\Qmeas$, $\dd S_t=(r-q)S_t\,\dd t+\sigma S_t\,\dd W_t$.
Montrer que $\E^{\Qmeas}[S_T\mid \F_t]=S_t e^{(r-q)(T-t)}$.

\emph{Solution.}
On sait que $S_T=S_t\exp((r-q-\frac12\sigma^2)(T-t)+\sigma (W_T-W_t))$.
Conditionnellement à $\F_t$, $W_T-W_t\sim\Normal(0,T-t)$.
Donc
\[
\E^{\Qmeas}[S_T\mid \F_t]=S_t e^{(r-q-\frac12\sigma^2)(T-t)}\E[e^{\sigma (W_T-W_t)}]
=S_t e^{(r-q-\frac12\sigma^2)(T-t)}e^{\frac12\sigma^2(T-t)}
=S_t e^{(r-q)(T-t)}.
\]

\paragraph{Exercice 2 --- retrouver la PDE par un raisonnement «martingale».}
Partir de $M_t=e^{-rt}V(t,S_t)$ et imposer que $M$ soit une martingale sous $\Qmeas$.
Retrouver la PDE de Black--Scholes.

\emph{Solution.}
On calcule $\dd M_t$ par Itô. L'absence de drift impose que le coefficient de $\dd t$ soit nul, ce qui donne exactement
$V_t+(r-q)SV_S+\tfrac12\sigma^2S^2V_{SS}-rV=0$.

\paragraph{Exercice 3 --- digital call et dérivée en strike.}
Considérons le payoff digital $\Phi(S_T)=\1_{\{S_T>K\}}$.
\begin{enumerate}[leftmargin=*]
\item Montrer que son prix est $D_0=e^{-rT}\Qmeas(S_T>K)$.
\item Montrer que $-\partial_K C(K,T)=e^{-rT}\Qmeas(S_T>K)$ (formule de Breeden--Litzenberger au 1er ordre).
\end{enumerate}

\emph{Solution.}
Le prix risque-neutre donne immédiatement
\[
D_0=e^{-rT}\E^{\Qmeas}[\1_{\{S_T>K\}}]=e^{-rT}\Qmeas(S_T>K).
\]
Pour le call $C(K,T)=e^{-rT}\E^{\Qmeas}[(S_T-K)^+]$, on utilise que la dérivée de $(s-K)^+$ en $K$ vaut $-\1_{\{s>K\}}$ (pour $s\neq K$).
Sous des hypothèses standard permettant d'échanger dérivée et espérance :
\[
\partial_K C(K,T)=e^{-rT}\E^{\Qmeas}[\partial_K(S_T-K)^+]
=-e^{-rT}\E^{\Qmeas}[\1_{\{S_T>K\}}],
\]
d'où la relation demandée.

\paragraph{Exercice 4 --- digital en Black--Scholes : formule fermée.}
Sous BS (taux constants), montrer que
\[
\Qmeas(S_T>K)=N(d_2)
\quad\Rightarrow\quad
D_0=e^{-rT}N(d_2).
\]

\emph{Solution.}
Sous BS, $\ln S_T\sim \Normal(m,\sigma^2T)$ avec
$m=\ln S_0+(r-q-\tfrac12\sigma^2)T$.
Alors
\[
\Qmeas(S_T>K)=\Qmeas(\ln S_T>\ln K)
=\Qmeas\!\left(\frac{\ln S_T-m}{\sigma\sqrt{T}}>\frac{\ln K-m}{\sigma\sqrt{T}}\right).
\]
Le terme à gauche est une $\Normal(0,1)$, donc la probabilité vaut $1-N(\cdot)$.
Or
\[
d_2=\frac{\ln(S_0/K)+(r-q-\tfrac12\sigma^2)T}{\sigma\sqrt{T}}=\frac{m-\ln K}{\sigma\sqrt{T}},
\]
donc $\frac{\ln K-m}{\sigma\sqrt{T}}=-d_2$ et
\[
\Qmeas(S_T>K)=1-N(-d_2)=N(d_2).
\]
Le prix suit : $D_0=e^{-rT}N(d_2)$.

\paragraph{Exercice 5 --- Novikov (cas constant) et validité de Girsanov.}
On suppose $\lambda$ constant et on définit
\[
Z_T=\exp\Big(-\lambda W_T^{\Pmeas}-\tfrac12\lambda^2T\Big).
\]
Montrer que $\E^{\Pmeas}[Z_T]=1$ et que la condition de Novikov est satisfaite.

\emph{Solution.}
Comme $W_T^{\Pmeas}\sim\Normal(0,T)$, on sait que
\[
\E[e^{-\lambda W_T^{\Pmeas}}]=\exp\Big(\tfrac12\lambda^2T\Big).
\]
Donc
\[
\E[Z_T]=e^{-\frac12\lambda^2T}\E[e^{-\lambda W_T}]
=e^{-\frac12\lambda^2T}e^{\frac12\lambda^2T}=1.
\]
La condition de Novikov (pour assurer que $Z_t$ est une martingale) est
\[
\E\!\left[\exp\Big(\tfrac12\int_0^T \lambda^2\,\dd t\Big)\right]
=\exp\Big(\tfrac12\lambda^2T\Big)<\infty,
\]
qui est vraie. Girsanov s'applique donc sur $[0,T]$.

\subsection*{(6) Bonus : résoudre la PDE Black--Scholes via l'équation de la chaleur (détaillé)}
La formule BS peut se dériver en résolvant la PDE, ce qui complète la dérivation probabiliste.
On part de la PDE (call/put européen) :
\[
V_t+(r-q)SV_S+\tfrac12\sigma^2S^2V_{SS}-rV=0,\qquad V(T,S)=\Phi(S).
\]
On va la transformer en équation de la chaleur.

\paragraph{Étape 1 : renverser le temps.}
Posons $\tau=T-t$ (donc $\tau=0$ à maturité). Alors $V_t=-V_\tau$ et la PDE devient :
\[
V_\tau=(r-q)SV_S+\tfrac12\sigma^2S^2V_{SS}-rV.
\]

\paragraph{Étape 2 : passer en log-spot.}
Posons $x=\ln(S/K)$ (on centre sur le strike) et écrivons $V(\tau,S)=K\,u(\tau,x)$ avec $S=Ke^x$.
On calcule les dérivées (règle de chaîne) :
\[
V_S=\frac{\partial (K u)}{\partial S}=K u_x\frac{\partial x}{\partial S}=K u_x\frac{1}{S},
\qquad
V_{SS}=K\frac{\partial}{\partial S}\!\left(u_x\frac{1}{S}\right)
=K\left(u_{xx}\frac{1}{S^2}-u_x\frac{1}{S^2}\right).
\]
Donc
\[
SV_S=K u_x,\qquad S^2V_{SS}=K(u_{xx}-u_x).
\]
En injectant dans la PDE :
\[
K u_\tau=(r-q)K u_x+\tfrac12\sigma^2K(u_{xx}-u_x)-rKu,
\]
et en divisant par $K$ :
\[
u_\tau=\tfrac12\sigma^2u_{xx}+\Big(r-q-\tfrac12\sigma^2\Big)u_x-r u.
\]

\paragraph{Étape 3 : supprimer le terme $-ru$ (actualisation).}
Posons $u(\tau,x)=e^{-r\tau}\,w(\tau,x)$. Alors
$u_\tau=e^{-r\tau}(w_\tau-rw)$, donc
\[
e^{-r\tau}(w_\tau-rw)=\tfrac12\sigma^2 e^{-r\tau}w_{xx}+\Big(r-q-\tfrac12\sigma^2\Big)e^{-r\tau}w_x-r e^{-r\tau}w.
\]
Les termes $-rw$ s'annulent, et il reste :
\[
w_\tau=\tfrac12\sigma^2 w_{xx}+\Big(r-q-\tfrac12\sigma^2\Big)w_x.
\]

\paragraph{Étape 4 : supprimer le terme en $w_x$ (translation).}
Posons
\[
y=x+\Big(r-q-\tfrac12\sigma^2\Big)\tau,
\qquad
w(\tau,x)=\widetilde w(\tau,y).
\]
Alors $w_x=\widetilde w_y$ et $w_{xx}=\widetilde w_{yy}$, tandis que
$w_\tau=\widetilde w_\tau+\Big(r-q-\tfrac12\sigma^2\Big)\widetilde w_y$.
En injectant :
\[
\widetilde w_\tau+\Big(r-q-\tfrac12\sigma^2\Big)\widetilde w_y
=\tfrac12\sigma^2\widetilde w_{yy}+\Big(r-q-\tfrac12\sigma^2\Big)\widetilde w_y,
\]
donc finalement :
\[
\boxed{\widetilde w_\tau=\tfrac12\sigma^2\widetilde w_{yy}.}
\]
C'est l'équation de la chaleur.

\paragraph{Condition initiale (à $\tau=0$).}
Pour un call, $\Phi(S)=(S-K)^+$, donc en $x=\ln(S/K)$ :
\[
V(0,S)=\Phi(S)\quad\Rightarrow\quad u(0,x)=\frac{\Phi(Ke^x)}{K}=(e^x-1)^+.
\]
Comme $u(0,x)=w(0,x)=\widetilde w(0,y)$, on a :
\[
\widetilde w(0,y)=(e^y-1)^+.
\]

\paragraph{Conclusion (intuition, sans refaire toute l'intégrale).}
La solution de l'équation de la chaleur est une convolution par une gaussienne :
\[
\widetilde w(\tau,y)=\int_{-\infty}^{\infty} (e^z-1)^+\,
\frac{1}{\sqrt{2\pi\sigma^2\tau}}\exp\!\left(-\frac{(y-z)^2}{2\sigma^2\tau}\right)\dd z.
\]
En revenant aux variables $(t,S)$, cette intégrale se calcule explicitement et redonne
$C=S_0e^{-qT}N(d_1)-Ke^{-rT}N(d_2)$.
Le point important est conceptuel : la normalité intervient car l'équation de la chaleur a pour noyau fondamental une gaussienne.


\begin{comment}
\section{D. Numérique / algorithmes}
\subsection*{Crank--Nicolson : pourquoi ce choix ?}
Crank--Nicolson est un schéma implicite d'ordre 2 en temps (en pratique stable) utilisé pour des PDE paraboliques. Dans l'app, il est encapsulé par la classe \codeid{CrankNicolsonBS}.

\end{comment}
\begin{comment}
\section{E. Implémentation dans le repo}
\subsection*{Solveur PDE}
\begin{itemize}[leftmargin=*]
\item Classe \codeid{CrankNicolsonBS} : \filepath{app/model/options/engines/pricing.py}.
\item La couche service l'utilise directement : \filepath{app/model/options/service.py}, fonction \codeid{compute_price}.
\end{itemize}

\subsection*{Extrait de code (création du solveur depuis le service)}
Dans \filepath{app/model/options/service.py}, \codeid{compute_price} construit le solveur PDE à partir des champs option/marché :

\begin{lstlisting}[caption={Création du solveur PDE : \filepath{app/model/options/service.py}, \codeid{compute_price} (extrait).},label={lst:compute-price}]
cpflag = "p" if cpflag.startswith("p") else "c"
type_raw = str(opt.get("style") or opt.get("exercise") or "Eu").strip().lower()
if type_raw.startswith("am"):
typeflag = "Am"
elif type_raw.startswith("bm"):
typeflag = "Bmd"
else:
typeflag = "Eu"

S0 = float(opt.get("S0") or mkt.get("spot") or 0.0)
K = float(opt.get("strike") or 0.0)
T = float(opt.get("T") or 0.0)
vol = float(opt.get("sigma") or opt.get("iv") or 0.0)
r = float(opt.get("r") or mkt.get("r") or 0.0)
q = float(opt.get("q") or mkt.get("q") or 0.0)

solver = CrankNicolsonBS(typeflag, cpflag, S0=S0, K=K, T=T, vol=vol, r=r, d=q)
price, _, _, _ = solver.CN_option_info()
\end{lstlisting}

\paragraph{Pourquoi ce bout de code est la traduction directe de la PDE ?}
\begin{itemize}[leftmargin=*]
\item \codeid{cpflag} encode la condition terminale call/put $\Phi(S)=(S-K)^+$ ou $(K-S)^+$.
\item \codeid{typeflag} sélectionne les contraintes d'exercice (Eu/Am/Bmd), ce qui correspond à une PDE avec \emph{contrainte d'obstacle} (américaine/bermudane).
\item \codeid{CrankNicolsonBS(...)} construit la grille et résout numériquement la PDE (dans \codeid{CN_option_info}).
\end{itemize}

\section{F. Exemple concret (chiffré)}
\begin{example}[Prix BS (numérique)]
Avec $S_0=100$, $K=100$, $r=2\%$, $q=0$, $\sigma=20\%$, $T=1$ an, le prix call BS vaut environ $8.916$ (et le put $6.936$).
Ces valeurs sont cohérentes avec la formule implémentée dans \filepath{app/model/options/engines/black_scholes.py} (\codeid{black_scholes_price}).
\end{example}

\section{G. Limites \& extensions}
\begin{itemize}[leftmargin=*]
\item Les hypothèses BS sont souvent violées (smile, sauts, vol stochastique) : la Partie II explique comment l'app traite cela via des modèles de surface.
\item En pratique, hedging discret + coûts : la Partie III (RL) propose une approche «décisionnelle», mais le DQN présent est un prototype.
\end{itemize}

\end{comment}
\begin{jurybox}
La PDE de Black--Scholes vient d'une couverture : choisir $\Delta=V_S$ annule le risque brownien. En pratique, la résolution numérique de cette PDE (p. ex. schémas implicites/Crank--Nicolson) est un outil robuste pour pricer quand on ne dispose pas d'une formule fermée ou pour vérifier des moteurs plus complexes.
\end{jurybox}

% ---------------------------------------------------------
\chapter{Options vanilles : BS, greeks, implied vol (inversion numérique)}

\section{A. Problème financier}
Dans une salle de marché, on cote souvent une option via sa volatilité implicite $\sigma_{\text{impl}}$ plutôt que son prix. On veut :
\begin{itemize}[leftmargin=*]
\item pricer (BS), obtenir les grecs (sensibilités),
\item inverser BS : prix $\to$ $\sigma_{\text{impl}}$,
\item comparer une surface marché à une surface modèle.
\end{itemize}

\section{B. Théorie}
\subsection*{Formule BS (call/put) avec dividende continu}
Pour $d_1=\frac{\ln(S_0/K)+(r-q+\tfrac12\sigma^2)T}{\sigma\sqrt{T}}$, $d_2=d_1-\sigma\sqrt{T}$ :
\[
C = S_0e^{-qT}N(d_1)-Ke^{-rT}N(d_2),\qquad
P = Ke^{-rT}N(-d_2)-S_0e^{-qT}N(-d_1).
\]

\subsection*{Grecs (rappel)}
\[
\Delta_{\text{call}}=e^{-qT}N(d_1),\quad
\Gamma = \frac{e^{-qT}\varphi(d_1)}{S_0\sigma\sqrt{T}},\quad
\text{Vega}=S_0e^{-qT}\varphi(d_1)\sqrt{T},
\]
où $\varphi$ est la densité $\Normal(0,1)$.

\subsection*{Volatilité implicite}
Définition : $\sigma_{\text{impl}}$ est la solution de
\[
C_{\text{BS}}(S_0,K,r,q,T,\sigma)=C_{\text{mkt}}.
\]
Pour un call européen, $C_{\text{BS}}$ est croissant en $\sigma$ (Vega $>0$), donc l'inversion est un problème 1D bien posé (si le prix respecte les bornes d'arbitrage).

\qa
{Standardiser la comparaison marché/modèle ; construire une surface (K,T) d'IV ; alimenter la calibration (Partie II).}
{Option européenne, modèle BS pour définir l'IV ; prix dans les bornes (pas d'arbitrage statique).}
{IV dépend du modèle : sous sauts/vol stochastique, l'IV résume un modèle plus riche. Pour certains produits (américaines), l'IV BS est un proxy.}
{Brique de marché : cotation en IV et inversion numérique (root-finding robuste) ; permet de comparer des modèles hétérogènes via une même «échelle» $\sigma_{\text{impl}}(K,T)$.}

\subsection*{Interprétation de $d_1$ et $d_2$ (lecture probabiliste)}
Sous Black--Scholes (taux constants), on peut écrire
\[
C = e^{-rT}\E^{\Qmeas}\big[(S_T-K)^+\big].
\]
Comme $S_T$ est log-normal, les intégrales se ramènent à des probabilités gaussiennes et $d_1,d_2$ apparaissent comme des variables normalisées.
Une lecture classique est :
\begin{itemize}[leftmargin=*]
\item $N(d_2)$ est proche d'une \emph{probabilité risque-neutre} que l'option finisse dans la monnaie (ITM), ajustée par actualisation ;
\item $N(d_1)$ est lié à la probabilité «pondérée par $S_T$», ce qui explique pourquoi $\Delta_{\text{call}}=e^{-qT}N(d_1)$.
\end{itemize}
Ce sont des heuristiques (les mesures changent), mais elles donnent une intuition robuste.

\subsection*{Parité call--put et identités de base}
La parité call--put (européenne) sous dividende continu est :
\[
C - P = S_0e^{-qT} - Ke^{-rT}.
\]
Elle sert de contrôle qualité : si des prix de marché violent fortement cette égalité (en tenant compte du bid/ask), l'IV implicite devient instable ou incohérente.

% =========================================================
% Compléments M1 — Options vanilles (B. Théorie)
% =========================================================

\subsection*{(1) Trois parametrisations utiles : spot, forward et log-moneyness}
En pratique, travailler directement avec $(S_0,K,T)$ n'est pas toujours le plus stable.
On introduit souvent :
\begin{itemize}[leftmargin=*]
\item le \textbf{forward} $F_{0,T}=S_0e^{(r-q)T}$ ;
\item la \textbf{moneyness} $m=K/F_{0,T}$ ;
\item la \textbf{log-moneyness} $k=\ln(K/F_{0,T})$.
\end{itemize}
\paragraph{Pourquoi c'est utile ?}
À taux et dividendes constants, $F_{0,T}$ absorbe le drift risque-neutre.
La surface $\sigma_{\mathrm{impl}}(k,T)$ (ou la variance totale $w(k,T)=\sigma_{\mathrm{impl}}(k,T)^2T$) est souvent plus «stationnaire» visuellement que $\sigma(K,T)$.

\subsection*{(2) Convexité et densité risque-neutre : ce que la vanilla révèle sur la loi de $S_T$}
Un fait central est que le prix d'un call européen est une fonction convexe de $K$ :
\[
K\mapsto C(K,T) \ \text{est décroissante et convexe}.
\]
Intuition : $C(K,T)=e^{-rT}\E^{\Qmeas}[(S_T-K)^+]$ et $(x-K)^+$ est convexe en $K$.
\paragraph{Breeden--Litzenberger (message).}
Sous des hypothèses de régularité, la densité risque-neutre $f_{S_T}^{\Qmeas}$ est reliée à la dérivée seconde du call :
\[
\frac{\partial^2 C}{\partial K^2}(K,T)=e^{-rT}f_{S_T}^{\Qmeas}(K).
\]
Donc une surface de vanilles \emph{encode} une loi terminale.
Cela explique pourquoi l'arbitrage statique (monotonicité/convexité) est une contrainte fondamentale pour toute interpolation/extrapolation de surface.

\subsection*{(3) Vol implicite : ce qu'elle résume (et ce qu'elle ne résume pas)}
La volatilité implicite $\sigma_{\mathrm{impl}}$ est un \emph{paramètre de re-paramétrisation} : on exprime un prix via le BS «comme si» la vol était constante.
En dehors de BS :
\begin{itemize}[leftmargin=*]
\item une même $\sigma_{\mathrm{impl}}$ peut provenir de sauts, vol stochastique, mélange de lois, etc. ;
\item la surface $(K,T)\mapsto \sigma_{\mathrm{impl}}(K,T)$ résume des \emph{risques de distribution} (skew, smile, kurtosis) ;
\item l'inversion price $\to$ IV est un problème numérique (mais monotone tant que Vega$>0$).
\end{itemize}
En finance quantitative, on utilise l'IV comme interface standardisée : on compare modèles/produits sur une échelle commune, tout en gardant à l'esprit qu'il s'agit d'un «résumé» (pas un paramètre structural).


\section{C. Dérivations (ultra détaillées)}
\subsection*{Bornes d'arbitrage et monotonicité}
Pour un call :
\[
\max(S_0e^{-qT}-Ke^{-rT},0)\le C \le S_0e^{-qT}.
\]
Si $C$ est hors bornes, \emph{aucune} volatilité ne peut l'expliquer dans BS $\Rightarrow$ on renvoie NaN.

\subsection*{Inversion numérique : méthode de Brent}
On résout $f(\sigma)=C_{\text{BS}}(\sigma)-C_{\text{mkt}}=0$.
Brent combine bissection (robuste) et interpolation (rapide) et requiert un encadrement $[\sigma_{\min},\sigma_{\max}]$ tel que $f(\sigma_{\min})f(\sigma_{\max})\le 0$.

\subsection*{Pourquoi l'inversion est bien posée : Vega $>0$}
On a (call européen) :
\[
\frac{\partial C_{\text{BS}}}{\partial \sigma}= \text{Vega} = S_0e^{-qT}\varphi(d_1)\sqrt{T} > 0 \quad (T>0).
\]
Donc $C_{\text{BS}}(\sigma)$ est strictement croissant en $\sigma$ et l'équation $C_{\text{BS}}(\sigma)=C_{\text{mkt}}$ a au plus une solution.
Cette propriété explique aussi un fait pratique : quand $T$ est très petit ou quand l'option est très loin de la monnaie, $\varphi(d_1)$ devient minuscule et l'inversion devient numériquement fragile (Vega quasi nulle).

\subsection*{Newton : dérivation de l'itération (et pourquoi on préfère Brent en production)}
Newton s'obtient en linéarisant $f(\sigma)$ :
\[
0=f(\sigma^\star)\approx f(\sigma_n)+f'(\sigma_n)(\sigma^\star-\sigma_n),
\]
donc
\[
\sigma_{n+1}=\sigma_n - \frac{f(\sigma_n)}{f'(\sigma_n)}=\sigma_n - \frac{C_{\text{BS}}(\sigma_n)-C_{\text{mkt}}}{\text{Vega}(\sigma_n)}.
\]
Newton est très rapide près de la solution, mais peut diverger si l'initialisation est mauvaise ou si Vega est faible ; Brent sécurise l'encadrement et évite ces pathologies.

\subsection*{Dérivation de la formule de Black--Scholes (call) : pas à pas}
Sous $\Qmeas$ (taux constants), on a
\[
\frac{\dd S_t}{S_t}=(r-q)\,\dd t+\sigma\,\dd W_t^{\Qmeas}
\quad\Rightarrow\quad
S_T=S_0\exp\Big((r-q-\tfrac12\sigma^2)T+\sigma\sqrt{T}\,Z\Big),
\]
où $Z\sim\Normal(0,1)$.
Le prix du call est
\[
C=e^{-rT}\E^{\Qmeas}[(S_T-K)^+]
  =e^{-rT}\Big(\E[S_T\1_{\{S_T>K\}}]-K\,\Qmeas(S_T>K)\Big).
\]
Écrivons l'événement $S_T>K$ en termes de $Z$ :
\[
\ln S_T > \ln K
\iff
Z > \frac{\ln(K/S_0)-(r-q-\tfrac12\sigma^2)T}{\sigma\sqrt{T}}
      =-d_2,
\]
avec
\[
d_2=\frac{\ln(S_0/K)+(r-q-\tfrac12\sigma^2)T}{\sigma\sqrt{T}},\qquad d_1=d_2+\sigma\sqrt{T}.
\]
Donc $\Qmeas(S_T>K)=\Prob(Z>-d_2)=N(d_2)$.
Pour le terme $\E[S_T\1_{\{S_T>K\}}]$, on utilise la forme exponentielle de $S_T$ :
\[
\E[S_T\1_{\{Z>-d_2\}}]=S_0e^{(r-q-\tfrac12\sigma^2)T}\E\big[e^{\sigma\sqrt{T}Z}\1_{\{Z>-d_2\}}\big].
\]
Une identité classique sur la normale (complétion du carré) donne :
\[
\E\big[e^{aZ}\1_{\{Z>b\}}\big]=e^{\tfrac12 a^2}\,N(a-b)\qquad (Z\sim\Normal(0,1)),
\]
ici avec $a=\sigma\sqrt{T}$ et $b=-d_2$, d'où
\[
\E[S_T\1_{\{Z>-d_2\}}]=S_0e^{(r-q)T}N(d_1).
\]
En assemblant :
\[
\boxed{C=S_0e^{-qT}N(d_1)-Ke^{-rT}N(d_2).}
\]
Le put s'obtient soit par un calcul analogue, soit via la parité call--put.

\subsection*{Dérivation (rapide) de quelques grecs}
On dérive $C(S_0,\sigma,\ldots)$ :
\begin{itemize}[leftmargin=*]
\item \textbf{Delta.} On obtient $\Delta=\partial_{S_0}C=e^{-qT}N(d_1)$ (le terme en $\partial d_1/\partial S_0$ s'annule grâce à une identité entre $N$ et $\varphi$).
\item \textbf{Gamma.} Une dérivation supplémentaire donne
\[
\Gamma=\partial_{S_0S_0}C=\frac{e^{-qT}\varphi(d_1)}{S_0\sigma\sqrt{T}}.
\]
\item \textbf{Vega.} La dérivée par rapport à $\sigma$ donne
\[
\text{Vega}=\partial_\sigma C=S_0e^{-qT}\varphi(d_1)\sqrt{T}.
\]
\end{itemize}
Ces expressions expliquent des faits pratiques : Gamma explose quand $T\to 0$ à l'ATM, et Vega s'effondre loin de la monnaie.

% =========================================================
% Compléments M1 — Options vanilles (C. Dérivations)
% =========================================================

\subsection*{(1) Identité clé : $S_0e^{-qT}\varphi(d_1)=Ke^{-rT}\varphi(d_2)$}
Cette identité est extrêmement utile pour dériver les grecs proprement (elle permet de simplifier des termes).
Rappel :
\[
d_1=\frac{\ln(S_0/K)+(r-q+\tfrac12\sigma^2)T}{\sigma\sqrt{T}},
\qquad
d_2=d_1-\sigma\sqrt{T}.
\]
On calcule :
\[
\frac{\varphi(d_1)}{\varphi(d_2)}
=\exp\Big(-\tfrac12 d_1^2+\tfrac12 d_2^2\Big)
=\exp\Big(\tfrac12(d_2^2-d_1^2)\Big)
=\exp\Big(\tfrac12(d_2-d_1)(d_2+d_1)\Big).
\]
Or $d_2-d_1=-\sigma\sqrt{T}$ et $d_2+d_1=2d_1-\sigma\sqrt{T}$, donc
\[
\tfrac12(d_2^2-d_1^2)=\tfrac12(-\sigma\sqrt{T})(2d_1-\sigma\sqrt{T})
=-\sigma\sqrt{T}\,d_1+\tfrac12\sigma^2T.
\]
En remplaçant $d_1$ :
\[
-\sigma\sqrt{T}\,d_1+\tfrac12\sigma^2T
=-\Big(\ln(S_0/K)+(r-q+\tfrac12\sigma^2)T\Big)+\tfrac12\sigma^2T
=-\ln(S_0/K)-(r-q)T.
\]
Donc
\[
\frac{\varphi(d_1)}{\varphi(d_2)}=\exp\big(-\ln(S_0/K)-(r-q)T\big)=\frac{K}{S_0}e^{-(r-q)T},
\]
ce qui équivaut à
\[
\boxed{S_0e^{-qT}\varphi(d_1)=Ke^{-rT}\varphi(d_2).}
\]

\subsection*{(2) Dérivation détaillée de $\Delta$ et Vega}
On part du call :
\[
C(S_0,\sigma)=S_0e^{-qT}N(d_1)-Ke^{-rT}N(d_2).
\]
\paragraph{Delta.}
On dérive en $S_0$ :
\[
\frac{\partial C}{\partial S_0}
=e^{-qT}N(d_1)+S_0e^{-qT}\varphi(d_1)\frac{\partial d_1}{\partial S_0}
-Ke^{-rT}\varphi(d_2)\frac{\partial d_2}{\partial S_0}.
\]
Or $\partial d_1/\partial S_0 = \frac{1}{S_0\sigma\sqrt{T}}$ et $\partial d_2/\partial S_0=\partial d_1/\partial S_0$.
Donc
\[
\frac{\partial C}{\partial S_0}
=e^{-qT}N(d_1)+e^{-qT}\frac{\varphi(d_1)}{\sigma\sqrt{T}}
-Ke^{-rT}\varphi(d_2)\frac{1}{S_0\sigma\sqrt{T}}.
\]
En utilisant l'identité $S_0e^{-qT}\varphi(d_1)=Ke^{-rT}\varphi(d_2)$, le dernier terme vaut exactement
$e^{-qT}\frac{\varphi(d_1)}{\sigma\sqrt{T}}$ et s'annule avec le second.
Ainsi :
\[
\boxed{\Delta_{\mathrm{call}}=\frac{\partial C}{\partial S_0}=e^{-qT}N(d_1).}
\]

\paragraph{Vega.}
On dérive en $\sigma$ :
\[
\frac{\partial C}{\partial \sigma}
=S_0e^{-qT}\varphi(d_1)\frac{\partial d_1}{\partial \sigma}
-Ke^{-rT}\varphi(d_2)\frac{\partial d_2}{\partial \sigma}.
\]
Comme $d_2=d_1-\sigma\sqrt{T}$, on a $\partial d_2/\partial\sigma=\partial d_1/\partial\sigma-\sqrt{T}$.
Donc
\[
\frac{\partial C}{\partial \sigma}
=\Big(S_0e^{-qT}\varphi(d_1)-Ke^{-rT}\varphi(d_2)\Big)\frac{\partial d_1}{\partial \sigma}
Ke^{-rT}\varphi(d_2)\sqrt{T}.
\]
Le terme entre parenthèses est nul par l'identité ci-dessus, donc
\[
\boxed{\mathrm{Vega}=\frac{\partial C}{\partial \sigma}=Ke^{-rT}\varphi(d_2)\sqrt{T}=S_0e^{-qT}\varphi(d_1)\sqrt{T}.}
\]
\paragraph{Message pratique.}
Vega $>0$ garantit que l'inversion prix $\to$ vol est bien posée tant qu'on reste dans les bornes d'arbitrage.

\subsection*{(3) Theta (sensibilité au temps) : formule et lecture}
Theta mesure l'effet du passage du temps (décroissance de valeur temps).
On peut montrer (call) :
\[
\Theta_{\mathrm{call}}
=-\frac{S_0e^{-qT}\varphi(d_1)\sigma}{2\sqrt{T}} + qS_0e^{-qT}N(d_1)-rKe^{-rT}N(d_2).
\]
\paragraph{Lecture.}
Le premier terme est la «theta diffusion» (lié à la convexité/Gamma), les autres proviennent du carry (dividendes) et de l'actualisation.

\subsection*{(4) Inversion IV : Newton sécurisé (pour comprendre Brent)}
Résoudre $f(\sigma)=C_{\mathrm{BS}}(\sigma)-C_{\mathrm{mkt}}=0$.
Newton donne :
\[
\sigma_{n+1}=\sigma_n-\frac{f(\sigma_n)}{f'(\sigma_n)}
=\sigma_n-\frac{C_{\mathrm{BS}}(\sigma_n)-C_{\mathrm{mkt}}}{\mathrm{Vega}(\sigma_n)}.
\]
\paragraph{Pourquoi Newton peut échouer ?}
Si Vega est très petit (options très OTM ou maturité très courte), le pas peut exploser.
On sécurise alors par :
\begin{itemize}[leftmargin=*]
\item \textbf{bracketing} (bornes $[\sigma_{\min},\sigma_{\max}]$) ;
\item \textbf{clipping} : limiter l'amplitude du pas ;
\item \textbf{fallback} : bisection/Brent si Newton sort des bornes.
\end{itemize}
Brent peut être vu comme un «Newton sécurisé» qui ne perd jamais la racine si le bracket est valide.

\subsection*{(5) Point subtil : initial guess et stabilité numérique}
Un bon point de départ accélère fortement l'inversion.
Pour un call proche de l'ATM-forward, une approximation simple (Brenner--Subrahmanyam) est :
\[
\sigma_{\mathrm{init}}\approx \sqrt{\frac{2\pi}{T}}\frac{C}{S_0e^{-qT}}.
\]
Ce n'est pas une formule universelle, mais elle donne une intuition : plus $C$ est grand (à moneyness donnée), plus la vol est grande.

\subsection*{(6) Exercices corrigés (BS / grecs / IV)}
\paragraph{Exercice 1 --- contrôle call--put parity.}
Avec $S_0=100$, $K=110$, $T=1$, $r=2\%$, $q=0$, on observe $C=4.2$ et $P=12.5$.
Vérifier la parité call--put et conclure.

\emph{Solution.}
La parité impose $C-P=S_0-Ke^{-rT}=100-110e^{-0.02}\approx 100-107.82=-7.82$.
Or $C-P=4.2-12.5=-8.3$, écart $\approx -0.48$.
Selon le bid/ask, c'est soit une incohérence (arbitrage théorique), soit un écart compatible avec spreads/erreurs de cotation.

\paragraph{Exercice 2 --- Vega et difficulté d'inversion.}
Expliquer pourquoi l'inversion d'IV est instable pour une option très OTM à courte maturité.

\emph{Solution.}
Pour une option très OTM, $|d_1|$ est grand, donc $\varphi(d_1)$ est exponentiellement petit.
Comme $\mathrm{Vega}=S_0e^{-qT}\varphi(d_1)\sqrt{T}$, Vega est proche de 0 : un petit bruit sur le prix correspond à une grande variation de $\sigma$, ce qui rend la racine numériquement instable.


\begin{comment}
\section{D. Numérique / algorithmes}
\begin{itemize}[leftmargin=*]
\item Robustesse : vérifier les bornes d'arbitrage avant de lancer le solveur.
\item Choix d'intervalle : typiquement $\sigma\in[10^{-4},5]$.
\item Complexité : $O(N_{\text{grid}}\times N_{\text{iter}})$ pour une grille d'IV.
\end{itemize}

\end{comment}
\begin{comment}
\section{E. Implémentation dans le repo}
\subsection*{Inversion d'IV (extrait)}
Dans \filepath{app/model/calibration/implied_vol.py} :

\begin{lstlisting}[caption={Inversion IV par Brent : \filepath{app/model/calibration/implied_vol.py}, \codeid{implied_vol_call} (extrait).},label={lst:iv-brent}]
def implied_vol_call(price, S0, K, t, r, q, vol_min=1e-4, vol_max=5.0):
if t <= 0 or S0 <= 0 or K <= 0:
return np.nan
intrinsic = max(S0 * math.exp(-q * t) - K * math.exp(-r * t), 0.0)
upper_bound = S0 * math.exp(-q * t)
if price < intrinsic - 1e-8 or price > upper_bound + 1e-8:
return np.nan

def f(vol):
    return bs_call_price(S0, K, t, r, q, vol) - price

try:
    return float(brentq(f, vol_min, vol_max, maxiter=100, xtol=1e-6))
except Exception:
    return np.nan
\end{lstlisting}

\paragraph{Traduction directe de la théorie.}
\begin{itemize}[leftmargin=*]
\item Les bornes \codeid{intrinsic} et \codeid{upper_bound} implémentent l'absence d'arbitrage statique.
\item \codeid{f(vol)} est $C_{\text{BS}}(\sigma)-C_{\text{mkt}}$.
\item \codeid{brentq} réalise l'inversion numérique robuste.
\end{itemize}

\subsection*{Grecs (proxy vanilla)}
Les grecs sont calculés via BS dans \filepath{app/model/options/core/greeks.py} (\codeid{compute_vanilla_greeks}). Pour des structures multi-jambes, \codeid{compute_option_greeks} prend la première jambe comme proxy (\NonImplem\ au sens «grecs exacts multi-jambes»).

\section{F. Exemple end-to-end}
\begin{example}[Surface modèle Heston $\to$ IV]
Le pricer Heston «V1» (\filepath{app/model/calibration/heston_pricer.py}, \codeid{call_price_cf}) produit des prix, puis \codeid{implied_vol_call} reconstruit l'IV. Par exemple, avec $(\kappa,\theta,\sigma,\rho,v_0)=(2,0.04,0.5,-0.7,0.04)$, $S_0=K=100$, $r=2\%$, $q=0$, $T=1$, le code donne un prix $\approx 8.664$ et une IV implicite $\approx 0.194$ (mécanisme : prix $\to$ Brent $\to$ IV).
\end{example}

\section{G. Limites \& extensions}
\begin{itemize}[leftmargin=*]
\item L'IV BS est un «résumé» : deux modèles différents peuvent produire la même IV sur un point mais diverger hors grille.
\item Pour les options américaines, l'IV BS n'est pas naturellement définie (il faut un pricer américain en forward). L'app utilise des pricers CRR/LSMC pour ces styles.
\end{itemize}

\end{comment}
\begin{jurybox}
La volatilité implicite est l'interface naturelle entre marché et modèle : elle transforme une grille hétérogène de prix en une représentation comparable $\sigma_{\text{impl}}(K,T)$. L'inversion Black--Scholes (root-finding) est la brique standard qui relie prix $\leftrightarrow$ IV et sert de passerelle vers la calibration.
\end{jurybox}

% =========================================================
% PARTIE II — Volatility surface \& calibration
% =========================================================
\part{Volatility surface \& calibration (c\oe{}ur de la modélisation IV)}

\chapter{Surface d'IV : smile/skew, axes, interpolation, arbitrage}

\section{A. Problème financier}
Le marché cote une \emph{surface} $\sigma_{\text{impl}}(K,T)$ (ou en moneyness $m=K/S_0$). Objectif :
\begin{itemize}[leftmargin=*]
\item représenter proprement les observations (données bruitées et incomplètes),
\item interpoler pour pricer des strikes non cotés,
\item calibrer un modèle paramétrique pour reproduire la surface.
\end{itemize}

\section{B. Théorie}
\subsection*{Définition formelle}
On appelle \emph{surface d'IV} la fonction
\[
(m,T)\mapsto \sigma_{\text{impl}}(m,T),
\]
où $m=K/S_0$ et $T$ est la maturité en années. En pratique, on standardise souvent ces deux axes pour obtenir une grille stable.

\subsection*{Intuition financière (smile/skew)}
\begin{itemize}[leftmargin=*]
\item \textbf{Skew} (actions) : les puts OTM sont chers $\Rightarrow$ IV plus élevée pour $K<S_0$ (demande de protection, levier, crash risk).
\item \textbf{Smile} (FX) : symétrie plus marquée due à la structure de risque différente.
\end{itemize}

\qa
{Construire un objet «marché» qui alimente calibration et pricing ; éviter de pricer sur des trous de données.}
{Suppose qu'on peut représenter le marché par des IV pointwise (ce qui est une convention).}
{Interpolation peut violer le non-arbitrage (convexité en strike, monotonicité en maturité) : une surface d'IV «jolie» n'est pas forcément économiquement admissible.}
{Pratique : projection sur une grille + masque d'observation, puis interpolation contrôlée (ou nearest) avant toute calibration ; l'objet «surface» devient une entrée standardisée pour les modèles.}

\subsection*{Choix des axes : strike, moneyness, log-moneyness, delta}
Une «surface» dépend fortement du système de coordonnées :
\begin{itemize}[leftmargin=*]
\item \textbf{Strike $K$} : naturel contractuellement, mais non stationnaire quand $S_0$ bouge.
\item \textbf{Moneyness $m=K/S_0$} : rend l'ATM proche de 1, stabilise les grilles.
\item \textbf{Log-moneyness $k=\ln(K/F_{0,T})$} : souvent préférable car symétrise les formules et relie directement aux développements asymptotiques.
\item \textbf{Delta} : utilisé en FX (surface cotée en delta) ; la conversion dépend du modèle (BS) et peut être circulaire.
\end{itemize}
Le point clé : \emph{l'interpolation est plus «douce» dans certains axes}, et les contraintes d'absence d'arbitrage s'expriment plus naturellement sur d'autres (souvent sur les prix ou la variance totale).

\subsection*{Arbitrage statique : contraintes minimales à respecter}
Pour des calls européens $C(K,T)$, l'absence d'arbitrage statique implique au minimum :
\begin{itemize}[leftmargin=*]
\item $C$ décroît en $K$ (plus le strike est haut, moins le call vaut) ;
\item $C$ est convexe en $K$ (densité risque-neutre non négative) ;
\item $C$ croît en $T$ (calendar spread : plus de temps ne doit pas rendre l'option moins chère, toutes choses égales par ailleurs).
\end{itemize}
Une interpolation naïve directement sur $\sigma_{\text{impl}}$ peut violer ces contraintes même si les points de départ (prix de marché) les respectent.

\subsection*{Variance totale $w(k,T)=\sigma_{\text{impl}}(k,T)^2T$}
La variance totale est souvent plus régulière que $\sigma_{\text{impl}}$ et est au c\oe{}ur de nombreuses paramétrisations (ex. SVI).
Elle permet aussi de discuter l'arbitrage «butterfly» et la croissance en maturité de manière plus structurée.

% =========================================================
% Compléments M1 — Surface d'IV (B. Théorie)
% =========================================================

\subsection*{(1) Trois niveaux de représentation : prix, IV, variance totale}
Une surface peut être représentée sous trois formes principales :
\begin{itemize}[leftmargin=*]
\item \textbf{Prix} $C(K,T)$ (ou $P(K,T)$) : c'est l'objet «économique» directement arbitrageable.
\item \textbf{Vol implicite} $\sigma_{\mathrm{impl}}(K,T)$ : convention de marché (échelle comparable).
\item \textbf{Variance totale} $w(k,T)=\sigma_{\mathrm{impl}}(k,T)^2T$ avec $k=\ln(K/F_{0,T})$ : souvent plus régulière.
\end{itemize}
\paragraph{Message.}
L'arbitrage statique s'exprime naturellement en \emph{prix}, mais les modèles paramétriques et les interpolations sont souvent plus «lisses» en variance totale.
Un bon pipeline de surface sait passer de l'une à l'autre sans perdre la cohérence.

\subsection*{(2) Arbitrage statique en strike : monotonicité, convexité, densité}
Pour un call européen $C(K,T)$ (taux constants) :
\begin{itemize}[leftmargin=*]
\item \textbf{Bornes} : $0\le C(K,T)\le S_0e^{-qT}$ et $C(K,T)\ge \max(S_0e^{-qT}-Ke^{-rT},0)$.
\item \textbf{Monotonicité} : $K_1<K_2 \Rightarrow C(K_1,T)\ge C(K_2,T)$.
\item \textbf{Convexité} : $C(\cdot,T)$ est convexe en $K$.
\end{itemize}
\paragraph{Interprétation.}
La convexité est l'expression «visible» du fait qu'une densité risque-neutre ne peut pas être négative.
Elle garantit qu'une stratégie de type butterfly ne crée pas un payoff positif pour un coût négatif.

\subsection*{(3) Arbitrage calendaire : cohérence en maturité}
Fixons $K$.
Sans arbitrage calendaire, un call de maturité plus longue ne doit pas être moins cher :
\[
T_1<T_2 \Rightarrow C(K,T_1)\le C(K,T_2).
\]
\paragraph{Pourquoi c'est non trivial ?}
Parce qu'une interpolation séparée en $T$ et en $K$ peut facilement violer cette propriété : on peut obtenir des surfaces «douces» mais économiquement impossibles.
\paragraph{Lien pratique avec $w(k,T)$.}
Dans beaucoup de paramétrisations, imposer que $T\mapsto w(k,T)$ soit croissante pour tout $k$ est une contrainte suffisante (pas toujours nécessaire) pour éviter le calendar arbitrage.
Cela motive la représentation en variance totale.

\subsection*{(4) Butterfly arbitrage : de la convexité aux contraintes sur la variance totale}
Le butterfly arbitrage se lit sur la convexité en $K$.
Mais si l'on travaille en $w(k,T)$, la condition devient une contrainte \emph{non linéaire} sur $w$ et ses dérivées en $k$.
Message M1 :
\begin{itemize}[leftmargin=*]
\item travailler en IV rend l'interpolation facile mais risque d'arbitrage ;
\item travailler en prix garantit l'interprétation économique mais est moins «lisse» ;
\item travailler en $w$ est un compromis utile, surtout avec des familles comme SVI.
\end{itemize}

\subsection*{(5) SVI (aperçu) : une paramétrisation standard de $w(k)$}
En equity, une paramétrisation très utilisée pour une maturité fixée est \textbf{SVI} :
\[
w(k)=a+b\Big(\rho(k-m)+\sqrt{(k-m)^2+\eta^2}\Big),
\]
avec paramètres $(a,b,\rho,m,\eta)$ et contraintes $b>0$, $|\rho|<1$, $\eta>0$.
\paragraph{Lecture qualitative.}
\begin{itemize}[leftmargin=*]
\item $a$ : niveau de variance (ATM) ;
\item $b$ : pente globale (amplitude du smile) ;
\item $\rho$ : asymétrie (skew) ;
\item $m$ : déplacement du centre ;
\item $\eta$ : courbure (arrondi du smile).
\end{itemize}
SVI est populaire car il est flexible, relativement stable en calibration, et on sait formuler des contraintes d'absence de butterfly arbitrage (au moins sous forme suffisante).

\subsection*{(6) Delta-surface (FX) : attention au caractère «modèle-dépendant»}
En FX, on cote souvent la surface en delta plutôt qu'en strike.
Mais \emph{la conversion delta $\leftrightarrow$ strike dépend du modèle} (souvent BS).
Cela crée un point subtil :
\begin{itemize}[leftmargin=*]
\item la surface en delta est pratique pour l'usage trader ;
\item la surface en strike est plus naturelle pour l'arbitrage statique (convexité en $K$) ;
\item passer de l'une à l'autre demande de choisir une convention (BS forward delta, premium-adjusted delta, etc.).
\end{itemize}
En M1, il faut surtout retenir que le choix des axes n'est pas innocent : il change la «géométrie» du problème d'interpolation et de calibration.


\section{C. Dérivations (ultra détaillées)}
\subsection*{Pourquoi passer en moneyness $m=K/S_0$ ?}
Si le sous-jacent change d'échelle (stock split, variation de prix), une grille en strike absolu devient instable. En moneyness, l'ATM est proche de $m=1$ et on stabilise l'espace de calibration.

\subsection*{Surface discrète et masque d'observation}
On observe des points $(m_i,T_i,\sigma_i)$ irréguliers. Pour calibrer, on construit une grille $(m_j){j=1..M}$, $(T_k){k=1..K}$ et on forme une matrice $\Sigma_{k,j}$.
Comme tous les points ne sont pas observés, on utilise un \emph{masque} $M_{k,j}\in{0,1}$ qui indique quels points proviennent d'observations. Cela permet :
\begin{itemize}[leftmargin=*]
\item de ne pas pénaliser un modèle sur des valeurs interpolées : on compare la perte uniquement sur les points réellement observés,
\item de distinguer «data» vs «remplissage».
\end{itemize}

\subsection*{Pourquoi l'arbitrage est une question de convexité (Breeden--Litzenberger)}
On part de
\[
C(K,T)=e^{-rT}\E^{\Qmeas}[(S_T-K)^+].
\]
En dérivant par rapport à $K$ (formellement, et sous conditions techniques standard), on obtient :
\[
\frac{\partial C}{\partial K}(K,T)=-e^{-rT}\Qmeas(S_T>K),
\]
et une seconde dérivée donne la densité risque-neutre :
\[
\frac{\partial^2 C}{\partial K^2}(K,T)=e^{-rT} f_{S_T}(K)\ge 0.
\]
Donc \textbf{convexité en $K$} $\Longleftrightarrow$ \textbf{densité non négative}.
Cette relation est la pierre angulaire des tests d'arbitrage sur une surface.

\subsection*{Lien (optionnel) : surface d'IV $\to$ volatilité locale (Dupire)}
Si l'on dispose de $C(K,T)$ suffisamment régulier, la formule de Dupire exprime une volatilité locale $\sigma_{\text{loc}}(K,T)$ telle que le modèle de diffusion locale reproduise exactement les prix vanilles :
\[
\sigma_{\text{loc}}^2(K,T)=\frac{\partial_T C(K,T)}{\tfrac12 K^2 \partial_{KK} C(K,T)}\quad \text{(forme simplifiée, taux constants)}.
\]
Cela illustre un point important : une surface d'IV contient \emph{beaucoup} d'information sur la distribution risque-neutre, mais toute interpolation doit préserver la régularité requise.

% =========================================================
% Compléments M1 — Surface d'IV (C. Dérivations)
% =========================================================

\subsection*{(1) Breeden--Litzenberger détaillé : du call à la densité}
On part du prix d'un call européen (taux constants pour simplifier) :
\[
C(K,T)=e^{-rT}\E^{\Qmeas}[(S_T-K)^+].
\]
Écrivons cette espérance sous forme intégrale, avec densité $f_{S_T}^{\Qmeas}$ :
\[
C(K,T)=e^{-rT}\int_0^{\infty}(s-K)^+\,f_{S_T}^{\Qmeas}(s)\,\dd s
=e^{-rT}\int_K^{\infty}(s-K)\,f_{S_T}^{\Qmeas}(s)\,\dd s.
\]
\paragraph{Première dérivée en $K$.}
On dérive (formellement, sous conditions de régularité permettant l'échange dérivée/intégrale) :
\[
\frac{\partial C}{\partial K}(K,T)
=e^{-rT}\frac{\partial}{\partial K}\int_K^{\infty}(s-K)\,f(s)\,\dd s,
\quad f=f_{S_T}^{\Qmeas}.
\]
Utilisons la règle de Leibniz :
\[
\frac{\partial}{\partial K}\int_K^{\infty}g(K,s)\,\dd s
=-g(K,K)+\int_K^{\infty}\partial_K g(K,s)\,\dd s,
\quad g(K,s)=(s-K)f(s).
\]
Ici, $g(K,K)=0$ et $\partial_K g(K,s)=-f(s)$. Donc
\[
\frac{\partial C}{\partial K}(K,T)
=e^{-rT}\int_K^{\infty}(-f(s))\,\dd s
=-e^{-rT}\Qmeas(S_T\ge K).
\]
\paragraph{Seconde dérivée.}
On dérive à nouveau :
\[
\frac{\partial^2 C}{\partial K^2}(K,T)
=-e^{-rT}\frac{\partial}{\partial K}\int_K^{\infty} f(s)\,\dd s
=e^{-rT}f(K).
\]
On obtient l'identité :
\[
\boxed{\frac{\partial^2 C}{\partial K^2}(K,T)=e^{-rT}f_{S_T}^{\Qmeas}(K)\ge 0.}
\]
\paragraph{Message.}
La convexité en $K$ équivaut à la positivité de la densité risque-neutre : c'est la base des tests d'arbitrage sur une surface.

\subsection*{(2) Dupire (forme simplifiée) : de $C(K,T)$ à une volatilité locale}
L'idée de Dupire est : \emph{si l'on connaît $C(K,T)$ pour tous $K,T$, on peut fabriquer un modèle de diffusion locale qui reproduit exactement ces prix vanilles}.

Pour garder l'algèbre lisible, on se place dans un cadre simplifié :
\[
r=0,\qquad q=0.
\]
Alors le prix d'un call est simplement
\[
C(K,T)=\E^{\Qmeas}[(S_T-K)^+].
\]
Sous un modèle de volatilité locale
\[
\dd S_t=\sigma_{\mathrm{loc}}(S_t,t)\,S_t\,\dd W_t^{\Qmeas},
\]
la densité $f(s,T)$ de $S_T$ satisfait une équation de Fokker--Planck (équation «forward»).
Dans ce cadre, Dupire montre que $C(K,T)$ satisfait la PDE en variables $(K,T)$ :
\[
\partial_T C(K,T)=\frac12\sigma_{\mathrm{loc}}^2(K,T)\,K^2\,\partial_{KK}C(K,T).
\]
\paragraph{Interprétation.}
Comme $\partial_{KK}C=e^{-rT}f$ (ici $r=0$), la dérivée seconde joue le rôle de densité, et $\partial_T C$ encode comment cette densité «diffuse» dans le temps.

En isolant $\sigma_{\mathrm{loc}}$ :
\[
\boxed{
\sigma_{\mathrm{loc}}^2(K,T)=\frac{\partial_T C(K,T)}{\tfrac12 K^2\,\partial_{KK}C(K,T)}.
}
\]
\paragraph{Remarque.}
Quand $r$ et $q$ sont non nuls, on obtient une formule plus longue (termes en $\partial_K C$ et $C$). Le message reste : la volatilité locale est un \emph{objet dérivé} de $C(K,T)$, donc très sensible au bruit et à l'interpolation.

\subsection*{(3) Contrôles discrets : convexité et densité via différences finies}
En pratique, on observe $C(K_i,T)$ sur une grille discrète de strikes.
Un test simple de convexité (butterfly arbitrage) utilise des différences secondes :
\[
\Delta^2 C_i \approx \frac{C(K_{i+1},T)-2C(K_i,T)+C(K_{i-1},T)}{(\Delta K)^2}.
\]
La condition «pas d'arbitrage butterfly» impose $\Delta^2 C_i\ge 0$ (à tolérance près).
Si $\Delta^2 C_i<0$ à de nombreux endroits, l'interpolation ou les données sont incohérentes économiquement.

\subsection*{(4) Mini-exercice corrigé}
\paragraph{Exercice --- densité implicite à partir de 3 strikes.}
On observe $C(K_{i-1},T)$, $C(K_i,T)$, $C(K_{i+1},T)$ avec un pas uniforme $\Delta K$.
Donner une approximation de la densité risque-neutre en $K_i$ et expliquer le signe attendu.

\emph{Solution.}
Par Breeden--Litzenberger, $f(K_i)\approx e^{rT}\partial_{KK}C(K_i,T)$.
On approxime $\partial_{KK}C$ par la différence seconde :
\[
f(K_i)\approx e^{rT}\frac{C(K_{i+1},T)-2C(K_i,T)+C(K_{i-1},T)}{(\Delta K)^2}.
\]
Comme une densité est positive, on attend un résultat $\ge 0$ (à bruit près). Un résultat fortement négatif indique un arbitrage butterfly (ou des erreurs de données/interpolation).


\begin{comment}
\section{D. Numérique / algorithmes}
\subsection*{Interpolation dans l'app}
\filepath{app/model/calibration/market_surface.py} :
\begin{itemize}[leftmargin=*]
\item associe chaque point observé au voisin le plus proche sur la grille (remplissage + masque),
\item tente une interpolation \codeid{griddata(method="linear")} si SciPy est disponible,
\item sinon «fallback» par nearest/median.
\end{itemize}

\end{comment}
\begin{comment}
\section{E. Implémentation dans le repo}
\begin{itemize}[leftmargin=*]
\item Grilles par défaut : \codeid{M_GRID_DEFAULT} et \codeid{T_GRID_DEFAULT} dans \filepath{app/model/calibration/market_surface.py}.
\item Construction : \codeid{build_fixed_grid} retourne \codeid{iv_grid}, \codeid{mask}, \codeid{m_grid}, \codeid{t_grid}.
\end{itemize}

\section{F. Exemple end-to-end}
Dans l'onglet «Calibration avancée» (\filepath{app/vue/tabs/tab_advanced_calibration.py}) :
\begin{enumerate}[leftmargin=*]
\item on charge une surface Yahoo (DataFrame avec colonnes \codeid{K,T,S0,iv,type}) via \codeid{CalibrationController.fetch_yahoo_iv_surface} (\filepath{app/controller/calibration_controller.py}) qui appelle \codeid{fetch_iv_surface} (\filepath{app/model/options/data/iv_surface.py}).
\item on canonicalise/filtre dans la Vue, puis on passe au contrôleur une table «surface marché» ;
\item le contrôleur appelle \codeid{make_fixed_grid} pour obtenir (matrice, masque).
\end{enumerate}

\section{G. Limites \& extensions}
\begin{itemize}[leftmargin=*]
\item \NonImplem\ : aucune contrainte explicite d'arbitrage statique (convexité en strike) n'est repérée dans les modules de calibration.
\item Extension crédible : calibration sous contraintes (convexité) ou paramétrisations arbitrage-free (SVI pour equity).
\end{itemize}

\end{comment}
\begin{jurybox}
Une surface d'IV n'est pas une figure : c'est un \textbf{objet de données} (fonction continue ou grille discrète) qui sert d'entrée à un problème inverse. Le point clé est de distinguer points observés vs valeurs interpolées (masque), car cela change la notion d'erreur et la stabilité de la calibration.
\end{jurybox}

\chapter{Calibration : principe, fonction objectif, pondérations, contraintes}

\section{A. Problème financier (objectif, inputs, outputs)}
On observe une surface de volatilité implicite de marché sur une grille (moneyness, maturité) :
\[
\Sigma^{\text{mkt}}{i,j}\approx \sigma{\text{impl}}(m_j,T_i),\qquad m_j=\frac{K_j}{S_0}.
\]
On veut identifier un modèle paramétrique $\mathcal{M}(\theta)$ (SABR, Heston, sauts, rough, \ldots) tel que sa surface
\[
\Sigma^{\text{mod}}(\theta)=\sigma_{\text{impl}}^{\mathcal{M}}(m_j,T_i;\theta)
\]
reproduise au mieux le marché.

\textbf{Inputs} : surface de marché (table de points), $(S_0,r,q)$, une grille $(m_j)$ et $(T_i)$, un choix de modèle, des contraintes (bornes, paramètres fixés), des réglages d'optimisation.
\textbf{Outputs} : paramètres calibrés $\hat\theta$, surface modèle, résidus (marché vs modèle), métriques (MAE/RMSE) et diagnostics de stabilité.

\section{B. Théorie (définitions, hypothèses)}
\subsection*{Définition formelle : calibration comme problème inverse}
On pose un problème d'ajustement (souvent moindres carrés) :
\[
\hat\theta \in \argmin_{\theta\in\Theta}; \mathcal{L}(\theta),
\]
où $\Theta$ encode des contraintes (bornes, signes, \ldots) et $\mathcal{L}$ mesure l'écart marché/modèle.

Deux choix standards :
\begin{itemize}[leftmargin=*]
\item \textbf{fit en volatilités} : $\mathcal{L}(\theta)=\sum w_{i,j}\big(\Sigma^{\text{mod}}{i,j}(\theta)-\Sigma^{\text{mkt}}{i,j}\big)^2$ ;
\item \textbf{fit en prix} : $\mathcal{L}(\theta)=\sum w_{i,j}\big(C^{\text{mod}}{i,j}(\theta)-C^{\text{mkt}}{i,j}\big)^2$,
avec $C^{\text{mkt}}$ reconstruit via BS à partir de l'IV observée.
\end{itemize}

\subsection*{Intuition financière}
\begin{itemize}[leftmargin=*]
\item Calibrer revient à «faire coïncider» le modèle avec la \emph{réalité cotée} (smile/skew) afin que les prix produits soient cohérents (au moins sur les points observés).
\item Le fit en \emph{prix} est souvent plus stable car la sensibilité d'un prix à la vol (Vega) varie fortement selon $K,T$ ; le fit en vol peut sur-pondérer des régions où un petit écart de prix correspond à un gros écart d'IV.
\end{itemize}

\qa
{Transformer une surface observée en paramètres de modèle utilisables pour pricing, scénarios et comparaison (surface marché vs surface modèle).}
{Qualité de la surface de marché ; modèle supposé «assez expressif» ; choix d'une fonction perte pertinente ; stabilité numérique.}
{Non-unicité (plusieurs $\theta$ peuvent fitter) ; minima locaux ; absence d'arbitrage non garantie ; dépendance au poids et au dataset.}
{Au c\oe{}ur : représenter la surface par une grille (et masque), définir une perte (prix vs IV, pondérations, pénalités) et résoudre un problème inverse souvent non convexe sous contraintes.}

\subsection*{Calibration en vol vs calibration en prix : le rôle central de la Vega}
Le choix «fit IV» vs «fit prix» est souvent résumé à une préférence, mais il a une lecture quantitative :
si un prix est perturbé de $\Delta C$, l'erreur correspondante en volatilité est approximativement
\[
\Delta \sigma \approx \frac{\Delta C}{\text{Vega}}.
\]
Donc :
\begin{itemize}[leftmargin=*]
\item minimiser une erreur en vol \emph{revient} à minimiser une erreur en prix \emph{pondérée par} $\text{Vega}^{-1}$ ;
\item près de l'ATM (Vega forte) une erreur de prix se traduit par une petite erreur d'IV ;
\item loin de la monnaie ou pour maturité très courte (Vega faible) l'IV devient numériquement instable.
\end{itemize}
Cette simple relation explique pourquoi de nombreux desks préfèrent calibrer en prix (ou en prix «normalisés») avec des poids adaptés.

\subsection*{Pondérations : bid/ask, Vega, et erreurs relatives}
Un schéma de pondération «idéal» utiliserait l'incertitude de mesure : si un prix de marché est observé avec un écart-type $\sigma_C(K,T)$ (lié au bid/ask), un maximum de vraisemblance gaussien donne des poids $w\propto 1/\sigma_C^2$.
En absence de modèle d'erreur explicite, on utilise souvent :
\begin{itemize}[leftmargin=*]
\item \textbf{erreur relative} : $(C^{\text{mod}}-C^{\text{mkt}})/\max(\varepsilon,C^{\text{mkt}})$ ;
\item \textbf{poids Vega} : erreurs en vol converties en erreurs de prix via Vega ;
\item \textbf{poids bid/ask} : $w\propto 1/(\mathrm{ask}-\mathrm{bid})^2$ (quand disponible).
\end{itemize}

\subsection*{Identifiabilité : pourquoi plusieurs paramètres peuvent «bien fitter»}
Le calibrage est un problème inverse : deux ensembles de paramètres différents peuvent produire des surfaces très proches sur le domaine observé.
Cela arrive surtout quand :
\begin{itemize}[leftmargin=*]
\item la surface est pauvre (peu de maturités, peu de strikes) ;
\item les options sont peu liquides (bruit important) ;
\item le modèle est trop flexible (beaucoup de paramètres).
\end{itemize}
Conséquence : il faut raisonner en termes de \emph{robustesse} (stabilité des paramètres) et pas seulement en termes de RMSE de surface.

% =========================================================
% Compléments M1 — Calibration (B. Théorie)
% =========================================================

\subsection*{(1) Ill-posedness : pourquoi «bien fitter» ne suffit pas}
La calibration est un problème inverse : on observe une surface (prix ou IV) et on remonte à des paramètres $\theta$.
Deux difficultés structurelles apparaissent presque toujours :
\begin{itemize}[leftmargin=*]
\item \textbf{Non convexité} : la fonction perte peut avoir de nombreux minima locaux (surtout avec FFT/CF/MC).
\item \textbf{Mauvais conditionnement} : certains paramètres ont un effet très faible sur la surface sur le domaine observé (directions «plates»), d'autres ont un effet très fort (directions «raides»).
\end{itemize}
Conséquence : deux jeux de paramètres peuvent donner une RMSE quasi identique mais des dynamiques très différentes hors grille (extrapolation, stress, hedging).

\subsection*{(2) Jacobien et «sloppiness» : lecture linéarisée du problème}
Notons $r(\theta)$ le vecteur des résidus (par exemple $C^{\mathrm{mod}}-C^{\mathrm{mkt}}$ sur les points observés).
Près d'un $\theta_0$, on linéarise :
\[
r(\theta_0+\Delta\theta)\approx r(\theta_0)+J(\theta_0)\Delta\theta,
\qquad
J_{ij}=\frac{\partial r_i}{\partial \theta_j}.
\]
Si $J^\top J$ a un spectre très étalé (grand rapport conditionnement), le problème est «sloppy» :
\begin{itemize}[leftmargin=*]
\item des directions de paramètres changent peu les prix $\Rightarrow$ paramètres instables ;
\item de petites erreurs de données peuvent provoquer de grandes variations sur ces directions.
\end{itemize}
Cette lecture explique pourquoi une calibration n'est pas seulement un «optimiseur qui converge» : c'est une question de \emph{géométrie}.

\subsection*{(3) Régularisation : stabiliser, au prix d'un biais contrôlé}
Une régularisation de type Tikhonov ajoute un terme :
\[
\mathcal{L}_{\mathrm{reg}}(\theta)=\mathcal{L}(\theta)+\alpha\|\theta-\theta_{\mathrm{prior}}\|^2,
\qquad \alpha>0.
\]
Interprétation :
\begin{itemize}[leftmargin=*]
\item \textbf{fréquentiste} : on pénalise les paramètres extrêmes pour améliorer le conditionnement ;
\item \textbf{bayésien} : c'est un MAP avec un prior gaussien centré en $\theta_{\mathrm{prior}}$.
\end{itemize}
En finance, cela correspond à une pratique raisonnable : préférer un modèle «plausible» et stable, plutôt qu'un fit parfait mais fragile.

\subsection*{(4) Contraintes : bornes, positivité, et conditions structurelles}
Les contraintes jouent deux rôles :
\begin{itemize}[leftmargin=*]
\item \textbf{économique} : éviter des dynamiques absurdes (variance négative, explosions).
\item \textbf{numérique} : éviter des zones où les pricers deviennent instables (intégrales oscillantes, CF mal conditionnée, etc.).
\end{itemize}
Exemples :
\begin{itemize}[leftmargin=*]
\item Heston : $\kappa>0$, $\theta>0$, $\sigma>0$, $v_0>0$, $|\rho|<1$, et (optionnel) pénalité Feller.
\item SABR : $\alpha>0$, $\nu>0$, $|\rho|<1$, et bornes sur $\beta$.
\end{itemize}

\subsection*{(5) Calibration «en prix» vs «en vol» : une relecture via la Vega}
La relation linéarisée $\,\Delta C\approx \mathrm{Vega}\cdot \Delta\sigma\,$ dit :
\[
\Delta\sigma \approx \frac{\Delta C}{\mathrm{Vega}}.
\]
Donc, minimiser une erreur en IV revient approximativement à minimiser une erreur en prix \emph{pondérée} par $1/\mathrm{Vega}^2$ :
\[
\sum (\Delta\sigma)^2 \approx \sum \left(\frac{\Delta C}{\mathrm{Vega}}\right)^2.
\]
Message : le choix «prix vs IV» est surtout un choix de \emph{pondération implicite} des régions de la surface.


\section{C. Dérivations (ultra détaillées)}
\subsection*{De la surface IV au prix de marché utilisé en calibration (fit en prix)}
Sur un point $(K,T)$, si l'on observe $\sigma_{\text{impl}}^{\text{mkt}}$, on peut définir un «prix marché cohérent» :
\[
C^{\text{mkt}}(K,T)=C_{\text{BS}}\big(S_0,K,r,q,T,\sigma_{\text{impl}}^{\text{mkt}}(K,T)\big).
\]
Cette transformation est standard : elle convertit une observation en IV (cotation) en un prix de référence, ce qui permet d'exprimer une perte en prix même quand le marché est fourni sous forme d'IV.

\subsection*{Pourquoi une pondération par l'échelle de prix ?}
Une perte brute $\sum (C^{\text{mod}}-C^{\text{mkt}})^2$ sur-pondère mécaniquement les prix élevés (ATM/ITM, maturités longues). Une normalisation simple est :
\[
\sum \Big(\frac{C^{\text{mod}}-C^{\text{mkt}}}{\text{scale}}\Big)^2,\qquad \text{scale}\approx \max(\varepsilon, C^{\text{mkt}}).
\]
Ce choix transforme une erreur absolue en une erreur \emph{relative} (avec plancher $\varepsilon$), ce qui évite qu'une poignée de points très chers domine la calibration.

\subsection*{Contraintes et pénalités : exemple Heston et condition de Feller}
Dans Heston, une condition suffisante (non nécessaire) pour que la variance reste strictement positive est la condition de Feller :
\[
2\kappa\theta \ge \sigma^2.
\]
Plutôt que d'imposer une contrainte dure, on peut ajouter une \emph{pénalité douce} dans la perte (ou dans les résidus) :
\[
\mathcal{L}(\theta)\gets \mathcal{L}(\theta) + \lambda\,\max(0,\sigma^2-2\kappa\theta)^2,
\]
où $\lambda>0$ règle l'intensité de la pénalisation. L'intérêt : guider l'optimiseur loin des zones numériquement instables, sans rendre le problème irréalisable si le marché «force» légèrement la contrainte.

\subsection*{Multi-start : pourquoi c'est nécessaire ?}
Les pertes sont non convexes (FFT/CF/MC + inversion IV) : on démarre l'optimisation depuis plusieurs initialisations $\theta^{(0)}$ et on garde la meilleure. Deux stratégies typiques :
\begin{itemize}[leftmargin=*]
\item \textbf{multi-start} : plusieurs points de départ (heuristiques + tirages aléatoires) ;
\item \textbf{design d'expériences} : explorer un espace de paramètres (souvent en log-échelle) puis raffiner autour des meilleurs candidats.
\end{itemize}
Le point M1 : le multi-start n'est pas un détail d'implémentation, c'est une réponse directe à la géométrie non convexe du problème inverse.

\subsection*{Lien quantitatif : convertir une erreur de prix en erreur d'IV}
Pour comprendre les métriques (RMSE en IV vs en prix), on peut linéariser Black--Scholes :
\[
C_{\text{BS}}(\sigma+\Delta\sigma)\approx C_{\text{BS}}(\sigma)+\text{Vega}\cdot \Delta\sigma.
\]
Donc, à premier ordre,
\[
\Delta\sigma \approx \frac{\Delta C}{\text{Vega}}.
\]
Cela justifie une pratique courante : si l'on calibre en prix mais que l'on veut comparer en IV, on doit tenir compte de la Vega (sinon on «sur-pondère» des points où l'IV est peu informative).

\subsection*{Robustesse : pertes robustes et pénalisation}
Les surfaces de marché contiennent des outliers (bid/ask large, erreurs de parsing, strikes aberrants).
Deux outils standards (au-delà du simple $L^2$) :
\begin{itemize}[leftmargin=*]
\item \textbf{perte de Huber} : quadratique près de 0, linéaire dans les queues ;
\item \textbf{régularisation} : ajouter des termes de pénalité sur des paramètres (ex. éloignement d'un prior raisonnable, violation de contraintes).
\end{itemize}
Ils stabilisent les paramètres au prix d'un léger biais, souvent acceptable en pratique.

% =========================================================
% Compléments M1 — Calibration (C. Dérivations)
% =========================================================

\subsection*{(1) Moindres carrés pondérés : gradient et Hessien (lecture Gauss--Newton)}
Considérons une perte de type moindres carrés pondérés :
\[
\mathcal{L}(\theta)=\frac12\sum_{i=1}^N w_i\,r_i(\theta)^2
=\frac12\,r(\theta)^\top W\,r(\theta),
\]
où $r_i(\theta)$ est un résidu (prix ou IV) et $W=\mathrm{diag}(w_i)$.
Notons $J(\theta)$ le jacobien $J_{ij}=\partial r_i/\partial \theta_j$.

\paragraph{Gradient.}
En dérivant :
\[
\nabla \mathcal{L}(\theta)=J(\theta)^\top W\,r(\theta).
\]
\paragraph{Hessien exact.}
Il s'écrit
\[
\nabla^2\mathcal{L}(\theta)=J^\top WJ + \sum_{i=1}^N w_i r_i(\theta)\,\nabla^2 r_i(\theta).
\]
\paragraph{Approximation Gauss--Newton.}
Près d'un bon fit, les résidus $r_i$ sont petits, donc on néglige le second terme :
\[
\nabla^2\mathcal{L}(\theta)\approx J^\top WJ.
\]
Cela justifie beaucoup d'algorithmes pratiques : on approxime la courbure par $J^\top WJ$.

\subsection*{(2) Levenberg--Marquardt : compromis entre Newton et descente de gradient}
Gauss--Newton propose de résoudre
\[
(J^\top WJ)\Delta\theta = -J^\top Wr.
\]
Si $J^\top WJ$ est mal conditionné, on ajoute un terme de régularisation (damping) :
\[
(J^\top WJ+\lambda I)\Delta\theta = -J^\top Wr,\qquad \lambda>0.
\]
\paragraph{Interprétation.}
\begin{itemize}[leftmargin=*]
\item $\lambda$ petit : comportement proche de Newton/Gauss--Newton (rapide si bien conditionné).
\item $\lambda$ grand : comportement proche d'une descente de gradient (robuste mais lent).
\end{itemize}
En calibration, ce mécanisme est crucial : les surfaces sont bruitées et les pricers peuvent être irréguliers, donc on a besoin de robustesse.

\subsection*{(3) Log-paramétrisation : rendre les contraintes «naturelles»}
Beaucoup de paramètres doivent être positifs : $\kappa,\theta,\sigma,v_0,\alpha,\nu,\dots$.
Plutôt que d'imposer des bornes dures, on peut optimiser sur des variables libres :
\[
\kappa=\exp(\widetilde\kappa),\quad \theta=\exp(\widetilde\theta),\quad \sigma=\exp(\widetilde\sigma),\ \dots
\]
Cette reparamétrisation :
\begin{itemize}[leftmargin=*]
\item garantit la positivité ;
\item homogénéise souvent les échelles ;
\item améliore le conditionnement (car les paramètres sont naturellement en échelle log en finance).
\end{itemize}

\subsection*{(4) Mini-exercice corrigé : «fit en vol» = «fit en prix pondéré» (au 1er ordre)}
Supposons que le marché soit donné en IV et qu'on compare des IV.
Montrer qu'au premier ordre, minimiser $\sum (\sigma^{\mathrm{mod}}-\sigma^{\mathrm{mkt}})^2$ revient à minimiser une somme de carrés de prix \emph{pondérée} par $1/\mathrm{Vega}^2$.

\emph{Solution.}
Sur un point, on a la linéarisation BS :
\[
C(\sigma^{\mathrm{mod}})\approx C(\sigma^{\mathrm{mkt}})+\mathrm{Vega}(\sigma^{\mathrm{mkt}})\cdot(\sigma^{\mathrm{mod}}-\sigma^{\mathrm{mkt}}).
\]
Donc
\[
\Delta C \approx \mathrm{Vega}\cdot \Delta\sigma
\quad\Rightarrow\quad
\Delta\sigma \approx \frac{\Delta C}{\mathrm{Vega}}.
\]
En sommant sur les points :
\[
\sum (\Delta\sigma)^2 \approx \sum \left(\frac{\Delta C}{\mathrm{Vega}}\right)^2,
\]
soit un fit en prix avec poids $w\propto 1/\mathrm{Vega}^2$.
Cela explique pourquoi «fit en vol non pondéré» sur-pondère mécaniquement les zones où Vega est faible.


\begin{comment}
\section{D. Numérique / algorithmes (stabilité, complexité, choix)}
\subsection*{Trois familles d'objectifs dans le dépôt}
\begin{enumerate}[leftmargin=*]
\item \textbf{Least squares en IV} (SABR) : rapide, mais dépend fortement du choix de paramétrisation d'IV (Hagan).
\item \textbf{Least squares en prix via FFT} (Heston/Merton/Bates/rHeston) : coût $O(N_T \cdot N \log N)$ pour produire une grille de prix (FFT par maturité), puis interpolation, puis inversion IV.
\item \textbf{MC design / proxy} (rBergomi/Volterra) : très coûteux, bruit stochastique ; le dépôt utilise une stratégie «surrogate» (rBergomi) ou «best-of-design» (Volterra).
\end{enumerate}

\subsection*{Dépendances optionnelles}
Plusieurs calibrateurs utilisent SciPy (\codeid{least_squares}, \codeid{griddata}, \codeid{RBFInterpolator}). Ils gèrent explicitement l'absence de SciPy en renvoyant une erreur contrôlée (voir, par exemple, \filepath{app/model/volatility_models/sabr/calibrator.py}).

\end{comment}
\begin{comment}
\section{E. Implémentation dans le repo (fichiers, fonctions, pipeline, paramètres)}
\subsection*{API unifiée de calibration (V2) : \codeid{SurfaceGrid} + calibrateurs}
L'objet d'entrée canonique est \codeid{SurfaceGrid} (dataclass) dans \filepath{app/model/calibration/base_calibrator.py}. Les réglages standards sont \codeid{CalibratorSettings} (\codeid{fit_to_observed_only}, \codeid{max_nfev}, \codeid{n_starts}, \codeid{seed}).

Le point d'orchestration côté UI (onglet «Calibration avancée») est :
\begin{itemize}[leftmargin=*]
\item \filepath{app/controller/calibration_controller.py}, méthode \codeid{run_advanced_surface_calibration},
\end{itemize}
qui route vers des calibrateurs concrets (SABR, Heston FFT, Merton, Bates, rHeston, rBergomi, Volterra).

\subsection*{Extrait de code : construction des bornes (contraintes) pour Heston V1}
Dans \filepath{app/model/calibration/heston_calibrator.py}, \codeid{build_bounds} traduit un dictionnaire JSON en bornes (et «fixe» un paramètre si besoin). Extrait :

\begin{lstlisting}[caption={Contraintes $\to$ bornes : \filepath{app/model/calibration/heston_calibrator.py}, \codeid{build_bounds} (extrait).},label={lst:build-bounds}]
def build_bounds(constraints: Dict[str, Any] | None = None):
lower = {k: float(v[0]) for k, v in DEFAULT_BOUNDS.items()}
upper = {k: float(v[1]) for k, v in DEFAULT_BOUNDS.items()}

if isinstance(constraints, dict):
    for name in PARAM_ORDER:
        if name not in constraints:
            continue
        mn, mx, fixed = _parse_param_constraint(constraints.get(name))
        if fixed is not None:
            lower[name] = fixed
            upper[name] = fixed
            continue
        if mn is not None:
            lower[name] = max(lower[name], mn)
        if mx is not None:
            upper[name] = min(upper[name], mx)

lb = np.array([lower[p] for p in PARAM_ORDER], dtype=float)
ub = np.array([upper[p] for p in PARAM_ORDER], dtype=float)
...
\end{lstlisting}

\paragraph{Pourquoi c'est la traduction directe de la théorie «contraintes» ?}
\begin{itemize}[leftmargin=*]
\item Le JSON «\codeid{"kappa":[min,max]}» devient une contrainte $\kappa\in[\kappa_{\min},\kappa_{\max}]$.
\item Un scalaire («\codeid{"rho": -0.7}») est interprété comme $\rho$ \emph{fixé} (min=max).
\item La boucle sur \codeid{PARAM_ORDER} garantit un ordre stable des paramètres dans l'optimiseur.
\end{itemize}

\subsection*{Attention : deux frameworks de calibration coexistent}
\begin{itemize}[leftmargin=*]
\item Framework V1 (scaffold) : \filepath{app/model/calibration/logic.py} (\codeid{run_calibration}) retourne explicitement \emph{Calibration non implémentée}.
\item Framework V2 (fonctionnel) : \filepath{app/controller/calibration_controller.py}, \codeid{run_advanced_surface_calibration} utilise des calibrateurs dans \filepath{app/model/volatility_models/}.
\end{itemize}

\section{F. Exemple end-to-end (lancement utilisateur)}
\begin{enumerate}[leftmargin=*]
\item L'utilisateur sélectionne un modèle (ex. \codeid{"heston_fft"}) dans \filepath{app/vue/tabs/tab_advanced_calibration.py}.
\item Il charge une surface et clique «Calibrate» : la Vue construit un payload et appelle \codeid{CalibrationController.run_advanced_surface_calibration}.
\item Le contrôleur :
\begin{itemize}[leftmargin=*]
\item charge \codeid{market_df} via \codeid{load_market_surface_csv_v2} (parsing),
\item calcule \codeid{iv_market, mask} via \codeid{make_fixed_grid},
\item instancie le calibrateur choisi,
\item renvoie \codeid{params, iv_model, iv_error, metrics}.
\end{itemize}
\item La Vue affiche \codeid{metrics} et les surfaces (marché/modèle/erreur).
\end{enumerate}

\section{G. Limites \& extensions crédibles}
\begin{itemize}[leftmargin=*]
\item \textbf{Arbitrage statique} : le dépôt ne contraint pas explicitement l'absence d'arbitrage (convexité en strike). Extension : intégrer des contraintes ou une paramétrisation arbitrage-free.
\item \textbf{Pondérations} : l'échelle \codeid{scale=max(1e-4, px)} est utile mais simpliste. Extension : poids $\propto$ Vega$^{-1}$ ou erreurs de bid-ask.
\item \textbf{Robustesse} : ajout de pertes robustes (Huber) et de diagnostics (outliers).
\end{itemize}

\end{comment}
\begin{jurybox}
La calibration est un \textbf{problème inverse} non convexe. À retenir : (i) représenter la surface de façon standardisée (grille + masque), (ii) imposer des bornes/contraintes, (iii) utiliser du multi-start, (iv) choisir une perte adaptée (prix ou IV, pondérations), et (v) ajouter des pénalités (ex. Feller) pour stabiliser l'optimisation.
\end{jurybox}

\chapter{Heston : volatilité stochastique (SDE, CF, pricing, calibration)}

\section{A. Problème financier (objectif, inputs, outputs)}
Sur actions, l'IV n'est pas plate : elle présente un skew marqué. Le modèle de Heston vise à capturer :
\begin{itemize}[leftmargin=*]
\item la \textbf{volatilité stochastique} (variance qui varie dans le temps),
\item une \textbf{corrélation} $\rho<0$ entre chocs de prix et chocs de variance (effet leverage),
\item une \textbf{reversion} de la variance vers un niveau de long terme.
\end{itemize}
\textbf{Inputs} : surface de marché, $(S_0,r,q)$, paramètres $\theta=(\kappa,\theta,\sigma,\rho,v_0)$ (attention : ici $\theta$ désigne le niveau de variance de long terme).
\textbf{Outputs} : prix (ou IV) par $(K,T)$, paramètres calibrés, surfaces modèle/erreur.

\section{B. Théorie (définitions, hypothèses)}
\subsection*{SDE sous la mesure risque-neutre}
Sous $\Qmeas$, Heston s'écrit :
\[
\frac{\dd S_t}{S_t}=(r-q)\dd t+\sqrt{V_t},\dd W_t^{(1)},\qquad
\dd V_t=\kappa(\theta - V_t)\dd t+\sigma\sqrt{V_t},\dd W_t^{(2)},
\]
avec $\Corr(\dd W^{(1)},\dd W^{(2)})=\rho$.

\subsection*{Hypothèses}
\begin{itemize}[leftmargin=*]
\item $V_t$ suit un processus de type CIR (racine carrée), reversion vers $\theta$.
\item Marché frictionless, taux $r$ et yield $q$ constants (dans les pricers de surface).
\item Pour la positivité stricte, condition de Feller : $2\kappa\theta\ge\sigma^2$ (souvent utilisée comme contrainte ou pénalité en calibration).
\end{itemize}

\subsection*{Intuition financière}
\begin{itemize}[leftmargin=*]
\item $\kappa$ : vitesse de retour de la vol vers «normal».
\item $\theta$ : niveau de vol long terme ($\sqrt{\theta}$).
\item $v_0$ : vol instantanée aujourd'hui.
\item $\rho<0$ : quand $S$ baisse, $V$ monte, générant un skew (puts chers).
\end{itemize}

\qa
{Modéliser la surface d'IV au-delà de BS, produire des prix cohérents (au moins pour calls européens) et calibrer des paramètres interprétables.}
{Vol stochastique de type CIR, continuité (pas de sauts), risque neutralité standard, stabilité numérique des intégrales/FFT.}
{Le modèle peut échouer sur des smiles très marqués ou des événements (earnings) ; calibration instable si surface pauvre ; absence d'arbitrage non garantie.}
{Deux approches de pricing sont classiques : (i) intégration 1D semi-fermée (Lewis/Heston) via la fonction caractéristique, (ii) FFT de Carr--Madan sur une grille de log-strikes pour obtenir de nombreux prix en une passe.}

\subsection*{Zoom sur la variance CIR : positivité, stationnarité, interprétation}
La variance suit une dynamique de type CIR :
\[
\dd V_t=\kappa(\theta-V_t)\dd t+\sigma\sqrt{V_t}\,\dd W_t^{(2)}.
\]
Quelques faits structurants :
\begin{itemize}[leftmargin=*]
\item \textbf{Réversion vers $\theta$.} Plus $\kappa$ est grand, plus $V_t$ revient vite vers $\theta$.
\item \textbf{Vol-of-vol $\sigma$.} Plus $\sigma$ est grand, plus la variance fluctue, ce qui accentue la courbure de la surface (smile).
\item \textbf{Condition de Feller.} La condition $2\kappa\theta\ge\sigma^2$ est suffisante pour éviter que $V_t$ touche 0 (strictement positif). En pratique, même si elle est violée, on peut encore pricer, mais certaines formules deviennent numériquement plus délicates.
\item \textbf{Stationnarité (intuition).} À long terme (sous conditions), $V_t$ admet une loi stationnaire de type Gamma : cela justifie l'idée d'une «vol long terme» autour de $\sqrt{\theta}$.
\end{itemize}

\subsection*{Pourquoi $\rho<0$ crée un skew}
Le terme $\rho$ corrèle les chocs de prix et de variance.
Si $\rho<0$, une baisse de $S$ tend à coïncider avec une hausse de $V$ : les scénarios de baisse sont plus volatils, ce qui rend les puts OTM plus chers et génère une IV plus élevée pour $K<S_0$ (skew).

\subsection*{Attention : paramètres sous $\Pmeas$ vs sous $\Qmeas$}
En général, la variance porte une \emph{prime de risque} : les paramètres de la dynamique sous $\Pmeas$ (historique) ne coïncident pas avec ceux sous $\Qmeas$ (pricing).
Le calibrage sur une surface d'options identifie donc des paramètres «pricing», pas forcément des paramètres «statistiques».

% =========================================================
% Compléments M1 — Heston (B. Théorie)
% =========================================================

\subsection*{(1) Pourquoi Heston est le «premier modèle sérieux» pour le skew equity}
Le smile/skew equity n'est pas seulement une irrégularité : il reflète des asymétries de distribution (crash risk, leverage effect).
Heston est un modèle minimal qui capture trois ingrédients :
\begin{itemize}[leftmargin=*]
\item \textbf{Volatilité stochastique} : la variance $V_t$ varie dans le temps, donc la distribution de $S_T$ est un mélange de log-normales.
\item \textbf{Leverage effect} : $\rho<0$ fait que les baisses de $S$ s'accompagnent de hausses de $V$.
\item \textbf{Mean-reversion} : $V_t$ revient vers $\theta$ à vitesse $\kappa$.
\end{itemize}

\subsection*{(2) Lecture des paramètres : «si je bouge tel paramètre, la surface fait quoi ?»}
En première approximation (à surface de marché fixée) :
\begin{itemize}[leftmargin=*]
\item $\kappa$ (vitesse) : grand $\Rightarrow$ la variance revient vite vers $\theta$ ; effet sur la structure en maturité (terme court vs long).
\item $\theta$ (niveau LT) : détermine la variance moyenne à long terme ; contrôle l'ATM long.
\item $\sigma$ (vol-of-vol) : plus $\sigma$ est grand, plus le smile est marqué (plus de dispersion de $V$).
\item $\rho$ : contrôle l'asymétrie (skew) ; $\rho<0$ accentue la pente côté puts.
\item $v_0$ : contrôle le niveau court terme (ATM court).
\end{itemize}
Cette lecture qualitative est cruciale en calibration : elle aide à choisir des bornes, des initialisations et à diagnostiquer une solution «non plausible».

\subsection*{(3) Condition de Feller : utilité pratique et limites}
La condition $2\kappa\theta\ge\sigma^2$ est suffisante pour garantir que $V_t$ reste strictement positive.
\paragraph{Pourquoi c'est utile ?}
Parce que les pricers basés sur CF/FFT peuvent devenir plus instables quand $V$ touche 0 (oscillations, pertes de précision).
\paragraph{Pourquoi ce n'est pas une loi ?}
Le marché peut pousser vers des paramètres qui violent légèrement Feller tout en gardant une dynamique «raisonnable» en pratique.
On l'utilise souvent comme pénalité douce plutôt que comme contrainte dure.

\subsection*{(4) Pourquoi une CF (fonction caractéristique) suffit pour pricer une vanilla}
Pour un payoff européen, le prix est une espérance sous $\Qmeas$.
Si l'on connaît la fonction caractéristique du log-prix $X_T=\ln S_T$ :
\[
\phi(u)=\E^{\Qmeas}[e^{iuX_T}],
\]
alors on peut reconstruire la loi (densité) et donc le prix par inversion de Fourier.
Deux approches classiques :
\begin{itemize}[leftmargin=*]
\item formulation $P_1/P_2$ : deux intégrales 1D donnant $C=S_0e^{-qT}P_1-Ke^{-rT}P_2$ ;
\item Carr--Madan : transformer le call en une fonction intégrable et FFT sur une grille de strikes.
\end{itemize}
\paragraph{Message M1.}
La CF est un «résumé» plus stable que la densité : elle évite de dériver numériquement et se combine bien avec des méthodes rapides (FFT).


\section{C. Dérivations (ultra détaillées)}
\subsection*{Pourquoi une fonction caractéristique ?}
Pour un payoff européen, le prix est une espérance sous $\Qmeas$. Plutôt que de simuler $(S_T,V_T)$, Heston admet une forme semi-analytique via la \emph{fonction caractéristique} du log-prix (ou log-return) :
\[
\phi(u)=\E^{\Qmeas}\big[e^{iu\ln(S_T)}\big]\quad \text{ou}\quad \E^{\Qmeas}\big[e^{iu\ln(S_T/S_0)}\big].
\]
Une fois $\phi$ connue, on peut obtenir le prix d'un call via une intégrale (Carr--Madan ou formulation $P_1/P_2$).

\subsection*{Formule $P_1/P_2$ (idée)}
Le prix d'un call peut s'écrire
\[
C = S_0e^{-qT}P_1 - Ke^{-rT}P_2,
\]
où $P_1$ et $P_2$ sont des probabilités «sous mesure transformée» calculées par intégrales de $\phi(u)$.
Cette décomposition met en évidence l'analogie avec Black--Scholes : un call est une différence entre une probabilité de «finir ITM» ($P_2$) et une probabilité pondérée ($P_1$), toutes deux calculables via la fonction caractéristique.

\subsection*{Carr--Madan (FFT) : principe}
Carr--Madan propose de calculer simultanément $C(K)$ sur un \emph{grillage de log-strikes} par FFT. On applique une atténuation $e^{\alpha k}$ (avec $k=\ln(K/S_0)$) pour rendre la transformée intégrable, puis on FFT.

\subsection*{Affinité et équations de Riccati (idée de dérivation de $\phi$)}
Heston est un modèle \emph{affine} : le log-prix et la variance entrent linéairement dans l'exposant de la fonction caractéristique.
On cherche une forme
\[
\phi(u)=\E^{\Qmeas}\big[e^{iu\ln S_T}\mid \F_t\big]=\exp\big(A(u,\tau)+B(u,\tau)V_t+iu\ln S_t\big),\quad \tau=T-t.
\]
En injectant cette forme dans l'équation de Kolmogorov (ou en appliquant Itô à l'exponentielle et en imposant la martingale), on obtient un système d'ODE couplées pour $A$ et $B$ (équations de Riccati) qui se résout analytiquement.
C'est ce qui rend Heston «semi-fermé» : la difficulté est déplacée dans une intégrale 1D (ou une FFT), plutôt que dans une simulation 2D.

\subsection*{Pourquoi un facteur d'atténuation $\alpha$ en FFT ?}
Le prix d'un call en fonction du log-strike $k$ ne décroît pas assez vite pour que sa transformée de Fourier soit intégrable.
Carr--Madan propose d'étudier $e^{\alpha k}C(e^k)$ pour un $\alpha>0$ adéquat, ce qui :
\begin{itemize}[leftmargin=*]
\item garantit l'intégrabilité ;
\item stabilise numériquement la FFT ;
\item introduit un compromis : trop petit $\alpha$ $\Rightarrow$ divergence ; trop grand $\alpha$ $\Rightarrow$ amplification d'erreurs numériques.
\end{itemize}
Le paramètre $\alpha$ est donc un hyperparamètre \emph{numérique} autant que mathématique.

% =========================================================
% Compléments M1 — Heston (C. Dérivations)
% =========================================================

\subsection*{(1) Réécrire la dynamique en log-prix : étape indispensable}
Sous $\Qmeas$, Heston est :
\[
\frac{\dd S_t}{S_t}=(r-q)\dd t+\sqrt{V_t}\,\dd W_t^{(1)},
\qquad
\dd V_t=\kappa(\theta-V_t)\dd t+\sigma\sqrt{V_t}\,\dd W_t^{(2)},
\qquad
\dd W^{(1)}\dd W^{(2)}=\rho\,\dd t.
\]
Posons $X_t=\ln S_t$. Par Itô (chapitre Itô) :
\[
\dd X_t=\frac{1}{S_t}\dd S_t-\frac12\frac{1}{S_t^2}(\dd S_t)^2
\]
or $(\dd S_t)^2=(\sqrt{V_t}S_t)^2(\dd W_t^{(1)})^2=V_t S_t^2\,\dd t$. Donc
\[
\boxed{\dd X_t=\Big(r-q-\tfrac12 V_t\Big)\dd t+\sqrt{V_t}\,\dd W_t^{(1)}.}
\]
Le couple $(X_t,V_t)$ est un processus de diffusion \emph{affine} (les coefficients sont au plus linéaires en $V$).

\subsection*{(2) La fonction caractéristique : définir l'objet et l'équation qu'il satisfait}
On cherche (pour $u\in\R$) la fonction caractéristique conditionnelle :
\[
\phi(t,x,v;u)=\E^{\Qmeas}\!\big[e^{iuX_T}\mid X_t=x,\ V_t=v\big].
\]
Par la théorie de Kolmogorov (ou Feynman--Kac), $\phi$ satisfait une PDE rétrograde :
\[
\partial_t \phi + \mathcal{L}\phi=0,
\qquad \phi(T,x,v;u)=e^{iux},
\]
où $\mathcal{L}$ est le générateur de $(X,V)$.

\subsection*{(3) Calcul du générateur $\mathcal{L}$ et ansatz affine}
À partir de
\[
\dd X=(r-q-\tfrac12 v)\dd t+\sqrt{v}\,\dd W^{(1)},
\qquad
\dd V=\kappa(\theta-v)\dd t+\sigma\sqrt{v}\,\dd W^{(2)},
\]
et $\dd W^{(1)}\dd W^{(2)}=\rho\,\dd t$, on obtient
\[
\mathcal{L}f
=(r-q-\tfrac12 v)\,\partial_x f+\kappa(\theta-v)\,\partial_v f
\tfrac12 v\,\partial_{xx}f+\tfrac12\sigma^2 v\,\partial_{vv}f+\rho\sigma v\,\partial_{xv}f.
\]
L'idée clé de Heston : chercher une solution de la forme affine
\[
\phi(t,x,v;u)=\exp\big(A(\tau;u)+B(\tau;u)\,v+iux\big),\qquad \tau=T-t.
\]
\paragraph{Dérivées.}
Si $\phi=\exp(A+Bv+iux)$, alors
\[
\partial_x\phi=iu\phi,\quad \partial_{xx}\phi=-u^2\phi,\quad
\partial_v\phi=B\phi,\quad \partial_{vv}\phi=B^2\phi,\quad
\partial_{xv}\phi=iuB\phi.
\]
\paragraph{Injection dans la PDE.}
On a $\partial_t\phi=-A'(\tau)\phi-B'(\tau)v\phi$ (car $\tau=T-t$).
En divisant la PDE par $\phi$ et en identifiant les coefficients constants et ceux en $v$, on obtient deux ODE :
\[
\boxed{
\begin{aligned}
B'(\tau)
&=\tfrac12\sigma^2 B(\tau)^2+(\rho\sigma iu-\kappa)B(\tau)-\tfrac12(u^2+iu),\\
A'(\tau)
&=iu(r-q)+\kappa\theta\,B(\tau),
\end{aligned}}
\qquad B(0)=0,\ A(0)=0.
\]
La première équation est une \textbf{Riccati} (non linéaire) ; la seconde est ensuite une intégration.

\subsection*{(4) Solution fermée de la Riccati : forme standard (à connaître)}
On pose
\[
d=\sqrt{(\kappa-\rho\sigma iu)^2+\sigma^2(u^2+iu)},
\qquad
g=\frac{\kappa-\rho\sigma iu-d}{\kappa-\rho\sigma iu+d}.
\]
Alors la solution s'écrit (forme «stable» utilisée en pratique) :
\[
\boxed{
B(\tau)
=\frac{\kappa-\rho\sigma iu-d}{\sigma^2}\cdot \frac{1-e^{-d\tau}}{1-ge^{-d\tau}}.
}
\]
Puis
\[
\boxed{
A(\tau)
=iu(r-q)\tau+\frac{\kappa\theta}{\sigma^2}\Big((\kappa-\rho\sigma iu-d)\tau-2\ln\frac{1-ge^{-d\tau}}{1-g}\Big).
}
\]
\paragraph{Remarque (branche complexe).}
Le choix de la racine carrée complexe et du logarithme est délicat numériquement (``little Heston trap'').
Le message M1 : la formule est semi-fermée, mais la stabilité dépend des conventions de branche.

\subsection*{(5) De la CF au prix : $P_1/P_2$ (structure) et intuition}
À $t=0$, $x=\ln S_0$ et $v=v_0$, donc
\[
\phi(u)=\exp\big(A(T;u)+B(T;u)v_0+iu\ln S_0\big).
\]
Le prix d'un call admet la structure :
\[
\boxed{
C=S_0e^{-qT}P_1-Ke^{-rT}P_2,
}
\]
où $P_1$ et $P_2$ sont calculés par des intégrales 1D impliquant $\phi$ (inversion de Fourier).
\paragraph{Intuition.}
La décomposition ressemble à BS :
\begin{itemize}[leftmargin=*]
\item $P_2$ joue le rôle d'une probabilité risque-neutre d'être ITM ;
\item $P_1$ est une probabilité sous une mesure «tiltée» (pondérée par $S_T$), d'où le terme $S_0e^{-qT}$.
\end{itemize}
Ce n'est pas qu'une analogie : on le démontre rigoureusement en manipulations de Fourier.

\subsection*{(6) Mini-exercice corrigé : retrouver BS comme cas limite}
\paragraph{Exercice.}
Montrer que si $V_t\equiv \bar v$ est constant (pas de vol-of-vol, pas de stochasticité de variance), alors $S$ suit un GBM de volatilité $\sqrt{\bar v}$ et on retrouve Black--Scholes.

\emph{Solution.}
Si $V_t\equiv\bar v$, la SDE devient
\[
\frac{\dd S_t}{S_t}=(r-q)\dd t+\sqrt{\bar v}\,\dd W_t,
\]
qui est exactement un GBM de volatilité $\sigma=\sqrt{\bar v}$.
Le pricing risque-neutre coïncide alors avec BS : la surface est plate (pas de skew), ce qui illustre que le skew vient de la stochasticité de $V$ et de la corrélation $\rho$.


\begin{comment}
\section{D. Numérique / algorithmes (stabilité, complexité, choix)}
\subsection*{Moteur V1 (intégrale numérique directe)}
Dans \filepath{app/model/calibration/heston_pricer.py}, l'intégrale est approchée par une somme de Riemann sur $u\in(0,u_{\max}]$ :
\[
\int_0^{u_{\max}} g(u),\dd u \approx \Delta u \sum_{n=1}^{N} g(n\Delta u),\qquad \Delta u=\frac{u_{\max}}{N}.
\]
Paramètres numériques exposés côté calibration : \codeid{u_max}, \codeid{n_integration}.

\subsection*{Moteur FFT (Carr--Madan)}
Complexité par maturité : $O(n\log n)$, où \codeid{n} est la taille FFT (par défaut \codeid{n=2048}). Les hyperparamètres sont regroupés dans \codeid{FFTConfig(alpha,n,eta)} dans \filepath{app/model/volatility_models/common/fft.py}.

\end{comment}
\begin{comment}
\section{E. Implémentation dans le repo (fichiers, fonctions, pipeline, paramètres)}
\subsection*{Heston V1 : CF et prix}
\begin{itemize}[leftmargin=*]
\item Fonction caractéristique : \codeid{heston_cf} dans \filepath{app/model/calibration/heston_pricer.py}.
\item Prix call semi-fermé : \codeid{call_price_cf} dans \filepath{app/model/calibration/heston_pricer.py}.
\item Grille de prix (moneyness $\times$ maturité) : \codeid{price_grid_from_params} dans le même fichier.
\item Calibration least squares (SciPy) : \codeid{calibrate_heston_least_squares} dans \filepath{app/model/calibration/heston_calibrator.py}.
\end{itemize}

\subsection*{Heston FFT : CF du log-return + Carr--Madan}
\begin{itemize}[leftmargin=*]
\item CF log-return Heston : \codeid{heston_log_return_cf} dans \filepath{app/model/volatility_models/heston/cf.py}.
\item FFT Carr--Madan : \codeid{carr_madan_fft_call_prices} dans \filepath{app/model/volatility_models/common/fft.py}.
\item Calibration : \codeid{HestonFFTCalibrator} dans \filepath{app/model/volatility_models/heston/calibrator_fft.py}.
\end{itemize}

\subsection*{Extrait de code : Carr--Madan FFT (noyau de calcul)}
Dans \filepath{app/model/volatility_models/common/fft.py} :

\begin{lstlisting}[caption={Carr--Madan FFT (extrait) : \filepath{app/model/volatility_models/common/fft.py}, \codeid{carr_madan_fft_call_prices}.},label={lst:carr-madan}]
u = eta * np.arange(n, dtype=float)
iu = 1j * u

shifted = u - 1j * (alpha + 1.0)
phi = cf_log_return(shifted, float(T))

disc = np.exp(-float(r) * float(T))
denom = (alpha * alpha + alpha - u * u) + 1j * (2.0 * alpha + 1.0) * u
denom = np.where(np.abs(denom) < 1e-16, 1e-16, denom)
psi = disc * phi / denom

lam = 2.0 * np.pi / (n * eta)
b = 0.5 * n * lam
k = -b + lam * np.arange(n, dtype=float)

w = _simpson_weights(n)
x = np.exp(1j * u * b) * psi * w * eta
y = np.fft.fft(x).real
\end{lstlisting}

\paragraph{Lecture «formule $\to$ code».}
\begin{itemize}[leftmargin=*]
\item \codeid{shifted = u - i(\alpha+1)} correspond au décalage de la transformée (atténuation Carr--Madan).
\item \codeid{psi = disc * phi / denom} est la transformée de la quantité atténuée (avec le dénominateur analytique).
\item \codeid{_simpson_weights} applique une quadrature de Simpson (stabilise l'intégrale discrétisée).
\item \codeid{np.fft.fft} calcule simultanément une grille de prix sur des log-strikes \codeid{k}.
\end{itemize}

\section{F. Exemple end-to-end}
Dans l'onglet «Calibration avancée» (\filepath{app/vue/tabs/tab_advanced_calibration.py}) :
\begin{enumerate}[leftmargin=*]
\item On sélectionne «Heston via FFT» (\codeid{model_key="heston_fft"}).
\item Le contrôleur \codeid{run_advanced_surface_calibration} instancie \codeid{HestonFFTCalibrator} et lance \codeid{least_squares}.
\item La sortie comprend : \codeid{params} (kappa/theta/sigma/rho/v0), \codeid{iv_model}, \codeid{iv_error}, \codeid{metrics}.
\item Un clic «Envoyer IV modèle vers Options» (dans \filepath{app/vue/tabs/tab_advanced_calibration.py}) pousse la surface dans \codeid{st.session_state["calib_model_surface_df"]} ; Options la lit via \filepath{app/vue/components/options/router.py}.
\end{enumerate}

\section{G. Limites \& extensions crédibles}
\begin{itemize}[leftmargin=*]
\item Le moteur FFT ici pricer des \textbf{calls européens}. La couverture d'options américaines sous Heston n'est pas implémentée dans ces calibrateurs (\NonImplem).
\item Extension : ajouter une structure de poids bid-ask et des contraintes d'arbitrage sur la surface.
\item Extension : calibration jointe avec taux (courbe) et dividendes : actuellement \codeid{get_q} renvoie 0.0 (placeholder) dans \filepath{app/model/yieldcurve/rates_utils.py}.
\end{itemize}

\end{comment}
\begin{jurybox}
Heston est le premier modèle «au-delà de BS» réellement codé : CF + intégrale (V1) et surtout FFT Carr--Madan (V2). La calibration ajuste des paramètres interprétables et applique une pénalité de type Feller pour éviter des variances pathologiques.
\end{jurybox}

\chapter{Jump diffusion : Merton et Bates (CF/FFT, calibration)}

\section{A. Problème financier (objectif, inputs, outputs)}
Les sauts modélisent des mouvements discontinus (earnings, news) qui génèrent des smiles plus marqués que Heston seul. Objectif : reproduire une surface d'IV avec un modèle comportant une composante de saut.

Deux modèles présents :
\begin{itemize}[leftmargin=*]
\item \textbf{Merton (1976)} : diffusion log-normale + sauts log-normaux (Poisson).
\item \textbf{Bates (1996)} : Heston + sauts Merton (vol stochastique + jumps).
\end{itemize}

\section{B. Théorie (définitions, hypothèses)}
\subsection*{Merton (log-returns avec sauts)}
Sous $\Qmeas$, on peut écrire le log-return $X_T=\ln(S_T/S_0)$ comme
\[
X_T = \Big(r-q-\tfrac12\sigma^2-\lambda \kappa_J\Big)T + \sigma W_T + \sum_{i=1}^{N_T} Y_i,
\]
où $N_T\sim\text{Poisson}(\lambda T)$, $Y_i\sim \Normal(\mu_J,\sigma_J^2)$ i.i.d., et
\[
\kappa_J = \E[e^{Y}]-1 = e^{\mu_J+\tfrac12\sigma_J^2}-1
\]
est l'ajustement de drift garantissant $e^{-(r-q)T}S_T$ martingale.

\subsection*{Bates}
Bates multiplie la composante Heston (vol stochastique) par une composante de sauts indépendante conditionnellement : le log-return admet une CF «produit».

\qa
{Capturer smile/skew extrêmes et événements via une composante de saut ; proposer une calibration alternative à Heston.}
{Sauts log-normaux, intensité constante $\lambda$, indépendance structurelle, pricing européen via CF/FFT.}
{Les sauts vrais peuvent être asymétriques ou dépendre du régime ; calibration instable si surface pauvre ; pas d'arbitrage non garanti.}
{Applications : pricing européen via fonction caractéristique et FFT, calibration d'une composante de saut $(\lambda,\mu_J,\sigma_J)$, et extension naturelle vers Bates (vol stochastique + sauts).}

\subsection*{Interprétation des paramètres de saut (et impact sur la surface)}
Le triplet $(\lambda,\mu_J,\sigma_J)$ contrôle des aspects distincts :
\begin{itemize}[leftmargin=*]
\item $\lambda$ (intensité) : fréquence moyenne des sauts ; augmente la probabilité d'événements rares.
\item $\mu_J$ : taille moyenne (en log) d'un saut ; si $\mu_J<0$, les sauts sont en moyenne baissiers (crash risk).
\item $\sigma_J$ : dispersion des tailles de saut ; contrôle l'épaisseur des queues.
\end{itemize}
Effet surface : des sauts négatifs fréquents et/ou dispersés rendent les puts OTM plus chers $\Rightarrow$ smile/skew plus prononcé qu'en diffusion pure.

\subsection*{Pourquoi la correction $-\lambda\kappa_J$ est nécessaire}
Sans correction, ajouter des sauts augmenterait mécaniquement la croissance moyenne de $S_t$ (car $\E[e^Y]>1$).
La condition «pas d'arbitrage» sous $\Qmeas$ impose que $e^{-(r-q)t}S_t$ soit une martingale, ce qui force l'ajustement
\[
\kappa_J=\E[e^Y]-1,\qquad \text{drift}\gets (r-q)-\lambda\kappa_J.
\]
Cette correction est un point non négociable : elle garantit que le forward reste $F_{0,T}=S_0e^{(r-q)T}$ malgré les sauts.

\section{C. Dérivations (ultra détaillées)}
\subsection*{CF de Merton}
Par indépendance (diffusion + Poisson), la fonction caractéristique du log-return est :
\[
\phi(u)=\exp\Big(iu\cdot \text{drift} -\tfrac12\sigma^2u^2T + \lambda T\big(\E[e^{iuY}]-1\big)\Big),
\]
et comme $Y\sim\Normal(\mu_J,\sigma_J^2)$,
\[
\E[e^{iuY}] = \exp\Big(iu\mu_J - \tfrac12\sigma_J^2u^2\Big).
\]
Donc
\[
\phi(u)=\exp\Big(iu\cdot \text{drift} -\tfrac12\sigma^2u^2T + \lambda T\big(e^{iu\mu_J - \frac12\sigma_J^2u^2}-1\big)\Big).
\]

\subsection*{Bates : produit}
Pour Bates, on écrit
\[
\phi_{\text{Bates}}(u)=\phi_{\text{Heston}}(u;\kappa,\theta,\sigma,\rho,v_0,\delta_{\text{drift}})\cdot \phi_{\text{Jumps}}(u;\lambda,\mu_J,\sigma_J),
\]
où le terme d'ajustement de drift vaut $\delta_{\text{drift}}=-\lambda\kappa_J$ pour préserver la martingale au niveau du drift (même mécanisme que Merton).

\subsection*{Merton : représentation en mélange (somme de Black--Scholes)}
Conditionnellement au nombre de sauts $N_T=n$, le log-return est normal :
\[
X_T\mid N_T=n \sim \Normal\Big(\big(r-q-\tfrac12\sigma^2-\lambda\kappa_J\big)T+n\mu_J,\ \sigma^2T+n\sigma_J^2\Big).
\]
Donc $S_T\mid N_T=n$ est log-normal, et le prix d'un call peut s'écrire comme une somme (mélange) :
\[
C=\sum_{n=0}^{\infty}\Prob(N_T=n)\,C^{(n)},
\quad \Prob(N_T=n)=e^{-\lambda T}\frac{(\lambda T)^n}{n!},
\]
où $C^{(n)}=e^{-rT}\E[(S_T-K)^+\mid N_T=n]$ admet une forme fermée (log-normale).
Cette décomposition est très pédagogique : \textbf{un modèle à sauts est un mélange de scénarios «$n$ sauts»}, chacun équivalent à une diffusion log-normale avec variance augmentée.
Numériquement, elle donne aussi une alternative à la FFT (tronquer la somme pour des $\lambda T$ modérés).

\begin{comment}
\section{D. Numérique / algorithmes (stabilité, complexité, choix)}
\begin{itemize}[leftmargin=*]
\item Pricing : Carr--Madan FFT identique au chapitre Heston (mêmes hyperparamètres \codeid{FFTConfig}).
\item Calibration : least squares sur prix (reconstruits depuis IV marché via BS) avec normalisation par \codeid{scale}.
\item Multi-start : plusieurs initialisations ; pour Bates, le code combine heuristique + aléatoire.
\end{itemize}

\end{comment}
\begin{comment}
\section{E. Implémentation dans le repo (fichiers, fonctions, pipeline, paramètres)}
\subsection*{Extrait de code : Bates CF (drift adjust + produit)}
Dans \filepath{app/model/volatility_models/jump_diffusion/cf.py} :

\begin{lstlisting}[caption={Bates : CF du log-return (extrait) : \filepath{app/model/volatility_models/jump_diffusion/cf.py}, \codeid{bates_log_return_cf}.},label={lst:bates-cf}]
kappa_j = np.exp(muj + 0.5 * sigj * sigj) - 1.0
drift_adjust = -lam * kappa_j

jump_cf = np.exp(1j * u * muj - 0.5 * sigj * sigj * u * u)
jump_part = np.exp(lam * T * (jump_cf - 1.0))

heston_part = heston_log_return_cf(
u, T, r=r, q=q, kappa=kappa, theta=theta, sigma=sigma, rho=rho, v0=v0,
drift_adjust=drift_adjust,
)
return heston_part * jump_part
\end{lstlisting}

\paragraph{Lecture financière.}
\begin{itemize}[leftmargin=*]
\item \codeid{kappa_j} est $\kappa_J=\E[e^{Y}]-1$.
\item \codeid{drift_adjust=-lamkappa_j} implémente la correction de drift (martingale).
\item \codeid{jump_part} est $\exp(\lambda T(\E[e^{iuY}]-1))$ (Poisson composé).
\item \codeid{heston_part} est la CF Heston ; le produit donne la CF Bates.
\end{itemize}

\subsection*{Calibration et exposition à l'UI}
\begin{itemize}[leftmargin=*]
\item Les modèles «Jump Diffusion (Merton) via FFT» et «Bates via FFT» sont exposés par \codeid{CalibrationController.get_advanced_models} (\filepath{app/controller/calibration_controller.py}).
\item Le runner unifié \codeid{run_advanced_surface_calibration} instancie \codeid{MertonJumpDiffusionCalibrator} et \codeid{BatesCalibrator}.
\end{itemize}

\section{F. Exemple end-to-end}
\begin{enumerate}[leftmargin=*]
\item Dans \filepath{app/vue/tabs/tab_advanced_calibration.py}, l'utilisateur choisit «Bates».
\item La surface de marché est convertie en \codeid{SurfaceGrid}.
\item Le calibrateur FFT produit une grille de prix call via Carr--Madan, interpole aux strikes voulus, calcule le résidu normalisé, et optimise.
\item Le résultat est converti en surface d'IV via \codeid{implied_vol_grid} (\filepath{app/model/calibration/implied_vol.py}).
\end{enumerate}

\section{G. Limites \& extensions crédibles}
\begin{itemize}[leftmargin=*]
\item Les paramètres de sauts peuvent être peu identifiables à partir d'une surface limitée en maturités ; extension : ajouter des contraintes ou un prior.
\item \NonImplem\ : pas de calibration conjointe à données «event» (earnings) ni de modèle d'intensité dépendant du temps.
\end{itemize}

\end{comment}
\begin{jurybox}
À retenir : une composante de saut de type Poisson composé se traduit par un terme $\lambda T(\E[e^{iuY}]-1)$ dans l'exposant de la fonction caractéristique. Cette structure se combine naturellement avec Heston (Bates) et rend le pricing européen efficace via FFT, puis la calibration se formule en moindres carrés sur des prix reconstruits depuis l'IV observée.
\end{jurybox}

\chapter{SABR : dynamique, formule de Hagan, calibration}

\section{A. Problème financier (objectif, inputs, outputs)}
SABR est un modèle de référence (rates/FX) pour reproduire une courbe de smile en fonction du strike. On l'utilise souvent comme \emph{générateur de surface d'IV} via une approximation analytique (Hagan) et on calibre \emph{par maturité}.

\textbf{Inputs} : $(S_0,r,q)$, grille $(m,T)$, paramètres SABR $(\alpha,\beta,\rho,\nu)$ (avec $\beta$ souvent fixé).
\textbf{Outputs} : surface d'IV modèle et paramètres calibrés (souvent $(\alpha,\rho,\nu)$ par maturité si $\beta$ est fixé).

\section{B. Théorie (définitions, hypothèses)}
\subsection*{Dynamique SABR (standard)}
Sur un forward $F_t$ :
\[
\dd F_t = \alpha_t F_t^\beta,\dd W_t,\qquad
\dd \alpha_t = \nu \alpha_t,\dd Z_t,\qquad
\Corr(\dd W_t,\dd Z_t)=\rho.
\]
Le but est de produire une volatilité implicite Black (sur forwards) approchée analytiquement.

\subsection*{Approximation de Hagan (2002) : idée}
On obtient une approximation fermée de $\sigma_{\text{Black}}(F,K,T)$ via un développement asymptotique en $T$ et en $\ln(F/K)$.

\qa
{Produire rapidement une surface d'IV (sans FFT) et calibrer de façon efficace (par maturité).}
{Validité de l'approximation Hagan (régime asymptotique), calibrage indépendant par maturité, $\beta$ fixé ou choisi.}
{Approximation moins fiable pour maturités longues ou strikes extrêmes ; ne garantit pas l'absence d'arbitrage ; calibration indépendante peut produire des courbes non lisses en $T$.}
{Applications : génération rapide de smiles par maturité, calibration en front-office, et construction de surfaces lisses (avec régularisation en $T$ si nécessaire).}

\subsection*{Interprétation des paramètres SABR}
\begin{itemize}[leftmargin=*]
\item $\beta\in[0,1]$ : élasticité (CEV). $\beta=1$ correspond à une dynamique log-normale, $\beta=0$ à une dynamique normale (en première approximation).
\item $\alpha$ : niveau de volatilité (à $t=0$), dimensionné pour que $\alpha F^\beta$ ait la dimension d'une volatilité sur $F$.
\item $\nu$ : \emph{vol-of-vol} (amplitude des fluctuations de $\alpha_t$) ; augmente la courbure du smile.
\item $\rho$ : corrélation entre mouvement de $F$ et mouvement de $\alpha$ ; contrôle l'asymétrie du smile (skew).
\end{itemize}
En pratique, on fixe souvent $\beta$ (par convention de marché) et on ajuste $(\alpha,\rho,\nu)$.

\subsection*{Pourquoi SABR se formule sur le forward}
Dans les marchés de taux/FX, les options sont souvent cotées en Black sur un forward $F_{t,T}$ (ou un swap rate).
Travailler en forward permet d'isoler l'effet de taux et de modéliser directement la dynamique de l'actif sous-jacent «de marché» pour l'option.

\section{C. Dérivations (ultra détaillées)}
\subsection*{Notations}
On travaille en forward $F=S_0 e^{(r-q)T}$ (taux et dividendes constants) et on note $k=\ln(F/K)$ le log-moneyness.

\subsection*{Formule Hagan : deux cas}
La formule de Hagan se présente naturellement en deux régimes :
\begin{itemize}[leftmargin=*]
\item \textbf{ATM} : $\ln(F/K)\approx 0$, formule simplifiée.
\item \textbf{général} : via
\[
z=\frac{\nu}{\alpha}(FK)^{\frac{1-\beta}{2}}\ln\frac{F}{K},
\quad
x(z)=\ln\frac{\sqrt{1-2\rho z+z^2}+z-\rho}{1-\rho},
\]
puis un facteur $\frac{z}{x(z)}$.
\end{itemize}
Le point clé est le ratio $z/x(z)$, qui corrige l'asymétrie induite par la corrélation $\rho$.

\subsection*{Limites et cas particuliers (contrôles de cohérence)}
Ces cas servent de tests «mentaux» sur la formule :
\begin{itemize}[leftmargin=*]
\item \textbf{$\nu\to 0$} : la volatilité devient (presque) déterministe et on retrouve une dynamique de type CEV ; le smile doit s'aplatir (moins de courbure).
\item \textbf{$\rho=0$} : smile plus symétrique autour de l'ATM (pas de skew induit par corrélation).
\item \textbf{$\beta=1$} : on se rapproche d'un cadre log-normal (Black/BS) ; les expressions se simplifient.
\item \textbf{ATM ($K\to F$)} : la formule générale contient des termes en $\ln(F/K)$ ; il faut une expression limite bien définie (d'où l'usage explicite de la limite ATM).
\end{itemize}

\subsection*{Pourquoi le terme $z/x(z)$ apparaît}
Le terme
\[
z=\frac{\nu}{\alpha}(FK)^{\frac{1-\beta}{2}}\ln\frac{F}{K}
\]
est une variable adimensionnée mesurant à la fois l'écart à l'ATM ($\ln(F/K)$) et l'amplitude de vol-of-vol ($\nu/\alpha$).
La fonction $x(z)$ vient d'une intégration qui «corrige» l'effet de corrélation $\rho$ ; le ratio $z/x(z)$ est un facteur d'asymétrie qui tend vers 1 quand $z\to 0$ (cohérence ATM).

\begin{comment}
\section{D. Numérique / algorithmes (stabilité, complexité, choix)}
\subsection*{Calibration par maturité}
Le calibrateur \codeid{SABRAnalyticCalibrator} (fichier \filepath{app/model/volatility_models/sabr/calibrator.py}) :
\begin{itemize}[leftmargin=*]
\item fixe $\beta$ (par défaut 0.5, override possible via contraintes),
\item pour chaque $T_i$ de la grille, ajuste $(\alpha_i,\rho_i,\nu_i)$ par least squares sur les IV (pas sur les prix),
\item utilise multi-start (x0 heuristique + tirages aléatoires).
\end{itemize}

\end{comment}
\begin{comment}
\section{E. Implémentation dans le repo (fichiers, fonctions, pipeline, paramètres)}
\subsection*{Points d'entrée}
\begin{itemize}[leftmargin=*]
\item Exposé à l'UI via \codeid{CalibrationController.get_advanced_models} : clé \codeid{"sabr"} (\filepath{app/controller/calibration_controller.py}).
\item Runner : \codeid{CalibrationController.run_advanced_surface_calibration} instancie \codeid{SABRAnalyticCalibrator}.
\end{itemize}

\subsection*{Extrait de code : Hagan (cas général)}
Dans \filepath{app/model/volatility_models/sabr/analytic.py} (extrait) :

\begin{lstlisting}[caption={Hagan SABR (extrait) : \filepath{app/model/volatility_models/sabr/analytic.py}, \codeid{hagan_black_iv}.},label={lst:hagan}]
fk_beta = (F * K) ** (0.5 * one_minus_beta)
z = (nu / alpha) * fk_beta * log_fk
sqrt_term = np.sqrt(max(1e-18, 1.0 - 2.0 * rho * z + z * z))
xz_num = sqrt_term + z - rho
xz_den = 1.0 - rho
xz = np.log(max(1e-18, xz_num / max(1e-18, xz_den)))

log_corr = 1.0 + (one_minus_beta * one_minus_beta / 24.0) * (log_fk * log_fk) + (
(one_minus_beta4) / 1920.0
) * (log_fk4)

z_over_xz = float(z / xz) if abs(xz) > 1e-14 else 1.0
iv = (alpha / (fk_beta * log_corr)) * z_over_xz * (1.0 + (term1 + term2 + term3) * T)
\end{lstlisting}

\paragraph{Lecture «math $\to$ code».}
\begin{itemize}[leftmargin=*]
\item \codeid{z} et \codeid{xz} implémentent la transformation non linéaire (gestion de $\rho$).
\item \codeid{log_corr} est la correction en $\ln(F/K)$.
\item \codeid{z_over_xz} est le facteur stabilisant, avec un fallback si $x(z)\approx 0$ (ATM).
\end{itemize}

\section{F. Exemple end-to-end}
Dans l'onglet «Calibration avancée», sélectionner SABR :
\begin{enumerate}[leftmargin=*]
\item la surface marché est groupée implicitement par maturité (via masque sur la ligne $T_i$),
\item pour chaque maturité, on minimise $\sum_k (\sigma_{\text{Hagan}}(K_k,T_i)-\sigma^{\text{mkt}}_k)^2$,
\item on reconstruit une surface complète \codeid{iv_model_full} via \codeid{SABRAnalyticModel.implied_vol_surface}.
\end{enumerate}

\section{G. Limites \& extensions crédibles}
\begin{itemize}[leftmargin=*]
\item \NonImplem\ : Nelson--Siegel--Svensson n'a pas de lien direct ici ; côté SABR, aucune variante (shifted SABR, SABR arbitrage-free) n'est repérée.
\item Extension : régulariser les courbes $(\alpha(T),\rho(T),\nu(T))$ via un lissage (Kalman/penalty). Un lisseur générique existe : \codeid{kalman_smooth} dans \filepath{app/controller/calibration_controller.py} (random-walk smoother dans \filepath{app/model/volatility_models/kalman/kalman_filter.py}).
\end{itemize}

\end{comment}
\begin{jurybox}
SABR via l'approximation de Hagan est une voie \textbf{très rapide} pour obtenir un smile plausible : calibration par maturité, multi-start si besoin, et paramètres interprétables. Contrepartie : c'est une approximation asymptotique (moins fiable sur maturités longues/strikes extrêmes).
\end{jurybox}

\chapter{Rough \& Volterra : rHeston (approx. markovienne), rBergomi (MC+surrogate), Volterra SDE (MC proxy)}

\section{A. Problème financier (objectif, inputs, outputs)}
Certaines surfaces d'IV exhibent une \textbf{rugosité} (décroissance lente de la corrélation de variance) incompatible avec Heston classique. Les modèles rough/Volterra visent à :
\begin{itemize}[leftmargin=*]
\item reproduire des dynamiques de volatilité à \emph{mémoire} (non markoviennes),
\item obtenir un meilleur fit sur l'échelle courte,
\item conserver une interprétation : paramètre de rugosité $H\in(0,1/2)$.
\end{itemize}

Trois familles de constructions illustrent ces idées :
\begin{itemize}[leftmargin=*]
\item rHeston via approximation markovienne + FFT,
\item rBergomi via Monte Carlo + surrogate,
\item un simulateur Volterra générique + calibration par recherche sur design.
\end{itemize}

\section{B. Théorie (définitions, hypothèses)}
\subsection*{Rugosité et noyau fractionnaire}
Un point de départ commun est le noyau fractionnaire (type Riemann--Liouville) :
\[
K(t)=\frac{t^{H-\frac12}}{\Gamma(H+\frac12)},\qquad H\in(0,1/2),
\]
qui induit une mémoire longue dans la variance.

\subsection*{Approximation markovienne (idée)}
On approxime $K(t)$ par une somme d'exponentielles :
\[
K(t)\approx \sum_{i=1}^{n} w_i e^{-x_i t},
\]
ce qui permet de représenter un système non markovien via un état augmenté (Markov) de dimension $n$.

\subsection*{rBergomi (idée)}
On construit une volatilité log-normale pilotée par une fBM (fractional Brownian motion) de paramètre $H$ (ici via une approximation Riemann--Liouville), puis on simule $S_t$.

\qa
{Modéliser le smile court terme et la mémoire de volatilité ; fournir des calibrations «research-grade» et des surfaces modèles.}
{Roughness $H\in(0,1/2)$, approximations numériques (convolutions $O(n^2)$ ou état augmenté), pricing européen par FFT/MC.}
{Coût computationnel élevé, objectives bruitées, identifiabilité faible ; risque d'overfitting ; absence d'arbitrage non garantie.}
{Applications : reproduire le skew court terme (lois de puissance), comparer des familles rough (rHeston, rBergomi, Volterra), et analyser le compromis précision/coût via markovianisation et surrogates.}

\subsection*{Qu'est-ce que «rough» ? (intuition statistique)}
Dire que la volatilité est «rough» signifie que ses trajectoires sont plus irrégulières qu'un brownien standard : localement, elles ressemblent à une fBM de paramètre de Hurst $H<1/2$.
Conséquences qualitatives :
\begin{itemize}[leftmargin=*]
\item la corrélation des incréments de volatilité décroît lentement et de façon non exponentielle ;
\item le skew implicite court terme suit souvent une loi de puissance en maturité (signature empirique) ;
\item les modèles deviennent \textbf{non markoviens} : l'état présent ne suffit plus à résumer le futur.
\end{itemize}
Attention : «rough» n'est pas forcément «long memory» au sens économétrique ; il s'agit d'une régularité locale (höldérienne).

\subsection*{Volterra : la bonne classe pour formaliser la mémoire}
Un processus de type Volterra s'écrit (schématiquement)
\[
Y_t = \int_0^t K(t-s)\,\dd B_s,
\]
où $K$ est un noyau. Si $K(t)\propto t^{H-1/2}$, on obtient une régularité contrôlée par $H$.
Les modèles rough sont souvent des modèles de volatilité construits à partir d'un tel $Y_t$ (puis exponentiation, reversion, etc.).

\section{C. Dérivations (ultra détaillées)}
\subsection*{(1) Approximation markovienne du noyau fractionnaire}
Le point technique central est d'approximer le noyau de loi de puissance par une somme d'exponentielles :
\[
K(t)=\frac{t^{H-\frac12}}{\Gamma(H+\frac12)} \approx \sum w_i e^{-x_i t},
\]
ce qui transforme une convolution (mémoire) en un état augmenté (Markov) de dimension $n$.

\paragraph{Idée via représentation de Laplace.}
Un noyau complètement monotone admet souvent une représentation de la forme
\[
K(t)=\int_0^{\infty} e^{-xt}\,\mu(\dd x),
\]
pour une mesure positive $\mu$. En approchant cette intégrale par une quadrature
\[
\int_0^\infty e^{-xt}\,\mu(\dd x)\approx \sum_{i=1}^n w_i e^{-x_i t},
\]
on obtient exactement la somme d'exponentielles. Les grilles géométriques (points $x_i$ espacés multiplicativement) sont naturelles car le noyau a une structure multi-échelles.

\subsection*{(2) rHeston : CF via état augmenté}
Une fois $K$ markovianisé, la variance (rough) devient fonction d'un vecteur de facteurs $(Y^{(1)},\ldots,Y^{(n)})$ qui résument la mémoire.
Dans les modèles affines (famille Heston), la fonction caractéristique garde une structure exponentielle-affine, mais les équations de Riccati deviennent un système de dimension $n$ (au lieu de 1 en Heston classique).
On retient :
\begin{itemize}[leftmargin=*]
\item \textbf{gain} : on récupère un état markovien donc des outils de pricing (CF/FFT) ;
\item \textbf{coût} : on remplace une ODE scalaire (Heston) par un système d'ODE (ou une discrétisation) en dimension $n$.
\end{itemize}

\subsection*{(3) rBergomi : construction d'une fBM Riemann--Liouville}
Une représentation Riemann--Liouville (RL) d'une fBM est :
\[
W_t^H \approx \sqrt{2H}\sum_{j<t} (t-t_j)^{H-\frac12}\Delta B_j,
\]
c'est une convolution discrète. Naïvement, le coût est $O(n_{\text{steps}}^2)$ (double boucle).
On peut accélérer via des approximations de noyau (sum-of-exponentials) ou via des méthodes FFT-convolution (selon la grille).

\subsection*{(4) Volterra SDE générique}
Une Volterra SDE pour la variance peut s'écrire, schématiquement,
\[
V_t=v_0+\int_0^t K(t-s)b(V_s)\,\dd s+\int_0^t K(t-s)\sigma(V_s)\,\dd B_s,
\]
où $K$ encode la mémoire. Sur une grille, cela devient une convolution discrète : pour $t_i$,
\[
V_{t_i}=v_0+\sum_{m<i} K(t_i-t_m),\Delta Z_m,
\]
où $\Delta Z_m$ combine drift (mean-reversion) et bruit ($\sqrt{V}\Delta B$). L'essentiel : l'état dépend de tout le passé via $K$.

\subsection*{Pourquoi une somme d'exponentielles rend le modèle markovien}
Si $K(t)\approx \sum_{i=1}^n w_ie^{-x_it}$, alors une convolution de type
\[
Y_t=\int_0^t K(t-s)\,\dd B_s
\]
se réécrit approximativement comme une somme de facteurs
\[
Y_t\approx \sum_{i=1}^n w_i Y_t^{(i)},\qquad Y_t^{(i)}=\int_0^t e^{-x_i(t-s)}\,\dd B_s.
\]
Chaque $Y^{(i)}$ vérifie une SDE \emph{markovienne} :
\[
\dd Y_t^{(i)}=-x_i Y_t^{(i)}\,\dd t+\dd B_t.
\]
Donc le vecteur $(Y^{(1)},\ldots,Y^{(n)})$ est un état markovien de dimension $n$ qui approxime la mémoire du noyau de puissance.
C'est exactement le mécanisme derrière la «markovianisation» de rHeston.

\subsection*{Complexité : pourquoi ces modèles sont coûteux}
Sans markovianisation, la convolution discrète de type
\[
W_{t_i}^H \approx \sum_{j<i} K(t_i-t_j)\Delta B_j
\]
coûte $O(n_{\text{steps}}^2)$ (double boucle).
La somme d'exponentielles réduit le coût à $O(n_{\text{steps}}\cdot n)$, au prix d'une approximation contrôlée par $n$.

\begin{comment}
\section{D. Numérique / algorithmes (stabilité, complexité, choix)}
\subsection*{rHeston : FFT + CF}
\begin{itemize}[leftmargin=*]
\item On obtient une CF approximative, puis on pricer via Carr--Madan FFT (comme Heston).
\item Calibration : least squares sur prix (comme Merton/Heston FFT) dans \filepath{app/model/volatility_models/rheston/calibrator_fft.py}.
\end{itemize}

\subsection*{rBergomi : surrogate}
\begin{itemize}[leftmargin=*]
\item On échantillonne un plan d'expériences (\codeid{latin_hypercube_samples}) dans l'espace des paramètres.
\item On évalue une perte (moyenne d'erreur de prix) par Monte Carlo.
\item On ajuste un surrogate RBF (\codeid{RBFInterpolator}) et on explore rapidement des candidats.
\item On réévalue par Monte Carlo au meilleur candidat surrogate.
\end{itemize}
Implémentation : \filepath{app/model/volatility_models/rbergomi/calibrator_mc_surrogate.py}.

\subsection*{Volterra : MC proxy search}
Le calibrateur Volterra (\filepath{app/model/volatility_models/volterra/calibrator_mc.py}) ne fait pas d'optimisation continue : il choisit le meilleur point sur un design LHS (avec un point d'ancrage).

\end{comment}
\begin{comment}
\section{E. Implémentation dans le repo (fichiers, fonctions, pipeline, paramètres)}
\subsection*{rHeston : noyau markovien et CF}
\begin{itemize}[leftmargin=*]
\item Approximation noyau : \codeid{fractional_kernel_markovian_approx} (\filepath{app/model/volatility_models/rheston/markovian_kernel.py}).
\item CF markovienne : \codeid{rheston_log_return_cf_markovian} (\filepath{app/model/volatility_models/rheston/cf_markovian.py}).
\item Calibrateur : \codeid{RHestonFFTMarkovianCalibrator} (\filepath{app/model/volatility_models/rheston/calibrator_fft.py}).
\end{itemize}

\subsection*{rBergomi : simulation et calibration}
\begin{itemize}[leftmargin=*]
\item Simulation : \codeid{simulate_rbergomi_paths} (\filepath{app/model/volatility_models/rbergomi/simulator.py}).
\item Pricing MC call : \codeid{price_call_grid_mc} (\filepath{app/model/volatility_models/rbergomi/pricing_mc.py}).
\item Calibration surrogate : \codeid{RBergomiMCSurrogateCalibrator} (\filepath{app/model/volatility_models/rbergomi/calibrator_mc_surrogate.py}).
\end{itemize}

\subsection*{Volterra : kernel + simulateur}
\begin{itemize}[leftmargin=*]
\item Noyaux : \codeid{fractional_kernel}, \codeid{exponential_kernel} (\filepath{app/model/volatility_models/volterra/kernels.py}).
\item Simulation : \codeid{simulate_volterra_sv_paths} (\filepath{app/model/volatility_models/volterra/simulator.py}).
\item Calibration : \codeid{VolterraSDECalibrator} (\filepath{app/model/volatility_models/volterra/calibrator_mc.py}).
\end{itemize}

\section{F. Exemple end-to-end}
\begin{example}[rBergomi (MC+surrogate) dans le runner unifié]
\begin{enumerate}[leftmargin=*]
\item L'utilisateur choisit \codeid{"rbergomi"} dans \filepath{app/vue/tabs/tab_advanced_calibration.py}.
\item Le contrôleur instancie \codeid{RBergomiMCSurrogateCalibrator}.
\item Le calibrateur :
\begin{itemize}[leftmargin=*]
\item transforme l'IV marché en prix via BS (\codeid{bs_call_price}),
\item génère un design LHS de paramètres,
\item simule des paths rough bergomi (O($n_{\text{steps}}^2$)) et pricer des calls,
\item ajuste un surrogate RBF et sélectionne un candidat,
\item reconstruit une surface d'IV via inversion (\codeid{implied_vol_grid}).
\end{itemize}
\end{enumerate}
\end{example}

\section{G. Limites \& extensions crédibles}
\begin{itemize}[leftmargin=*]
\item Le dépôt signale ces modèles comme «expensive» mais ne met pas en place de parallélisation ; extension : vectorisation/numba ou exécution parallèle.
\item \NonImplem\ : pas de stratégie «production» de calibration rough (avec convergence statistique contrôlée, seeds multiples, CI).
\item Extension : approximation markovienne plus riche pour Volterra (somme d'exponentielles) afin de passer de $O(n^2)$ à $O(n)$.
\end{itemize}

\end{comment}
\begin{jurybox}
Rough/Volterra sont des modèles \textbf{coûteux} et souvent «research-grade» : (i) rHeston rend le rough plus \emph{calibrable} via une approximation markovienne et FFT, (ii) rBergomi contourne le coût MC via des surrogates, (iii) une Volterra SDE générique se simule mais la calibration nécessite souvent une recherche sur design (proxy).
\end{jurybox}

\chapter{Réseaux de neurones dans la calibration (surrogates pour Heston)}

\section{A. Problème financier (objectif, inputs, outputs)}
Calibrer Heston par optimisation non linéaire est coûteux (intégrales/FFT + inversion IV) et non convexe. Une alternative est d'apprendre une \emph{inversion directe} :
\[
\text{surface IV} ;\mapsto; (\kappa,\theta,\sigma,\rho,v_0),
\]
pour accélérer :
\begin{itemize}[leftmargin=*]
\item l'initialisation (meilleur $x_0$),
\item ou une calibration «en un coup» (surrogate).
\end{itemize}

\textbf{Input} : une matrice d'IV sur une grille fixe.
\textbf{Output} : un dictionnaire de paramètres Heston.

\section{B. Théorie (définitions, hypothèses)}
\subsection*{Régression supervisée}
On définit un dataset synthétique ${(X^{(n)},y^{(n)})}$ :
\begin{itemize}[leftmargin=*]
\item $X^{(n)}$ : surface d'IV générée par Heston pour un paramètre $y^{(n)}$,
\item $y^{(n)}=(\kappa,\theta,\sigma,\rho,v_0)$.
\end{itemize}
On entraîne un réseau $f_\omega$ en minimisant une MSE :
\[
\min_\omega \frac1N\sum_{n=1}^N |f_\omega(X^{(n)})-y^{(n)}|^2.
\]

\subsection*{Intuition}
Le réseau apprend un inverse approximatif du pipeline :
\[
(\kappa,\theta,\sigma,\rho,v_0)\to \text{prix}\to \text{IV}.
\]
Il ne remplace pas nécessairement une calibration «exacte», mais fournit une estimation rapide.

\qa
{Accélérer la calibration (ou fournir un point de départ robuste) en apprenant l'inversion surface $\to$ paramètres.}
{Dataset synthétique représentatif ; grille fixe (moneyness, maturités) ; stabilité de l'inversion IV.}
{Risque de biais (domain shift : marché réel vs synthétique) ; mauvaise généralisation hors domaine ; incertitude non quantifiée.}
{Application : estimation instantanée des paramètres (ou initialisation) à partir d'une surface, puis éventuel raffinement par optimisation classique.}

\subsection*{Pourquoi l'inversion «surface $\to$ paramètres» est difficile}
Le problème inverse est (souvent) \textbf{mal posé} :
\begin{itemize}[leftmargin=*]
\item \textbf{Non-unicité} : différentes valeurs de $(\kappa,\theta,\sigma,\rho,v_0)$ peuvent produire des surfaces très proches sur un domaine fini $(K,T)$.
\item \textbf{Sensibilités hétérogènes} : certains paramètres modifient surtout le skew, d'autres la term structure ; si la grille n'observe pas assez ces dimensions, l'estimation devient instable.
\item \textbf{Bruit de marché} : une surface issue de quotes bid/ask n'est pas une fonction lisse.
\end{itemize}
Un réseau de neurones ne «résout» pas ces problèmes : il apprend une inverse \emph{effective} sur un domaine de données donné.

\subsection*{Surrogate vs inversion directe : deux usages différents}
\begin{itemize}[leftmargin=*]
\item \textbf{Surrogate direct} : apprendre $(\theta)\mapsto \Sigma(\theta)$ pour accélérer le pricing durant l'optimisation.
\item \textbf{Inverse direct} (ici) : apprendre $\Sigma\mapsto \theta$ pour obtenir une estimation instantanée ou une initialisation.
\end{itemize}
Les deux peuvent être combinés : l'inverse donne un bon point de départ et le surrogate direct accélère les itérations de calibration.

\section{C. Dérivations (ultra détaillées)}
\subsection*{Normalisation des entrées}
Le réseau consomme une surface 2D ; on gère les valeurs manquantes (masque, interpolation simple, ou remplissage neutre), puis on normalise :
\[
X_{\text{norm}} = \frac{X-\mu}{\sigma+\varepsilon},
\]
où $(\mu,\sigma)$ sont estimés sur le dataset (standardisation).

\subsection*{Contraintes de sortie}
Pour respecter le domaine des paramètres :
\begin{itemize}[leftmargin=*]
\item paramètres positifs via \emph{softplus} : $\text{softplus}(x)=\ln(1+e^x)$,
\item corrélation $\rho\in(-1,1)$ via \emph{tanh}.
\end{itemize}
On applique donc une post-transformation \emph{différentiable} aux sorties du réseau, ce qui stabilise l'apprentissage et évite des paramètres non physiques.

\subsection*{Pourquoi la normalisation améliore l'apprentissage (lecture optimisation)}
Sans normalisation, les IV peuvent avoir des ordres de grandeur très différents selon $(K,T)$, ce qui détériore le conditionnement du problème et rend la descente de gradient instable (gradients très disparates).
La standardisation
\[
X_{\text{norm}}=\frac{X-\mu}{\sigma+\varepsilon}
\]
ramène les entrées à une échelle comparable et permet :
\begin{itemize}[leftmargin=*]
\item des pas de gradient plus homogènes ;
\item une meilleure utilisation de la capacité du réseau ;
\item une convergence plus rapide et plus stable.
\end{itemize}

\subsection*{Contraintes : projection différentiable sur un domaine admissible}
La post-transformation (softplus/tanh) est une forme de projection \emph{différentiable} :
elle empêche le réseau de prédire des paramètres physiquement absurdes (variance négative, corrélation hors $[-1,1]$).
On peut aller plus loin en ajoutant des pénalités (ex. Feller) dans la loss ; c'est une extension naturelle si l'on veut des sorties plus «physiques».

\begin{comment}
\section{D. Numérique / algorithmes (stabilité, complexité, choix)}
\begin{itemize}[leftmargin=*]
\item Dépendance : PyTorch est \emph{optionnel}. \codeid{TORCH_AVAILABLE} est testé, et sinon les fonctions renvoient une erreur contrôlée.
\item Dataset synthétique : coût élevé (pricing CF + inversion IV sur grille) ; le code limite par \codeid{max_attempts} et filtre les surfaces trop «NaN».
\item Sauvegarde : poids dans \filepath{app/model/calibration/weights/heston_surface_net.pt} (chemin \codeid{WEIGHTS_PATH}).
\end{itemize}

\end{comment}
\begin{comment}
\section{E. Implémentation dans le repo (fichiers, fonctions, pipeline, paramètres)}
\subsection*{Génération du dataset (vérifiée)}
Dans \filepath{app/model/calibration/heston_nn.py}, la fonction \codeid{_generate_synthetic_dataset} :
\begin{itemize}[leftmargin=*]
\item tire aléatoirement des paramètres dans des bornes codées en dur (ex. $\kappa\in[0.5,5]$, $\rho\in[-0.9,0.3]$),
\item calcule une grille de prix call par \codeid{call_price_cf} (\filepath{app/model/calibration/heston_pricer.py}),
\item reconstruit une grille d'IV via \codeid{implied_vol_grid} (\filepath{app/model/calibration/implied_vol.py}),
\item normalise et stocke $(X,y)$.
\end{itemize}

\subsection*{Extrait de code : contraintes de sortie (softplus/tanh)}
Dans \filepath{app/model/calibration/heston_nn.py} :

\begin{lstlisting}[caption={Sorties contraintes : \filepath{app/model/calibration/heston_nn.py}, \codeid{postprocess_tensor} (extrait).},label={lst:postprocess}]
def postprocess_tensor(raw: torch.Tensor) -> torch.Tensor:
kappa = torch.nn.functional.softplus(raw[..., 0]) + EPS
theta = torch.nn.functional.softplus(raw[..., 1]) + EPS
sigma = torch.nn.functional.softplus(raw[..., 2]) + EPS
rho = torch.tanh(raw[..., 3]) * 0.999
v0 = torch.nn.functional.softplus(raw[..., 4]) + EPS
return torch.stack([kappa, theta, sigma, rho, v0], dim=-1)
\end{lstlisting}

\paragraph{Pourquoi c'est important ?}
\begin{itemize}[leftmargin=*]
\item Sans cela, le réseau pourrait proposer $\kappa<0$ ou $v_0<0$, rendant le pricer instable.
\item La contrainte sur $\rho$ empêche des valeurs non physiques ($|\rho|\ge 1$) qui casseraient certaines racines carrées/CF.
\end{itemize}

\section{F. Exemple end-to-end}
\begin{enumerate}[leftmargin=*]
\item L'utilisateur entraîne les poids via \codeid{CalibrationController.train_heston_nn_weights} (\filepath{app/controller/calibration_controller.py}).
\item Pour une surface donnée, \codeid{predict_params} renvoie un dictionnaire \codeid{params}.
\item Le contrôleur peut ensuite reconstruire une surface modèle via \codeid{price_grid_from_params} puis \codeid{implied_vol_grid} (pipeline explicite dans \codeid{run_heston_nn_from_yahoo}).
\end{enumerate}

\section{G. Limites \& extensions crédibles}
\begin{itemize}[leftmargin=*]
\item Le réseau est entraîné sur un dataset \emph{synthétique} : extension cruciale = calibrer/valider sur surfaces réelles, et quantifier l'incertitude (ensembles, dropout).
\item \NonImplem\ : pas de mécanisme d'early stopping, pas de split train/val explicite dans \codeid{train_heston_surface_net}.
\end{itemize}

\end{comment}
\begin{jurybox}
Un NN inverseur (surface $\to$ paramètres) fournit une estimation rapide avec contraintes de domaine (softplus/tanh). C'est une brique d'accélération (initialisation ou surrogate inverse), mais un usage «production» exige validation hors synthétique et quantification d'incertitude.
\end{jurybox}

% =========================================================
% PARTIE III — RL Hedger
% =========================================================
\part{RL Hedger (théorie $\to$ design $\to$ pratique)}

\chapter{Couverture RL : formulation MDP, DQN, et aspects pratiques}

\section{A. Problème financier (objectif, inputs, outputs)}
\subsection*{Objectif}
Une position options a une exposition dynamique (delta, gamma, vega). On veut \textbf{réduire un risque} (ici : résidu de delta) en rebalançant une position de couverture (action) de façon \emph{discrète} et sous contraintes (coûts, slippage, limites d'inventaire).

\subsection*{Choix minimal d'un objectif (exemple pédagogique)}
Pour illustrer la démarche RL, on peut partir d'un objectif volontairement simple :
\begin{itemize}[leftmargin=*]
\item état = (exposition delta proxy, position action),
\item action = {hold, buy, sell, flatten},
\item reward = $-| \text{delta résiduel} |$ (éventuellement pénalisé par des coûts de transaction).
\end{itemize}

\textbf{Inputs} : définition d'état/action, simulateur (ou données) pour générer des transitions, et un critère de coût/risque.
\textbf{Outputs} : une politique (ou une valeur $Q$) et des recommandations d'ajustement de couverture.

\section{B. Théorie (définitions, hypothèses)}
\subsection*{(1) Hedging discret : PnL et coûts}
En discret, la PnL d'un hedge delta approximatif inclut :
\[
\Delta \Pi \approx \Delta V - \Delta_t \Delta S - \text{coûts},
\]
où $\Delta_t$ est la quantité d'action détenue. Les coûts (spread, commissions, slippage) rendent le hedging \emph{problème de contrôle} plutôt qu'égalité de réplication.

\subsection*{(2) MDP}
Un \emph{processus de décision markovien} (MDP) est $(\mathcal{S},\mathcal{A},P,R,\gamma)$ :
\begin{itemize}[leftmargin=*]
\item état $s\in\mathcal{S}$,
\item action $a\in\mathcal{A}$,
\item transition $P(s'|s,a)$,
\item récompense $R(s,a,s')$,
\item discount $\gamma\in(0,1)$.
\end{itemize}
On cherche une politique $\pi(a|s)$ maximisant $\E[\sum_t \gamma^t r_t]$.

\subsection*{(3) SARSA / Expected SARSA / Q-learning}
\textbf{Théorie générale (cadre standard RL).}
\begin{itemize}[leftmargin=*]
\item SARSA (on-policy) :
\[
Q(s,a)\leftarrow Q(s,a)+\alpha\big(r+\gamma Q(s',a')-Q(s,a)\big).
\]
\item Expected SARSA :
\[
Q(s,a)\leftarrow Q(s,a)+\alpha\big(r+\gamma \E_{a'\sim\pi(\cdot|s')}[Q(s',a')]-Q(s,a)\big).
\]
\item Q-learning (off-policy) :
\[
Q(s,a)\leftarrow Q(s,a)+\alpha\big(r+\gamma \max_{a'}Q(s',a')-Q(s,a)\big).
\]
\end{itemize}
Dans un DQN, la cible utilise typiquement le terme $\max_{a'} Q_{\text{target}}(s',a')$ (réseau cible) : c'est une forme de Q-learning approximé.

\subsection*{(4) Importance sampling}
En off-policy, on corrige le décalage entre politique comportement $\mu$ et politique cible $\pi$ via un ratio :
\[
\rho_t=\frac{\pi(a_t|s_t)}{\mu(a_t|s_t)},
\]
au prix d'une variance potentiellement élevée.

\qa
{Passer d'une couverture «réplication continue» (BS) à une couverture \emph{décisionnelle} en discret, potentiellement sous coûts et contraintes.}
{Markovianité approximative de l'état ; stationnarité ; reward représentatif du risque réel ; accès aux transitions (simulation).}
{Dépend fortement du choix d'état/reward ; risque d'overfitting à un simulateur ; sécurité/exécution réelle.}
{Applications : formulation MDP d'un hedge discret, apprentissage d'une stratégie sous coûts, et comparaison de variantes (Q-learning, Double DQN, distributional RL) selon le niveau de risque/cout voulu.}

\subsection*{De la réplication (théorie) au contrôle (pratique)}
En Black--Scholes \emph{idéal} (continu, sans coûts), on peut répliquer un payoff européen : le hedge est imposé (delta).
En pratique, trois écarts rendent la réplication impossible ou sous-optimale :
\begin{itemize}[leftmargin=*]
\item \textbf{Temps discret} : rebalancer toutes les $\Delta t$ introduit une erreur de discrétisation (hedging error) qui dépend de $\Gamma$ et de la volatilité.
\item \textbf{Coûts de transaction} : rebalancer plus souvent réduit l'erreur mais augmente les coûts ; on obtient un compromis.
\item \textbf{Contraintes} : inventaire limité, taille minimale d'ordre, interdiction de vendre à découvert, etc.
\end{itemize}
Le problème devient alors un \textbf{problème de contrôle stochastique} : choisir une action (rebalancement) pour minimiser un critère de risque.

\subsection*{Choix du reward : ce que l'on optimise vraiment}
Un reward de type $-|\\Delta_{\text{résiduel}}|$ est un proxy simple : il encourage une exposition proche de 0.
Mais d'autres critères plus proches d'une gestion de risque existent :
\begin{itemize}[leftmargin=*]
\item minimiser $\E[(\text{hedging error})^2]$ (mean-variance) ;
\item maximiser une utilité $\E[U(\text{PNL})]$ (aversion au risque) ;
\item pénaliser l'inventaire et les coûts (L1/L2) pour contrôler le turnover.
\end{itemize}
Le point pédagogique : \textbf{le RL n'est pas «magique»} ; la performance dépend directement de la qualité de l'environnement et du reward.

\section{C. Dérivations (ultra détaillées)}
\subsection*{Bellman optimalité}
Pour une politique optimale, la fonction $Q^\star$ vérifie :
\[
Q^\star(s,a)=\E\big[r+\gamma\max_{a'}Q^\star(s',a')\mid s,a\big].
\]
DQN approxime $Q^\star(s,a)$ par un réseau $Q_\omega(s,a)$.

\subsection*{Perte TD (DQN)}
On définit une cible (avec réseau cible $\omega^-$) :
\[
y = r + \gamma(1-d)\max_{a'} Q_{\omega^-}(s',a'),
\]
et on minimise
\[
\mathcal{L}(\omega)=\E\big[(Q_\omega(s,a)-y)^2\big].
\]
C'est exactement la structure standard d'un DQN.

\subsection*{Pourquoi le replay buffer stabilise l'apprentissage}
Les trajectoires $(s_t,a_t,r_t,s_{t+1})$ sont fortement corrélées dans le temps.
Apprendre directement dessus rend l'optimisation instable (les gradients «poursuivent» un signal non stationnaire).
Le replay buffer casse cette corrélation en rééchantillonnant des transitions passées, ce qui rapproche l'hypothèse i.i.d. des méthodes de régression.

\subsection*{Biais d'overestimation et Double DQN (perspective)}
La cible $\max_{a'}Q_{\omega^-}(s',a')$ a tendance à surestimer les valeurs (max sur des estimateurs bruités).
Une extension standard est \emph{Double DQN} : sélectionner l'action avec le réseau en ligne et l'évaluer avec le réseau cible.
C'est une extension naturelle si l'on veut améliorer la stabilité sans changer le cadre global.

\begin{comment}
\section{D. Numérique / algorithmes (stabilité, complexité, choix)}
Les stabilisateurs DQN présents :
\begin{itemize}[leftmargin=*]
\item \textbf{Experience replay} : réduire corrélation temporelle des samples (\codeid{ReplayBuffer}).
\item \textbf{Target network} : stabiliser la cible (hard update périodique).
\item \textbf{Epsilon-greedy} : exploration durant l'entraînement (\codeid{_epsilon_at_step}).
\item \textbf{Optimiseur} : Adam (implémentation numpy).
\end{itemize}

\end{comment}
\begin{comment}
\section{E. Implémentation dans le repo (fichiers, fonctions, pipeline, paramètres)}
\subsection*{Ce qui est \emph{exactement} codé}
\begin{itemize}[leftmargin=*]
\item Environnement \textbf{réel} (Alpaca) : \codeid{HedgingEnv} dans \filepath{app/model/hedger_v2/env.py}.
\item Environnement \textbf{synthétique} (training offline) : \codeid{SyntheticHedgingEnv} dans le même fichier.
\item Agent DQN numpy : \codeid{DQNAgent}, \codeid{ReplayBuffer}, \codeid{NumpyMLP} dans \filepath{app/model/hedger_v2/dqn_hedger.py}.
\item Persistence : \filepath{cache/HedgerDQN/dqn-v1-numpy.npz} et \filepath{cache/HedgerDQN/dqn-v1-numpy.json} (chemins construits via \codeid{CACHE_CSV_DIR} dans \filepath{app/utils/paths.py}).
\item Intégration UI : \filepath{app/vue/tabs/tab_hedging_systems.py} appelle \filepath{app/controller/hedger_v2_controller.py}.
\end{itemize}

\subsection*{Important : delta \emph{proxy}, pas delta réel}
Dans \codeid{HedgingEnv._current_state} (\filepath{app/model/hedger_v2/env.py}), la variable \codeid{net_delta_proxy} est calculée par :
\begin{itemize}[leftmargin=*]
\item $+1$ par call, $-1$ par put, multiplié par la quantité,
\end{itemize}
ce qui n'est pas le delta BS ni un delta de marché. Il faut donc lire ce module comme un \textbf{prototype pédagogique}.

\subsection*{Extrait de code : update DQN (cible Q-learning)}
Dans \filepath{app/model/hedger_v2/dqn_hedger.py}, \codeid{train_on_batch} :

\begin{lstlisting}[caption={DQN TD target (extrait) : \filepath{app/model/hedger_v2/dqn_hedger.py}, \codeid{train_on_batch}.},label={lst:dqn-train}]
q_pred_all, cache = self.q_net.forward(states, retain_cache=True)
q_pred = q_pred_all[np.arange(actions.shape[0]), actions]

q_next = self.target_net.predict(next_states)
max_q_next = np.max(q_next, axis=1)
q_target = rewards + self.config.gamma * (1.0 - dones) * max_q_next

td = (q_pred - q_target).astype(np.float32)
loss = float(np.mean(td * td))
\end{lstlisting}

\paragraph{Traduction directe de la formule TD.}
\begin{itemize}[leftmargin=*]
\item \codeid{q_target} implémente $y=r+\gamma(1-d)\max_{a'}Q_{\omega^-}(s',a')$.
\item \codeid{loss = mean((q_pred-q_target)^2)} est la MSE TD.
\end{itemize}

\subsection*{Ce qui n'est pas implémenté}
\begin{itemize}[leftmargin=*]
\item SARSA / Expected SARSA : \NonImplem\ (aucun module dédié trouvé).
\item Importance sampling off-policy : \NonImplem\ (pas de ratio $\pi/\mu$).
\item Mesures de risque «finance» (CVaR, pénalité de gamma/vega, PnL complet) : \NonImplem\ dans le DQN actuel.
\end{itemize}

\subsection*{Proposition «production-grade» (non présente dans le code)}
\textbf{Design recommandé} :
\begin{itemize}[leftmargin=*]
\item \textbf{État} : $\Delta,\Gamma,\text{Vega}$ estimés (BS ou modèle calibré), temps à maturité, inventaire, cash, volatilité réalisée, spread.
\item \textbf{Action} : ajustement continu/discret borné, contraintes de position.
\item \textbf{Reward} : $-\text{RiskMeasure}(\text{PnL})-\lambda\cdot\text{costs}$, où RiskMeasure peut être CVaR.
\item \textbf{Environnement} : simulation sous modèle calibré (Heston/Bates/rough), coûts réalistes, slippage.
\item \textbf{Algo} : DQN double + dueling + prioritized replay, ou actor-critic (SAC) selon action continue.
\end{itemize}
Cette section est une \emph{extension crédible} demandée, pas une description du dépôt.

\section{F. Exemple end-to-end}
Dans l'UI «Hedging Systems» :
\begin{enumerate}[leftmargin=*]
\item L'utilisateur clique «Train / update DQN» (\filepath{app/vue/tabs/tab_hedging_systems.py}) $\to$ \codeid{ctrl.train_dqn_model}.
\item Le modèle est entraîné offline sur \codeid{SyntheticHedgingEnv} et sauvegardé (\filepath{app/model/hedger_v2/dqn_hedger.py}, \codeid{load_or_train_dqn_model}).
\item L'utilisateur clique «Get DQN hedge suggestion» $\to$ \codeid{suggest_hedge_action}. Le DQN choisit l'action argmax.
\item Optionnel : «Execute» soumet un ordre Alpaca via \codeid{AlpacaHedgerClient.submit_market_order} (\filepath{app/model/hedger_v2/alpaca_client.py}).
\end{enumerate}

\section{G. Limites \& extensions crédibles}
\begin{itemize}[leftmargin=*]
\item Le modèle est entraîné sur un \textbf{environnement synthétique} très éloigné d'un portefeuille options réel : risque majeur de non-transférabilité.
\item Le proxy de delta est grossier : en production, il faut des grecs calculés (BS ou modèle calibré) et une PnL complète.
\item Extension : backtesting sur historique (non implémenté pour RL hedger).
\end{itemize}

\end{comment}
\begin{jurybox}
Un DQN de couverture est typiquement un \textbf{prototype} au départ : il démontre la chaîne MDP $\to$ estimation de $Q$ $\to$ recommandation d'action, mais la performance dépend surtout du choix d'état/reward et de la qualité du simulateur. Le point clé est la validation : stabilité hors échantillon et robustesse aux coûts/risques.
\end{jurybox}

% =========================================================
% PARTIE IV — Yield Curve
% =========================================================
\part{Yield Curve}

\chapter{Courbe de taux : DF, taux zéro, forwards, Nelson--Siegel}

\section{A. Problème financier (objectif, inputs, outputs)}
Le pricing et l'actualisation nécessitent une courbe de taux : discount factors $DF(T)$, taux zéro $z(T)$, forwards $f(t_1,t_2)$. On veut :
\begin{itemize}[leftmargin=*]
\item construire une courbe à partir de n\oe{}uds (tenors, zéros, DF) ;
\item interpoler de façon stable et sans arbitrage grossier ;
\item fournir $r(T)$, $DF(T)$ et des forwards cohérents ;
\item proposer un lissage paramétrique (Nelson--Siegel) si les données sont bruitées ou rares.
\end{itemize}

\textbf{Inputs} : n\oe{}uds (tenor, $T$, zero rate ou DF), devise, maturité de référence.
\textbf{Outputs} : courbe (DF/zéros/forwards), taux instantanés (si dérivables) et paramètres Nelson--Siegel (si ajustés).

\section{B. Théorie (définitions, hypothèses)}
\subsection*{Définitions formelles}
\begin{itemize}[leftmargin=*]
\item \textbf{Discount factor} : $DF(T)$, valeur aujourd'hui de 1 unité à $T$.
\item \textbf{Taux zéro (comp. continue)} : $z(T)$ défini par
\[
DF(T)=e^{-z(T)T}\quad\Rightarrow\quad z(T)=-\frac{1}{T}\ln DF(T).
\]
\item \textbf{Forward} $f(t_1,t_2)$ (comp. continue) :
\[
f(t_1,t_2)=-\frac{\ln DF(t_2)-\ln DF(t_1)}{t_2-t_1}.
\]
\item \textbf{Forward instantané} :
\[
f(t)=-\frac{\dd}{\dd t}\ln DF(t).
\]
\end{itemize}

\subsection*{Intuition financière}
\begin{itemize}[leftmargin=*]
\item $DF(T)$ est l'objet «primitif» le plus sûr : il doit être positif et décroissant.
\item Interpoler en $\ln DF$ est souvent plus stable (et garantit $DF>0$).
\end{itemize}

\qa
{Fournir une actualisation cohérente, calculer forwards et alimenter les pricers en $r(T)$.}
{Zéro-coupons cohérents, interpolation stable, données de n\oe{}uds de qualité.}
{Si la courbe est pauvre, l'interpolation devient fragile : on peut être amené à régulariser, imposer une structure paramétrique, ou accepter une approximation (courbe plate) en tant que \emph{proxy}.}
{Applications : actualisation cohérente, extraction de forwards, et alimentation des pricers (PDE/MC/FFT) en taux de financement.}

\subsection*{Conventions : capitalisation, day count, et cohérence}
En pratique, les marchés utilisent des conventions (ACT/360, 30/360, comp. simple vs continue, etc.).
Dans un dossier théorique, on choisit souvent la capitalisation continue car elle rend les formules plus lisibles :
\[
DF(T)=e^{-\int_0^T r(u)\,\dd u}.
\]
Ce choix n'est pas «la vérité», mais une convention mathématique ; l'important est d'être \textbf{cohérent} dans toutes les conversions.

\subsection*{Arbitrage taux : contraintes minimales sur $DF$}
Une courbe sans arbitrage doit au minimum satisfaire :
\begin{itemize}[leftmargin=*]
\item $DF(T)>0$ ;
\item $DF$ décroissant en $T$ si les taux sont non négatifs (dans un monde où les taux peuvent être négatifs, la monotonie peut être localement violée, mais $DF$ reste positif) ;
\item régularité suffisante pour dériver des forwards (éviter des oscillations d'interpolation).
\end{itemize}
Interpoler en $\ln DF$ est une façon simple de garantir $DF>0$.

\section{C. Dérivations (ultra détaillées)}
\subsection*{DF $\leftrightarrow$ zéro-rate}
Si $DF(T)=e^{-z(T)T}$, alors :
\[
z(T)=-\frac{1}{T}\ln DF(T).
\]
Inversement, si $z(T)$ est donné :
\[
DF(T)=e^{-z(T)T}.
\]
Ces conversions sont le point de départ de tout pipeline de courbe : on passe d'un format de données (zéros, DF) à un autre en restant cohérent sur la convention.

\subsection*{Forward rate}
À partir de deux discount factors, on dérive :
\[
DF(t_2)=DF(t_1),e^{-f(t_1,t_2)(t_2-t_1)}
\quad\Rightarrow\quad
f(t_1,t_2)=-\frac{\ln DF(t_2)-\ln DF(t_1)}{t_2-t_1}.
\]
Cette expression définit le forward continu moyen sur l'intervalle $(t_1,t_2)$.

\subsection*{Lien entre zéro-rate et forward instantané}
En définissant $z(T)=-\frac1T\ln DF(T)$ et $f(T)=-\frac{\dd}{\dd T}\ln DF(T)$, on obtient :
\[
\ln DF(T)=-Tz(T)\quad\Rightarrow\quad f(T)=z(T)+Tz'(T).
\]
Cette relation explique un phénomène empirique : une courbe de zéros «douce» peut cacher des forwards plus volatils si $z'(T)$ oscille.
Elle justifie l'intérêt de paramétrisations lisses (Nelson--Siegel) quand on veut des forwards stables.

\subsection*{Nelson--Siegel (et extension Svensson)}
La paramétrisation Nelson--Siegel :
\[
z(T)=\beta_0+\beta_1\frac{1-e^{-T/\tau}}{T/\tau}+\beta_2\Big(\frac{1-e^{-T/\tau}}{T/\tau}-e^{-T/\tau}\Big).
\]
Deux faits importants pour la calibration :
\begin{itemize}[leftmargin=*]
\item \textbf{linéarité conditionnelle} : si $\tau$ est fixé, la relation est linéaire en $(\beta_0,\beta_1,\beta_2)$, donc ajustable par moindres carrés ;
\item \textbf{choix de $\tau$} : on peut faire une recherche sur grille (robuste) ou une optimisation non linéaire (plus délicate).
\end{itemize}

\paragraph{Extension Nelson--Siegel--Svensson (NSS).}
Une extension classique ajoute un second terme de courbure (échelle $\tau_2$) :
\[
z(T)=\beta_0+\beta_1\frac{1-e^{-T/\tau_1}}{T/\tau_1}+\beta_2\Big(\frac{1-e^{-T/\tau_1}}{T/\tau_1}-e^{-T/\tau_1}\Big)
\;+\;\beta_3\Big(\frac{1-e^{-T/\tau_2}}{T/\tau_2}-e^{-T/\tau_2}\Big).
\]
NSS est utile quand la courbe présente plusieurs bosses (flexibilité accrue), au prix d'une calibration plus sensible (plus de paramètres).

\subsection*{Lecture économique des paramètres $(\beta_0,\beta_1,\beta_2,\tau)$}
\begin{itemize}[leftmargin=*]
\item $\beta_0$ : niveau de long terme (quand $T\to\infty$, les deux autres chargements tendent vers 0).
\item $\beta_1$ : pente court terme (agit surtout sur les maturités courtes).
\item $\beta_2$ : courbure (charge autour d'une maturité caractéristique contrôlée par $\tau$).
\item $\tau$ : échelle de temps : déplace la zone où la courbure est maximale.
\end{itemize}
Cette lecture justifie le calibrage par grille sur $\tau$ : fixer $\tau$ rend le modèle linéaire en $(\beta_0,\beta_1,\beta_2)$, donc calibrable par moindres carrés.

\begin{comment}
\section{D. Numérique / algorithmes (stabilité, complexité, choix)}
\begin{itemize}[leftmargin=*]
\item \textbf{Interpolation DF} : log-linéaire en discount factors (\codeid{log_discount_interp}) ($\Rightarrow DF>0$).
\item \textbf{Interpolation zéro} : fallback linéaire (\codeid{linear_zero_interp}).
\item \textbf{NS calibration} : grid search (robuste, peu coûteux) plutôt qu'optimisation non linéaire.
\end{itemize}

\end{comment}
\begin{comment}
\section{E. Implémentation dans le repo (fichiers, fonctions, pipeline, paramètres)}
\subsection*{Courbe de taux consommée par les pricers}
Le pricer MC européen dans \filepath{app/model/options/mc_engine.py} appelle \codeid{get_r(T)} depuis \filepath{app/model/yieldcurve/rates_utils.py}. Cette fonction :
\begin{itemize}[leftmargin=*]
\item tente de charger la courbe active via \codeid{get_active_curve} (\filepath{app/model/yieldcurve/service.py}),
\item sinon renvoie \codeid{DEFAULT_RF} (env \codeid{DEFAULT_RF_RATE}, défaut 2\%).
\end{itemize}

\subsection*{Chargement de données : nodes \& cache}
\begin{itemize}[leftmargin=*]
\item N\oe{}uds utilisateurs : \filepath{data/yield_curves/<CCY>_nodes.(csv|json)} (chargés via \codeid{load_nodes_from_file} dans \filepath{app/model/yieldcurve/parsing.py}).
\item Cache CSV : \filepath{cache/yield_curve.csv} chargé via \filepath{app/model/yieldcurve/loader.py}. Par défaut, si absent, \codeid{download_yield_curve_to_cache} crée un placeholder (taux constants 2.0).
\item Optionnel : un downloader FRED complet existe dans \filepath{app/model/yieldcurve/builder.py} (téléchargement \emph{réel} de séries \codeid{DGS}), mais n'est pas automatiquement appelé par défaut dans \codeid{get_active_curve} (car \codeid{allow_api=False} par défaut).
\end{itemize}

\section{F. Exemple end-to-end}
\begin{enumerate}[leftmargin=*]
\item L'utilisateur ouvre l'onglet \filepath{app/vue/tabs/tab_yieldcurve.py} et importe un fichier nodes (CSV/JSON).
\item La Vue appelle \codeid{yc.import_curve_nodes_upload} (\filepath{app/controller/yieldcurve_controller.py}) qui route vers \codeid{save_curve_nodes} (\filepath{app/model/yieldcurve/service.py}).
\item \codeid{get_curve_snapshot} renvoie : n\oe{}uds, DF/zc sur une grille, forwards, et (si $\ge3$ points) un lissage Nelson--Siegel.
\item Les autres modules (options/MC) peuvent ensuite appeler \codeid{get_r(T)} et utiliser ce $r(T)$.
\end{enumerate}

\section{G. Limites \& extensions crédibles}
\begin{itemize}[leftmargin=*]
\item Le cache \codeid{yield_curve.csv} peut être un placeholder constant (voir \filepath{app/model/yieldcurve/loader.py}) : en production, activer une source robuste (nodes, FRED).
\item \NonImplem\ : Nelson--Siegel--Svensson, bootstrapping instrument-based (swap curve) et arbitrage-free constraints.
\end{itemize}

\end{comment}
\begin{jurybox}
La courbe de taux dans l'app est un module cohérent : n\oe{}uds $\to$ DF/zc via interpolation log-DF, forwards et lissage Nelson--Siegel (si assez de points). Le fallback (courbe plate) est explicite, ce qui est essentiel pour l'audit.
\end{jurybox}

% =========================================================
% PARTIE V — Portfolio allocation \& risk
% =========================================================
\part{Portfolio allocation \& risk}

\chapter{Allocation de portefeuille : Markowitz, Risk Parity, Eigen-portfolios, rebalancement}

\section{A. Problème financier (objectif, inputs, outputs)}
On veut déterminer des poids cibles $w\in\R^n$ sur $n$ actifs à partir de rendements historiques, puis générer un plan de rebalancement (ordres) vers ces poids.

\textbf{Inputs} : matrice de rendements (historique), contraintes implicites (poids $\ge 0$, somme=1), portefeuille courant (valeurs de marché), prix spot.
\textbf{Outputs} : poids cibles, plan de rebalancement (achats/ventes, quantités) et indicateurs de risque/turnover associés.

\section{B. Théorie (définitions, hypothèses)}
\subsection*{(1) Markowitz}
On modélise les rendements $R\in\R^n$ par une moyenne $\mu$ et une covariance $\Sigma$. Le problème min-variance (sans contraintes additionnelles) est :
\[
\min_{w}; w^\top \Sigma w \quad\text{s.c.}\quad \sum_i w_i=1,; w_i\ge 0.
\]
Le «max Sharpe» (simplifié sans taux sans risque explicite) vise un $w$ proportionnel à $\Sigma^{-1}\mu$ (puis projeté sur le simplex).

\subsection*{(2) Risk parity (ERC)}
On cherche des contributions au risque égales. Une version intuitive : ajuster $w$ pour que
\[
\text{RC}_i \approx \frac{1}{n},\quad
\text{RC}_i=\frac{w_i(\Sigma w)_i}{\sqrt{w^\top \Sigma w}}.
\]

\subsection*{(3) Eigen-portfolios}
La diagonalisation $\Sigma = V\Lambda V^\top$ donne des portefeuilles «facteurs» $v_k$. Selon le choix du vecteur propre, on obtient :
\begin{itemize}[leftmargin=*]
\item un portefeuille direction «variance maximale» (plus grand $\lambda$),
\item ou «variance minimale» (plus petit $\lambda$).
\end{itemize}
Le choix dépend de l'objectif (exposition aux facteurs dominants vs portefeuille plus market-neutral) et doit être stabilisé (PCA sensible au bruit).

\qa
{Transformer des données historiques en poids cibles, puis en ordres, dans un cadre simple et audit-able.}
{Stationnarité relative des rendements, covariance estimable, absence de contraintes complexes (levier, turnover) dans le solveur.}
{Sensibilité aux fenêtres (lookback), instabilité de la covariance, coûts ignorés, contraintes simplistes.}
{Applications : allocation long-only simple, rebalancement périodique, et comparaison de familles (min-variance, risk parity, PCA) sous contraintes de base.}

\subsection*{Estimation : le vrai problème de Markowitz}
La théorie suppose $\mu$ et $\Sigma$ connus. En réalité :
\begin{itemize}[leftmargin=*]
\item $\mu$ est extrêmement bruitée (erreur d'estimation énorme) ;
\item $\Sigma$ est instable si $n$ est grand et l'historique court ;
\item les contraintes (long-only, turnover) dominent souvent la solution.
\end{itemize}
Conséquence pédagogique : beaucoup d'implémentations utilisent Markowitz surtout pour \textbf{min-variance} (moins dépendant de $\mu$) et ajoutent de la régularisation (shrinkage) implicitement (ici via $\Sigma+\varepsilon I$).

\subsection*{Risk parity : «ne pas dépendre de $\mu$»}
Le risk parity (ERC) vise une allocation dont le risque est réparti entre actifs/facteurs.
Il est souvent plus robuste que «max Sharpe» quand l'estimation de $\mu$ est peu fiable.
Interprétation : on choisit $w$ pour que chaque actif contribue de manière comparable à la volatilité du portefeuille.

\subsection*{Eigen-portfolios : lire la covariance comme des facteurs}
Décomposer $\Sigma$ revient à identifier des directions de risque (PCA).
Les vecteurs propres associés aux grandes valeurs propres correspondent à des facteurs dominants ; les petits peuvent correspondre à des spreads/portefeuilles plus «market-neutral».
Mais attention : les vecteurs propres sont eux-mêmes sensibles au bruit (instabilité de PCA en haute dimension).

\section{C. Dérivations (ultra détaillées)}
\subsection*{Markowitz min-variance (forme fermée sans contrainte de positivité)}
Si on ignore $w_i\ge 0$, le Lagrangien :
\[
\mathcal{L}(w,\lambda)=w^\top \Sigma w - \lambda(\bm{1}^\top w - 1).
\]
Condition d'optimalité :
\[
2\Sigma w - \lambda \bm{1}=0 \Rightarrow w = \frac{\lambda}{2}\Sigma^{-1}\bm{1}.
\]
En imposant $\bm{1}^\top w=1$ :
\[
\frac{\lambda}{2}\bm{1}^\top\Sigma^{-1}\bm{1}=1 \Rightarrow
w^\star=\frac{\Sigma^{-1}\bm{1}}{\bm{1}^\top\Sigma^{-1}\bm{1}}.
\]
En long-only, on ne peut pas garder cette forme fermée : on projette (ou on résout un QP) sous la contrainte $w\ge 0$ et $\bm{1}^\top w=1$ (simplex).

\subsection*{Max Sharpe (simplifié)}
Sans taux sans risque explicite, une heuristique est :
\[
w \propto \Sigma^{-1}\mu,
\]
puis normalisation sur le simplex (et éventuellement bornes, levier, turnover).

\subsection*{Max Sharpe (avec taux sans risque) : portefeuille tangent}
Si l'on introduit un taux sans risque $r_f$ et des rendements excédentaires $\tilde\mu=\mu-r_f\bm{1}$, alors le portefeuille tangent (sans contraintes) vérifie
\[
w^\star \propto \Sigma^{-1}\tilde\mu.
\]
La constante de proportionnalité se fixe par la contrainte $\bm{1}^\top w=1$ (ou par un levier choisi).
Cette dérivation montre que «max Sharpe» est une direction, pas une solution unique : la normalisation (simplex, levier) est un choix de contraintes.

\subsection*{Risk parity (descente de gradient simplifiée)}
Une itération de gradient (schématique) consiste à mettre à jour :
\[
w \leftarrow w - \eta(\text{RC}(w)-\tfrac1n\bm{1}),
\]
puis projette sur le simplex.

\subsection*{Condition ERC (équations) : ce que l'on cherche vraiment}
En notant $\sigma_p=\sqrt{w^\top\Sigma w}$, la contribution au risque de l'actif $i$ est
\[
\text{RC}_i=\frac{w_i(\Sigma w)_i}{\sigma_p}.
\]
La condition «Equal Risk Contribution» s'écrit $\text{RC}_i=\text{RC}_j$ pour tout $i,j$ (ou $\text{RC}_i=1/n$ après normalisation).
C'est un système non linéaire, d'où l'usage d'itérations (gradient) plutôt qu'une formule fermée.

\begin{comment}
\section{D. Numérique / algorithmes (stabilité, complexité, choix)}
\begin{itemize}[leftmargin=*]
\item Estimation covariance : \codeid{np.cov} puis régularisation $\Sigma+\varepsilon I$ (\codeid{1e-6}).
\item Projection simplex : assure $w_i\ge 0$ et $\sum w_i=1$ (\codeid{_project_simplex}).
\item Risk parity : itérations fixes, pas adaptatif constant (\codeid{lr=0.01} par défaut).
\end{itemize}

\end{comment}
\begin{comment}
\section{E. Implémentation dans le repo (fichiers, fonctions, pipeline, paramètres)}
\subsection*{Pipeline complet (UI)}
\begin{itemize}[leftmargin=*]
\item Rendements : \codeid{compute_returns_matrix} appelle \codeid{AlpacaPortfolioClient.get_history} puis \codeid{pct_change} (\filepath{app/model/portfolio_allocation/engine.py}).
\item Optimisation : \codeid{markowitz_optimize}, \codeid{risk_parity_optimize}, \codeid{eigen_portfolio_optimize}.
\item Ordres : \codeid{compute_rebalance_orders} compare valeurs cibles vs actuelles et divise par le prix courant.
\item Exécution : \codeid{execute_rebalance_orders} envoie des market orders via Alpaca.
\end{itemize}

\subsection*{Remarque : eigen-portfolio dans ce module}
Dans \filepath{app/model/portfolio_allocation/engine.py}, \codeid{eigen_portfolio_optimize} choisit \codeid{idx = np.argmin(vals)} (plus petite valeur propre), donc un \emph{portefeuille de variance minimale} dans la base des vecteurs propres (puis projection non négative).
Un autre module \filepath{app/model/portfolio/stats.py} calcule des «eigen orders» via la \emph{plus grande} valeur propre (principal), mais ce module n'est pas branché à l'UI (\NonImplem\ au sens «utilisé dans le flux principal»).

\section{F. Exemple end-to-end}
Dans \filepath{app/vue/tabs/tab_portfolio_and_risk.py} :
\begin{enumerate}[leftmargin=*]
\item L'utilisateur choisit «Markowitz min var» et clique «Compute target allocation».
\item La Vue appelle \codeid{ctrl.compute_allocation} (\filepath{app/controller/portfolio_and_risk_controller.py}).
\item Le contrôleur récupère l'historique via Alpaca, calcule les rendements, appelle l'optimiseur, puis renvoie \codeid{target_weights}.
\item «Generate rebalance plan» appelle \codeid{compute_rebalance_orders}, puis affiche les ordres.
\item «Execute rebalance» envoie les ordres à Alpaca (dangereux, explicitement signalé).
\end{enumerate}

\section{G. Limites \& extensions crédibles}
\begin{itemize}[leftmargin=*]
\item Coûts de transaction et liquidité ne sont pas modélisés dans l'optimisation ni le plan ; extension : pénalités de turnover, contraintes de lot, slippage.
\item Risque : aucune contrainte de volatilité cible ni budget de risque explicite ; extension : optimiser sous contrainte de volatilité.
\end{itemize}

\end{comment}
\begin{jurybox}
À retenir : en allocation, l'essentiel n'est pas la formule fermée mais l'\textbf{estimation} ($\mu,\Sigma$), la \textbf{régularisation} et les \textbf{contraintes} (long-only, turnover, coûts). Les solveurs (min-variance, risk parity, PCA) ne sont que des façons différentes de transformer ces estimations en poids cibles.
\end{jurybox}

\chapter{Risk management : exposition, PnL, VaR-lite, alertes}

\section{A. Problème financier (objectif, inputs, outputs)}
Un tableau de bord de risque doit répondre à des questions basiques :
\begin{itemize}[leftmargin=*]
\item Quelle exposition totale et nette ?
\item Quelle concentration par position ?
\item Quel PnL (réalisé / latent) ?
\item Une VaR indicative (même simplifiée) ?
\item Quelles alertes opérationnelles ?
\end{itemize}

\textbf{Inputs} : positions, historique de prix, paramètres (horizon, niveau de confiance, fenêtres).
\textbf{Outputs} : métriques de risque et alertes (exposition, concentration, PnL, indicateurs quantiles).

\section{B. Théorie (définitions, hypothèses)}
\subsection*{Exposition}
\[
\text{Exposure}=\sum_i |MV_i|,\qquad
\text{NetExposure}=\sum_i \text{sign}_i\cdot |MV_i|.
\]
\subsection*{VaR historique (idée)}
Une VaR historique classique utiliserait des rendements et des poids. Une version «VaR-lite» consiste à prendre un quantile d'une série agrégée de variations et à le convertir en montant via la valeur du portefeuille.

\qa
{Donner une visibilité opérationnelle rapide sur le risque d'un portefeuille, sans prétendre à une VaR institutionnelle.}
{Données de prix disponibles ; positions correctement valorisées ; historique suffisant.}
{La VaR-lite n'est pas une VaR standard (dépend de la construction des retours) ; pas de stress tests ; pas de facteurs.}
{Applications : dashboard de risque «lite» (exposition, concentration, PnL), quantiles indicatifs, et base de discussion pour des extensions (stress, ES, facteurs).}

\subsection*{VaR : définition standard (pour situer «VaR-lite»)}
Pour un horizon $h$ et un niveau de confiance $c\in(0,1)$, la Value-at-Risk d'un PnL $L$ (perte positive) est
\[
\mathrm{VaR}_{c}(L)=\inf\{\ell:\Prob(L\le \ell)\ge c\}.
\]
Si l'on travaille en \emph{rendement} $R$ (perte $L=-R$), la VaR correspond à un quantile de la distribution des rendements.
Trois familles classiques :
\begin{itemize}[leftmargin=*]
\item \textbf{VaR paramétrique} : hypothèse de loi (souvent normale) sur $R$.
\item \textbf{VaR historique} : quantile empirique sur des rendements observés.
\item \textbf{VaR Monte Carlo} : simulation de scénarios via un modèle.
\end{itemize}
La VaR-lite est un indicateur opérationnel qui se rapproche d'une VaR historique, mais avec une construction de «retours» simplifiée.

\subsection*{VaR vs ES (Expected Shortfall)}
La VaR ne dit rien sur la gravité au-delà du quantile.
L'Expected Shortfall (ES) complète l'information :
\[
\mathrm{ES}_c(L)=\E[L\mid L\ge \mathrm{VaR}_c(L)].
\]
C'est une extension naturelle si l'on souhaite un indicateur plus sensible aux queues.

\section{C. Dérivations (ultra détaillées)}
\subsection*{Exposition et net}
Pour chaque position $i$, on note $MV_i$ sa valeur de marché (positive en valeur absolue) et $\text{sign}_i\in\{-1,+1\}$ son sens (short/long). On définit :
\begin{itemize}[leftmargin=*]
\item \textbf{exposition brute} : $\sum_i |MV_i|$ ;
\item \textbf{exposition nette} : $\sum_i \text{sign}_i\,|MV_i|$.
\end{itemize}
Ces définitions sont directes (aucune hypothèse de covariance) : elles mesurent la taille du portefeuille et le biais directionnel global.

\subsection*{VaR-lite : ce qui est exactement fait}
Une construction possible (très simplifiée) est :
\begin{itemize}[leftmargin=*]
\item construire un panel de prix (séries temporelles) ;
\item calculer des variations $\Delta S_t$ (différences) au lieu de rendements $\Delta S_t/S_t$ ;
\item agréger ces variations (moyenne ou somme pondérée) pour obtenir une série scalaire ;
\item prendre un quantile $\alpha=1-\text{confidence}$ ;
\item convertir en montant via une échelle de portefeuille (valeur totale, notionnel).
\end{itemize}
C'est une métrique indicative : elle donne un ordre de grandeur, mais ne coïncide pas avec une VaR standard (qui travaillerait sur rendements et pondérations).

\subsection*{Pourquoi les différences changent l'interprétation}
Un rendement standard est $\Delta S / S$ (ou $\ln(S_{t+\Delta t}/S_t)$), ce qui est sans dimension et comparable entre actifs.
Une différence $S_{t+\Delta t}-S_t$ dépend de l'unité de prix : un actif à 300\$ et un actif à 30\$ ne sont pas comparables.
Donc la VaR-lite telle qu'écrite mesure davantage une \emph{variation moyenne en dollars} qu'un quantile de rendement de portefeuille.
C'est acceptable comme indicateur «lite» si l'on le lit comme tel, mais ce n'est pas une VaR au sens standard.

\subsection*{VaR delta-normale (rappel pédagogique)}
Si l'on suppose un rendement de portefeuille $R_p$ normal, $R_p\sim\Normal(\mu_p,\sigma_p^2)$, alors le quantile s'écrit
\[
q_\alpha(R_p)=\mu_p+\sigma_p\,\Phi^{-1}(\alpha),
\]
où $\Phi^{-1}$ est l'inverse de la CDF normale.
La VaR en valeur monétaire est ensuite $\mathrm{VaR}= -q_\alpha(R_p)\cdot V$ (avec $V$ la valeur du portefeuille), selon les conventions de signe.
Ce rappel est utile pour comprendre les variantes «lite» : elles essaient toutes de fabriquer un proxy de distribution pour en prendre un quantile.

\begin{comment}
\section{D. Numérique / algorithmes}
\begin{itemize}[leftmargin=*]
\item Robustesse : nombreux fallbacks si offline ou données vides.
\item Le module s'arrête à des calculs $O(Tn)$ simples (pas de solveur).
\end{itemize}

\end{comment}
\begin{comment}
\section{E. Implémentation dans le repo}
\begin{itemize}[leftmargin=*]
\item Moteur : \filepath{app/model/risk_management/engine.py}.
\item Contrôleur : \filepath{app/controller/portfolio_and_risk_controller.py} (\codeid{get_risk_summary} assemble exposition, VaR-lite, alertes, PnL).
\item UI : \filepath{app/vue/tabs/tab_portfolio_and_risk.py}.
\end{itemize}

\section{F. Exemple end-to-end}
\begin{enumerate}[leftmargin=*]
\item Ouverture de l'onglet \filepath{app/vue/tabs/tab_portfolio_and_risk.py}.
\item Appel \codeid{ctrl.get_risk_summary()}.
\item Affichage : exposition (camembert), net exposure (bar), PnL \& VaR-lite, alertes, puis table détaillée.
\end{enumerate}

\section{G. Limites \& extensions crédibles}
\begin{itemize}[leftmargin=*]
\item Remplacer \codeid{diff()} par des rendements % et intégrer des poids de portefeuille.
\item Ajouter stress tests (scénarios), CVaR, contributions au risque par position.
\end{itemize}

\end{comment}
\begin{jurybox}
Un dashboard de risque «lite» doit être lu comme un instrument de pilotage : exposition, concentration, alertes, et quantiles indicatifs. Le point clé est la transparence des hypothèses (rendements vs différences, pondérations, fenêtre), car ces choix changent radicalement l'interprétation.
\end{jurybox}

% =========================================================
% ANNEXES
% =========================================================
\appendix

\chapter{Annexe : lien avec le projet PaperTradingApp}
Cette annexe regroupe le \textbf{pont} entre le dossier (théorie/dérivations) et le projet :
\begin{itemize}[leftmargin=*]
\item inventaire des payoffs et des moteurs de pricing présents ;
\item figures de payoffs pour vérifier les conventions ;
\item table de correspondance module $\leftrightarrow$ notion théorique $\leftrightarrow$ implémentation ;
\item bibliographie minimale.
\end{itemize}

\chapter{Pricing des options présentes dans le repo}

\section*{A. Inventaire (payoff \& méthode)}
Les familles de payoffs et moteurs présents sont repérables via :
\begin{itemize}[leftmargin=*]
\item Payoffs \& vues (payoff/profit) : \filepath{app/model/options/core/pricing_lib.py} (\codeid{payoff_}, \codeid{view_*}).
\item Black--Scholes (formules) : \filepath{app/model/options/engines/black_scholes.py}.
\item Arbres CRR (Am/Bmd) : \filepath{app/model/options/engines/crr.py}.
\item Monte Carlo exotiques (Asian, lookback, barrier, cliquet) : \filepath{app/model/options/core/trees.py}.
\item PDE BS Crank--Nicolson : \filepath{app/model/options/engines/pricing.py} (\codeid{CrankNicolsonBS}).
\end{itemize}

\section{Payoffs et pricing (table synthèse)}
\begin{longtable}{@{}p{3.6cm}p{6.0cm}p{5.2cm}@{}}
\toprule
\textbf{Produit} & \textbf{Payoff} & \textbf{Méthode(s) repérées}\\
\midrule
\endhead
Call européen & $(S_T-K)^+$ & BS : \filepath{app/model/options/engines/black_scholes.py} (\codeid{black_scholes_price}); PDE : \filepath{app/model/options/engines/pricing.py} (\codeid{CrankNicolsonBS}); MC : \filepath{app/model/options/mc_engine.py} (\codeid{price_mc_european})\\
Put européen & $(K-S_T)^+$ & Idem (BS/PDE/MC)\\
Americain (CRR) & exercice optimal & \filepath{app/model/options/engines/crr.py} (\codeid{price_american_crr}); aussi LSMC : \filepath{app/model/options/mc_engine.py} (\codeid{price_mc_lsmc})\\
Bermudan (CRR) & exercice sur dates discrètes & \filepath{app/model/options/engines/crr.py} (\codeid{price_bermuda_crr}); aussi LSMC via \filepath{app/model/options/mc_engine.py}\\
Digital (cash-or-nothing) & $\1_{{S_T>K}}$ (ou variantes) & BS : \filepath{app/model/options/engines/black_scholes.py} (\codeid{price_digital}); MC barrière digital : \filepath{app/model/options/core/trees.py} (\codeid{price_barrier_digital})\\
Asset-or-nothing & $S_T\1_{{S_T>K}}$ & \filepath{app/model/options/engines/black_scholes.py} (\codeid{price_asset_or_nothing})\\
Barrier (vanilla/digital) & actif si barrière non touchée (ou touchée) & MC : \filepath{app/model/options/core/trees.py} (\codeid{price_barrier_vanilla}, \codeid{price_barrier_digital}); PDE barrière : \filepath{app/model/options/engines/pricing.py} (\codeid{CN_Barrier_option})\\
Asian (arith/geo) & $(\bar S - K)^+$ ou $(\tilde S - K)^+$ & MC : \filepath{app/model/options/core/trees.py} (\codeid{price_asian_mc}, \codeid{price_asian_geo_mc})\\
Lookback & fonctions de min/max & MC : \filepath{app/model/options/core/trees.py} (\codeid{price_lookback_mc}); module dédié : \filepath{app/model/options/exotic/lookback/lookback_call.py} (\codeid{lookback_call_option})\\
Cliquet & somme de rendements cap/floor & MC : \filepath{app/model/options/core/trees.py} (\codeid{price_cliquet})\\
Spreads (straddle/strangle/etc.) & combinaisons linéaires de calls/puts & Payoffs \& vues : \filepath{app/model/options/core/pricing_lib.py}; pricing BS direct : \filepath{app/model/options/engines/black_scholes.py} (\codeid{price_straddle}, \codeid{price_strangle}, \ldots)\\
\bottomrule
\end{longtable}

\begin{jurybox}
Le dépôt contient plusieurs moteurs (BS, PDE, CRR, MC) et une bibliothèque de payoffs/visualisation. Pour un audit, les points d'entrée «métier» sont \filepath{app/model/options/service.py} et \filepath{app/controller/options_controller.py}.
\end{jurybox}

\chapter{Graphes de payoff (PGFPlots)}

\section{B. Graphes (payoff brut vs profit)}
Les figures ci-dessous sont des \textbf{illustrations} : on fixe $S$ (prix à maturité) sur $[50,150]$ et des strikes.
Pour les graphes «profit», on soustrait une prime \emph{exemple} (valeurs illustratives, pas un output automatique du dépôt).

\subsection*{Notations}
\[
\text{Payoff}(S) \quad\text{et}\quad \text{Profit}(S)=\text{Payoff}(S)-\text{Premium}.
\]

\begin{figure}[h]
\centering
\begin{tikzpicture}
\begin{axis}[
width=0.92\textwidth,height=6.0cm,
xlabel={$S_T$}, ylabel={Valeur},
title={Call (K=100) : payoff vs profit (premium exemple = 8.92)},
legend style={at={(0.02,0.98)},anchor=north west},
grid=both
]
\addplot[thick,blue,domain=50:150,samples=300]{max(x-100,0)};
\addlegendentry{Payoff call}
\addplot[thick,red,domain=50:150,samples=300]{max(x-100,0)-8.916};
\addlegendentry{Profit call}
\end{axis}
\end{tikzpicture}
\caption{Call : payoff brut et profit (prime illustrative).}
\end{figure}

\begin{figure}[h]
\centering
\begin{tikzpicture}
\begin{axis}[
width=0.92\textwidth,height=6.0cm,
xlabel={$S_T$}, ylabel={Valeur},
title={Put (K=100) : payoff vs profit (premium exemple = 6.94)},
legend style={at={(0.02,0.98)},anchor=north west},
grid=both
]
\addplot[thick,blue,domain=50:150,samples=300]{max(100-x,0)};
\addlegendentry{Payoff put}
\addplot[thick,red,domain=50:150,samples=300]{max(100-x,0)-6.936};
\addlegendentry{Profit put}
\end{axis}
\end{tikzpicture}
\caption{Put : payoff brut et profit (prime illustrative).}
\end{figure}

\begin{figure}[h]
\centering
\begin{tikzpicture}
\begin{axis}[
width=0.92\textwidth,height=6.0cm,
xlabel={$S_T$}, ylabel={Valeur},
title={Straddle (K=100) : payoff vs profit (premium exemple = 15.85)},
legend style={at={(0.02,0.98)},anchor=north west},
grid=both
]
\addplot[thick,blue,domain=50:150,samples=300]{max(x-100,0)+max(100-x,0)};
\addlegendentry{Payoff straddle}
\addplot[thick,red,domain=50:150,samples=300]{max(x-100,0)+max(100-x,0)-15.852};
\addlegendentry{Profit straddle}
\end{axis}
\end{tikzpicture}
\caption{Straddle.}
\end{figure}

\begin{figure}[h]
\centering
\begin{tikzpicture}
\begin{axis}[
width=0.92\textwidth,height=6.0cm,
xlabel={$S_T$}, ylabel={Valeur},
title={Strangle (Kp=95,Kc=105) : payoff vs profit (premium exemple = 10)},
legend style={at={(0.02,0.98)},anchor=north west},
grid=both
]
\addplot[thick,blue,domain=50:150,samples=300]{max(95-x,0)+max(x-105,0)};
\addlegendentry{Payoff strangle}
\addplot[thick,red,domain=50:150,samples=300]{max(95-x,0)+max(x-105,0)-10};
\addlegendentry{Profit strangle}
\end{axis}
\end{tikzpicture}
\caption{Strangle (primes illustratives).}
\end{figure}

\begin{figure}[h]
\centering
\begin{tikzpicture}
\begin{axis}[
width=0.92\textwidth,height=6.0cm,
xlabel={$S_T$}, ylabel={Valeur},
title={Call spread (long 95, short 105) : payoff vs profit (premium exemple = 4)},
legend style={at={(0.02,0.98)},anchor=north west},
grid=both
]
\addplot[thick,blue,domain=50:150,samples=300]{max(x-95,0)-max(x-105,0)};
\addlegendentry{Payoff call spread}
\addplot[thick,red,domain=50:150,samples=300]{max(x-95,0)-max(x-105,0)-4};
\addlegendentry{Profit call spread}
\end{axis}
\end{tikzpicture}
\caption{Call spread.}
\end{figure}

\begin{figure}[h]
\centering
\begin{tikzpicture}
\begin{axis}[
width=0.92\textwidth,height=6.0cm,
xlabel={$S_T$}, ylabel={Valeur},
title={Put spread (long 105, short 95) : payoff vs profit (premium exemple = 4)},
legend style={at={(0.02,0.98)},anchor=north west},
grid=both
]
\addplot[thick,blue,domain=50:150,samples=300]{max(105-x,0)-max(95-x,0)};
\addlegendentry{Payoff put spread}
\addplot[thick,red,domain=50:150,samples=300]{max(105-x,0)-max(95-x,0)-4};
\addlegendentry{Profit put spread}
\end{axis}
\end{tikzpicture}
\caption{Put spread.}
\end{figure}

\begin{figure}[h]
\centering
\begin{tikzpicture}
\begin{axis}[
width=0.92\textwidth,height=6.0cm,
xlabel={$S_T$}, ylabel={Valeur},
title={Butterfly (90,100,110) : payoff vs profit (premium exemple = 2)},
legend style={at={(0.02,0.98)},anchor=north west},
grid=both
]
\addplot[thick,blue,domain=50:150,samples=300]{max(x-90,0)-2*max(x-100,0)+max(x-110,0)};
\addlegendentry{Payoff butterfly}
\addplot[thick,red,domain=50:150,samples=300]{max(x-90,0)-2*max(x-100,0)+max(x-110,0)-2};
\addlegendentry{Profit butterfly}
\end{axis}
\end{tikzpicture}
\caption{Butterfly.}
\end{figure}

\begin{figure}[h]
\centering
\begin{tikzpicture}
\begin{axis}[
width=0.92\textwidth,height=6.0cm,
xlabel={$S_T$}, ylabel={Valeur},
title={Condor (90,95,105,110) : payoff vs profit (premium exemple = 1.5)},
legend style={at={(0.02,0.98)},anchor=north west},
grid=both
]
\addplot[thick,blue,domain=50:150,samples=300]{max(x-90,0)-max(x-95,0)-max(x-105,0)+max(x-110,0)};
\addlegendentry{Payoff condor}
\addplot[thick,red,domain=50:150,samples=300]{max(x-90,0)-max(x-95,0)-max(x-105,0)+max(x-110,0)-1.5};
\addlegendentry{Profit condor}
\end{axis}
\end{tikzpicture}
\caption{Condor.}
\end{figure}

\begin{figure}[h]
\centering
\begin{tikzpicture}
\begin{axis}[
width=0.92\textwidth,height=6.0cm,
xlabel={$S_T$}, ylabel={Valeur},
title={Iron condor (puts 90/95, calls 105/110) : payoff vs profit (premium exemple = 2)},
legend style={at={(0.02,0.98)},anchor=north west},
grid=both
]
% long put 90, short put 95, short call 105, long call 110
\addplot[thick,blue,domain=50:150,samples=300]{max(90-x,0)-max(95-x,0)-max(x-105,0)+max(x-110,0)};
\addlegendentry{Payoff iron condor}
\addplot[thick,red,domain=50:150,samples=300]{max(90-x,0)-max(95-x,0)-max(x-105,0)+max(x-110,0)-2};
\addlegendentry{Profit iron condor}
\end{axis}
\end{tikzpicture}
\caption{Iron condor.}
\end{figure}

\begin{jurybox}
Les payoffs sont des combinaisons linéaires de calls/puts. Le dépôt encode explicitement ces payoffs (via \codeid{np.maximum}) et construit des vues «payoff vs profit» dans \filepath{app/model/options/core/pricing_lib.py}. Les figures ci-dessus illustrent cette logique.
\end{jurybox}

\chapter{Table \og Code map\fg{} (module $\leftrightarrow$ théorie $\leftrightarrow$ implémentation)}

\begin{longtable}{@{}p{3.2cm}p{4.8cm}p{5.4cm}p{3.8cm}@{}}
\toprule
\textbf{Concept théorique} & \textbf{Module} & \textbf{Fichiers / fonctions} & \textbf{Sorties} \\
\midrule
\endhead
GBM / Brownien & Options (MC) & \filepath{app/model/options/core/trees.py} : \codeid{simulate_gbm_paths} ; \filepath{app/model/options/logic.py} : \codeid{simulate_gbm_paths} & Paths, prix MC \\
Itô / PDE BS & Options (PDE) & \filepath{app/model/options/engines/pricing.py} : \codeid{CrankNicolsonBS} ; \filepath{app/model/options/service.py} : \codeid{compute_price} & Prix, (Delta,Gamma,Theta) \\
Grecs BS & Options (greeks) & \filepath{app/model/options/core/greeks.py} : \codeid{compute_vanilla_greeks} & delta/gamma/vega/theta/rho \\
IV (inversion) & Calibration (IV) & \filepath{app/model/calibration/implied_vol.py} : \codeid{implied_vol_call}, \codeid{implied_vol_grid} & IV pointwise / grille \\
Surface (grille+masque) & Calibration (surface) & \filepath{app/model/calibration/market_surface.py} : \codeid{build_fixed_grid}, \codeid{make_fixed_grid} & \codeid{iv_grid}, mask \\
Moindres carrés + contraintes & Calibration & \filepath{app/model/calibration/heston_calibrator.py} : \codeid{calibrate_heston_least_squares} ; \filepath{app/model/volatility_models//calibrator.py} & params, runs, metrics \\
Heston CF (V1) & Calibration Heston & \filepath{app/model/calibration/heston_pricer.py} : \codeid{call_price_cf} & Prix call \\
Carr--Madan FFT & Vol models & \filepath{app/model/volatility_models/common/fft.py} : \codeid{carr_madan_fft_call_prices} & Grille (K,C) \\
Heston FFT & Vol models & \filepath{app/model/volatility_models/heston/} : \codeid{HestonFFTCalibrator} & params, \codeid{iv_model} \\
Merton/Bates CF & Vol models & \filepath{app/model/volatility_models/jump_diffusion/cf.py} & CF log-return \\
SABR Hagan & Vol models & \filepath{app/model/volatility_models/sabr/analytic.py} : \codeid{hagan_black_iv} & IV analytique \\
rHeston markovien & Vol models & \filepath{app/model/volatility_models/rheston/} : \codeid{fractional_kernel_markovian_approx}, \codeid{rheston_log_return_cf_markovian} & CF approx, \codeid{iv_model} \\
rBergomi MC surrogate & Vol models & \filepath{app/model/volatility_models/rbergomi/calibrator_mc_surrogate.py} & params, \codeid{iv_model} \\
Volterra MC proxy & Vol models & \filepath{app/model/volatility_models/volterra/calibrator_mc.py} & params, \codeid{iv_model} \\
Kalman (lissage) & Vol models & \filepath{app/model/volatility_models/kalman/kalman_filter.py} & \codeid{x_smooth} \\
DQN (RL) & Hedger v2 & \filepath{app/model/hedger_v2/dqn_hedger.py} : \codeid{DQNAgent}, \codeid{ReplayBuffer} & suggestion, modèle cache \\
MDP env (synthetic) & Hedger v2 & \filepath{app/model/hedger_v2/env.py} : \codeid{SyntheticHedgingEnv} & transitions, rewards \\
Yield curve (DF/forwards) & YieldCurve & \filepath{app/model/yieldcurve/curves.py} : \codeid{NodeYieldCurve} & z(T), DF(T), f(t1,t2) \\
Nelson--Siegel & YieldCurve & \filepath{app/model/yieldcurve/engine.py} : \codeid{fit_nelson_siegel_grid_search} & beta0,beta1,beta2,tau \\
Markowitz / Risk parity & Allocation & \filepath{app/model/portfolio_allocation/engine.py} : \codeid{markowitz_optimize}, \codeid{risk_parity_optimize} & weights \\
Rebalancement & Allocation & \filepath{app/model/portfolio_allocation/engine.py} : \codeid{compute_rebalance_orders} & ordres buy/sell \\
Risque (exposition, VaR-lite) & Risk mgmt & \filepath{app/model/risk_management/engine.py} : \codeid{compute_exposure}, \codeid{compute_var_lite} & métriques, alertes \\
\bottomrule
\end{longtable}

\begin{jurybox}
Cette table est l'outil d'audit : chaque concept du dossier renvoie à un fichier et une fonction précise. Si un concept est traité théoriquement mais absent du code (SARSA, importance sampling), il est explicitement marqué \NonImplem\ dans la Partie III.
\end{jurybox}

\chapter{Bibliographie minimale}

\begin{thebibliography}{99}

\bibitem{ito}
K.\ Itô.
\newblock \emph{On stochastic differential equations}.
\newblock Memoirs of the American Mathematical Society, 1951.

\bibitem{blackscholes}
F.\ Black and M.\ Scholes.
\newblock The pricing of options and corporate liabilities.
\newblock \emph{Journal of Political Economy}, 81(3), 1973.

\bibitem{merton73}
R.\ C.\ Merton.
\newblock Theory of rational option pricing.
\newblock \emph{Bell Journal of Economics}, 4(1), 1973.

\bibitem{feynman-kac}
M.\ Kac.
\newblock On distributions of certain {W}iener functionals.
\newblock \emph{Transactions of the American Mathematical Society}, 65(1), 1949.

\bibitem{heston}
S.\ L.\ Heston.
\newblock A closed-form solution for options with stochastic volatility with applications to bond and currency options.
\newblock \emph{Review of Financial Studies}, 6(2), 1993.

\bibitem{carrmadan}
P.\ Carr and D.\ Madan.
\newblock Option valuation using the fast {F}ourier transform.
\newblock \emph{Journal of Computational Finance}, 2(4), 1999.

\bibitem{mertonjump}
R.\ C.\ Merton.
\newblock Option pricing when underlying stock returns are discontinuous.
\newblock \emph{Journal of Financial Economics}, 3(1--2), 1976.

\bibitem{bates}
D.\ S.\ Bates.
\newblock Jumps and stochastic volatility: exchange rate processes implicit in {D}eutsche {M}ark options.
\newblock \emph{Review of Financial Studies}, 9(1), 1996.

\bibitem{hagan}
P.\ Hagan, D.\ Kumar, A.\ Lesniewski, and D.\ Woodward.
\newblock Managing smile risk.
\newblock \emph{Wilmott Magazine}, 2002.

\bibitem{roughvol}
J.\ Gatheral, T.\ Jaisson, and M.\ Rosenbaum.
\newblock Volatility is rough.
\newblock \emph{Quantitative Finance}, 18(6), 2018.

\end{thebibliography}

\end{document}
